\chapter{Source Code}
\section{Information}
I verify that I am the sole author of the programs contained in this folder, except where explicitly stated to the contrary - {Oliver H\''ohn}, \today.
\section{Table of Contents}
\begin{enumerate}
	\item src/main/java/backend: Package for the logic of the system.
		\begin{enumerate}
			\item helpers: Package for utility code used in different components in the system.
			\begin{enumerate}
				\item Sort.java: class that holds sorting methods used in the project.
				\item ToJSON.java: class used to turn a list of Results into a JSON String. Uses the GSON library.
				\item ToPDF.java: class used to save a list of Results into a PDF File. Uses the Apache PDFBox and Apache Commons libraries.
			\end{enumerate}
			\item process: Package for modules to process documents.
			\begin{enumerate}
				\item CallbackResult.java: interface for frontend classes that want to be notified when the files have been processed.
				\item Engine.java: class to process the text in a document to generate Result objects. Uses the Stanford CoreNLP library.
				\item FileData.java: class that represents data needed from a File.
				\item ProcessFileCallback.java: interface to inform observer when a File has been processed.
				\item ProcessFiles.java: class to process a list of Files (documents) and produce a set of Results. Uses the Apache PDFBox and Apache POI libraries.
				\item Result.java: class that represents the data for an event in the timeline.
				\item TimelineData: class that parses NER dates and holds start and end dates for an event. Uses the JODA Time library.
			\end{enumerate}
			\item ranges: Package for modules relating to placing Results in Ranges.
			\begin{enumerate}
				\item ProduceRanges.java: class to produce a list of Ranges (using the algorthm presented in the Design Chapter) from Results.
				\item Range.java: class that represents a node in a Tree, where dates encapsulate Results.
			\end{enumerate}
			\item system: Package for modules relating to the state and settings of the system.
			\begin{enumerate}
				\item BackEndSystem.java: class to hold the data needed by the entire Back-Entire (i.e. settings used, system state, etc.).
				\item Settings.java: class to represent the settings to be used in the system to process files, and build the UI. It saves the settings in a File. Uses the Stanford CoreNLP library.
				\item SystemState.java: enumeration of the different states the system can be in when a file is being processed.
			\end{enumerate}
		\end{enumerate}
	\item src/main/java/frontend: Package for the UI of the system. Use the JavaFX libraries.
	\begin{enumerate}
		\item controllers: Package of the view controllers of the system. 
		\begin{enumerate}
			\item CustomResultRowController.java: controller for the custom Result rows in the ListView of the RangeData layout.
			\item CustomTimelineRow.java: class for the custom row in the Traditional Timeline.
			\item DocumentLoadedRowController.java: controller for each row in the Documents-Loaded ListView.
			\item DocumentReaderController.java: controller for the layout of the Document Reader window. Uses the RichTextFX library.
			\item EditEventController.java: controller for the dialog used to edit events.
			\item ListViewController.java: controller for the layout where the ListView is shown.
			\item MenuBarControllerInter.java: interface for observers of a controller that has a menu bar.
			\item RangeDataController.java: controller for the layout that represents the data held by a Range.
			\item StartUpController.java: controller for the startup scene (i.e. the controller for the main window of the UI).
			\item TimelineRowController.java: controller class for each row in the traditional Timeline ListView.
		\end{enumerate}
		\item dialogs: Package for the dialogs used in the system.
		\begin{enumerate}
			\item AboutDialog.java: class to create the dialog providing information of the system.
			\item EditEventDialog.java: class that creates a dialog to edit an event.
			\item FileConfirmationDialog.java: class that creates a dialog to edit data of the files selected to process.
			\item LoadingDialog.java: class used to show/hide the loading dialog when data is being processed.
			\item RemoveConfirmationDialog.java: class to confirm the removal of a document from the timeline.
			\item SettingsDialog.java: class to show and edit the settings used in the system.
		\end{enumerate}
		\item helpers: Package for the helpers needed in the front-end of the system.
		\begin{enumerate}
			\item TextFieldState.java: enumeration class to represent the state the data in a TextField is in (i.e. is it validated)
		\end{enumerate}
		\item observers: Package for the Observer interfaces used to interact with the different layouts in the system.
		\begin{enumerate}
			\item DocumentReaderObserver.java: interface for the observer of the DocumentReaderController.
			\item DocumentsLoadedObserver.java: interface for the observer of the DocumentLoadedRowController.
			\item EditEventDialogObserver.java: interface for the observer of the EditEventDialog.
			\item MenuBarObserver.java: interface for the observer of layouts that have menu bars (and their items).
			\item StartUpObserver.java: interface for the observer of the StartUpController.
			\item TimelineObserver.java: interface for the observer of the Timeline layout (includes both kinds of timelines).
			\item TimelineRowObserver.java: interface for the observer of the TimelineRowController.
		\end{enumerate}
		\item Main.java: main class that is ran to begin the system. It prepares the back-end and shows the UI.
	\end{enumerate}
	\item src/main/resources: Resources used throughout the system.
	\begin{enumerate}
		\item controllers: Package for the resource files used by the controllers of the system.
		\begin{enumerate}
			\item customErrorFields.css: CSS file for input fields.
			\item customResultRow.fxml: FXML file for layout of rows in the Result.
			\item documentLoadedRow.fxml: FXML file for layout of rows of the documents loaded.
			\item editEventDialog.fxml: FXML file for the layout of an edit event dialog.
			\item icaddcircleoutline.png: PNG used to add subjects to an event.
			\item iccloseblack.png: PNG used to remove a document from the timeline.
			\item listViewTheme.css: CSS for the traditional timeline view.
			\item listViewThemeTimeline.css: CSS for the Range timeline view.
			\item rangeDataLayout.fxml: FXML file for the layout of a Range.
			\item timelineRowEven.fxml: FXML file for the layout of an even row in the traditional timeline.
			\item timelineRowOdd.fxml: FXML file for the layout of an odd row in the traditional timeline.
		\end{enumerate}
		\item dialogs: Package for the resource files used by the dialogs of the system.
		\begin{enumerate}
			\item customErrorFields.css: CSS file for input fields in the dialogs.
			\item loadingDialog.css: CSS file for the loading dialog when processing files.
		\end{enumerate}
		\item documentListViewTheme.css: CSS file for the documents loaded ListView.
		\item listview.fxml: FXML file for the layout of the window with the timelines and documents loaded.
		\item startup.fxml: FXML file for the layout of the initial window shown when the system is launched.
	\end{enumerate}
	\item src/test/backend: Unit Tests for the logic of the system. Uses the JUnit library.
	\begin{enumerate}
		\item EngineTest.java: test class for the processing done by the Engine.java class.
		\item ProcessFileTest.java: test class for the work done by the ProcessFiles.java class.
		\item ProduceRangesTest.java: test class for the work done by the ProduceRanges.java class.
		\item SystemStateTest.java: test class for how the state of the system changes when it is processing documents.
		\item TimelineDateTest.java: test class for the parsing of the TimelineDate.java class.
		\item ToJSONTest.java: test class for the production of JSONs for a list of Results.
	\end{enumerate}
	\item src/test/resources: Package for the test files used in the backend tests. Mostly used in ProcessFileTest.java, to test how well files are being processed.
	\begin{enumerate}
		\item testfile1.txt: sample text to test the extraction and processing of .txt documents. 
		\item testfile2.txt: sample text to test the extraction and processing of .txt documents. 
		\item testfile3.txt: sample text to test the extraction and processing of .txt documents. 
		\item testfile4.docx: sample text to test the extraction and processing of .docx documents. 
		\item testfile5.pdf: sample text to test the extraction and processing of .pdf documents. 
	\end{enumerate}
	\item build.gradle: a settings file for Gradle to set the location of the main method to run the system, the libraries required at different compilations (incluiding for tests), and define tasks for debugging.
\end{enumerate}

\section{Instructions}
Complete source code listings must be submitted as an appendix to the report. The project source codes are usually spread out over several files/units. You should try to help the reader to navigate through your source code by providing a ``table of contents'' (titles of these files/units and one line descriptions). The first page of the program listings folder must contain the following statement certifying the work as your own: ``I verify that I am the sole author of the programs contained in this folder, except where explicitly stated to the contrary''. Your (typed) signature and the date should follow this statement.

All work on programs must stop once the code is submitted to KEATS. You are required to keep safely several copies of this version of the program and you must use one of these copies in the project examination. Your examiners may ask to see the last-modified dates of your program files, and may ask you to demonstrate that the program files you use in the project examination are identical to the program files you have uploaded to KEATS. Any attempt to demonstrate code that is not included in your submitted source listings is an attempt to cheat; any such attempt will be reported to the KCL Misconduct Committee.

\textbf{You may find it easier to firstly generate a PDF of your source code using a text editor and then merge it to the end of your report. There are many free tools available that allow you to merge PDF files.}

\section{src/main/java/backend}
\subsection{Helpers}
\subsubsection{Sort.java}
\begin{lstlisting}
package backend.helpers;

import backend.process.Result;

import java.util.Collections;
import java.util.Comparator;
import java.util.List;

/**
 * Class that holds the sorting methods used throughout the project.
 */
public class Sort {
    /**
     * Sort the given List by their Date1 (in their TimelineDate)
     *
     * @param resultList the given List.
     * @return the given List sorted in ascending order.
     */
    public static List<Result> sortByDate1(List<Result> resultList) {
        Collections.sort(resultList, new Comparator<Result>() {//will sort in ascending order
            @Override
            public int compare(Result o1, Result o2) {
                if (o1.getTimelineDate().getDate1() != null && o2.getTimelineDate().getDate1() != null) {
                    return o1.getTimelineDate().getDate1().compareTo(o2.getTimelineDate()
                    		.getDate1());
                }
                if (o1.getTimelineDate().getDate1() == null) {
                    return -1;
                }
                return 1;
            }
        });
        return resultList;
    }
}
\end{lstlisting}
\subsubsection{ToJSON.java}
\begin{lstlisting}
package backend.helpers;

import backend.process.FileData;
import backend.process.Result;
import com.google.gson.*;

import java.lang.reflect.Type;
import java.util.List;

/**
 * Class used to turn a list of Results into a JSON String.
 */
public class ToJSON {
    /**
     * For the given List of Result objects, produce a JSON String of an array, where each index corresponds to one Result
     * in the list. Each Result is given by its range (date1, date2), its subjects (array of subjects), its event, and its
     * FileData (the filename, and base date used for processing the file), represented as a JsonObject from.
     * <p>
     * Whenever a value of data in a Result object is null, its corresponding key-pair will not be included in the final
     * JSON string. Such that, a completely empty Result object would be represented by: {subjects:[], event:"", from:{}}.
     *
     * @param results the given List of Result objects.
     * @return the JSON String representing the list of Result objects.
     */
    public static String toJSON(List<Result> results) {
        //sort the list in ascending order
        List<Result> sortedList = Sort.sortByDate1(results);
        String toReturn;
        final GsonBuilder gsonBuilder = new GsonBuilder();
        gsonBuilder.excludeFieldsWithoutExposeAnnotation();
        gsonBuilder.registerTypeAdapter(Result.class, new JsonSerializer<Result>() {
            @Override
            public JsonElement serialize(Result src, Type typeOfSrc, JsonSerializationContext context) {
                JsonObject jsonObject = new JsonObject();
                //adding the range dates (which can be null)
                //whenever a value is null, the key-pair will not be included in the final JSON
                jsonObject.addProperty("date1", src.getTimelineDate().getDate1FormattedDayMonthYear());
                jsonObject.addProperty("date2", src.getTimelineDate().getDate2FormattedDayMonthYear());
                //adding the subjects
                JsonArray subjectJsonArray = new JsonArray();
                for (String subject : src.getSubjects()) {
                    subjectJsonArray.add(subject);
                }
                jsonObject.add("subjects", subjectJsonArray);
                //adding the event
                jsonObject.addProperty("event", src.getEvent());
                //adding the file data (excluding the path, since this can be used on other system where files are elsewhere)
                JsonObject fromJsonObject = new JsonObject();
                FileData fileData = src.getFileData();
                if (fileData != null) {
                    fromJsonObject.addProperty("filename", fileData.getFileName());
                    fromJsonObject.addProperty("baseDate", fileData.getCreationDateFormattedDayMonthYear());
                }
                jsonObject.add("from", fromJsonObject);

                return jsonObject;
            }
        });
        //gsonBuilder.setPrettyPrinting();
        final Gson gson = gsonBuilder.create();
        toReturn = gson.toJson(sortedList);
        return toReturn;
    }
}
\end{lstlisting}
\subsubsection{ToPDF.java}
\begin{lstlisting}
package backend.helpers;

import backend.process.Result;
import org.apache.commons.lang3.text.WordUtils;
import org.apache.pdfbox.pdmodel.PDDocument;
import org.apache.pdfbox.pdmodel.PDPage;
import org.apache.pdfbox.pdmodel.PDPageContentStream;
import org.apache.pdfbox.pdmodel.common.PDRectangle;
import org.apache.pdfbox.pdmodel.font.PDType1Font;

import java.awt.*;
import java.io.File;
import java.io.IOException;
import java.util.List;

/**
 * Class used to represent a list of Results in a PDF.
 */
public class ToPDF {
    private static int widthOfRectangle = 275;
    private static int heightOfRectangle = 150;
    private static int strokeWidthOfLine = 2;
    private static int padding = 7;
    private static int fontSize = 15;
    private static int maxNoOfCharacters = 40;//max number of characters per line in the event box
    private int currentX = 0;
    private int currentY = 0;
    private float widthOfPage;
    private float heightOfPage;
    private PDDocument pdDocument;
    private PDPage currentPage;
    private PDPageContentStream contentStream;

    /**
     * Constructor called to prepare for the creation of the PDF.
     */
    public ToPDF() {
        pdDocument = new PDDocument();
        try {
            reset();
        } catch (IOException e) {
            e.printStackTrace();
        }
    }

    /**
     * Will set all the global variables, and close any open streams.
     *
     * @throws IOException when closing a stream that could be open or closed.
     */
    private void reset() throws IOException {
        currentPage = new PDPage(PDRectangle.A4);
        pdDocument.addPage(currentPage);
        if (contentStream != null) {
            contentStream.close();
            contentStream = null;
        }
        contentStream = new PDPageContentStream(pdDocument, currentPage);
        contentStream.setStrokingColor(Color.black);
        contentStream.setLineWidth(strokeWidthOfLine);

        widthOfPage = currentPage.getMediaBox().getWidth();
        heightOfPage = currentPage.getMediaBox().getHeight();

        currentX = 0;
        currentY = (int) heightOfPage;

        contentStream.moveTo(widthOfPage / 2, heightOfPage);
        contentStream.lineTo(widthOfPage / 2, 0);
        contentStream.stroke();
    }


    /**
     * For the given Results, produce a PDF File that has all the given Results displayed in a timeline (like on the
     * Application but without the buttons).
     *
     * @param results the list of given Results.
     * @param file    the File where we are storing the timeline.
     * @throws IOException due to working with streams.(I.e. trying to read or remove a file that is already open).
     */
    public void saveToPDF(List<Result> results, File file) throws IOException {
        List<Result> sortedList = Sort.sortByDate1(results);//sort the list in ascending order

        int counterOfEvents = 0;
        for (int i = 0; i < sortedList.size(); i++) {
            if (counterOfEvents >= 5) {//start a new page
                reset();
                counterOfEvents = 0;
            }
            if (i \% 2 == 0) {//even
                drawEvenEvent(sortedList.get(i), contentStream, i);
            } else {//odd
                drawOddEvent(sortedList.get(i), contentStream, i);
            }
            counterOfEvents++;
        }
        if (contentStream != null) {
            contentStream.close();
        }
        pdDocument.save(file);
        pdDocument.close();
    }

    /**
     * Called to draw the event in the timeline for a Result that is at a Odd index in the list.
     *
     * @param result        the given Result object for which to draw this event.
     * @param contentStream the stream to which we are drawing.
     * @param position      the position of the Result in the timeline (to show in the event box)
     * @throws IOException due to working with streams.
     */
    private void drawOddEvent(Result result, PDPageContentStream contentStream, int position) throws IOException {
        //initially y is the top right where this needs to be shown, x starts from the middle
        currentY -= padding; //add some padding to y
        currentX = (int) widthOfPage / 2;
        contentStream.moveTo(currentX, currentY);
        //write the text for the Event
        int lengthOfHorLine = (int) ((widthOfPage / 2) - (padding + widthOfRectangle));
        currentX += lengthOfHorLine;
        writeText(result, contentStream, currentX, position);
        //draw the rectangle to surround the text
        drawRectangle(contentStream, currentX, currentY - heightOfRectangle);
        //draw the horizontal line connecting event and timeline
        currentY -= heightOfRectangle / 2;
        contentStream.moveTo(currentX, currentY);
        contentStream.lineTo(widthOfPage / 2, currentY);
        contentStream.stroke();

        currentY -= (heightOfRectangle / 2) + padding;
    }

    /**
     * Called to draw the event in the timeline for a Result that is at a Even index in the list.
     *
     * @param result        the given Result object for which to draw this event.
     * @param contentStream the stream to which we are drawing.
     * @param position      the position of the Result in the timeline (to show in the event box)
     * @throws IOException due to working with streams.
     */
    private void drawEvenEvent(Result result, PDPageContentStream contentStream, int position) throws IOException {
        //initially x and y are top right in the page (0, pageHeight)
        //give some spacing between events
        currentY -= padding;//add some vertical small padding
        currentX = padding;
        contentStream.moveTo(currentX, currentY);
        //write the text for the Event
        writeText(result, contentStream, 0, position);
        //draw the rectangle to surround the text
        drawRectangle(contentStream, currentX, currentY - heightOfRectangle);
        //draw the horizontal line connecting event and timeline
        currentX += widthOfRectangle;
        currentY -= heightOfRectangle / 2;
        contentStream.moveTo(currentX, currentY);
        contentStream.lineTo(widthOfPage / 2, currentY);
        contentStream.stroke();

        currentY -= (heightOfRectangle / 2) + padding;//could set a string padding here

        //now connected the event (text surrounded by rectangle) to the timeline
    }

    /**
     * For the given Result and stream, write at the current x and y position (in new lines) the data held by the
     * Result object.
     *
     * @param result        the given Result.
     * @param contentStream the stream we are drawing to.
     * @param xOffset       an x-value offset from which the text is written from (starting at currentX)
     * @param position      the position of the Result in the timeline (to show in the event box)
     * @throws IOException due to using streams.
     */
    private void writeText(Result result, PDPageContentStream contentStream, int xOffset, int position) throws IOException {
        contentStream.beginText();
        contentStream.newLineAtOffset(2 * padding + xOffset, currentY - (fontSize + padding));//pad it horizontally and give vertical space for text
        contentStream.setFont(PDType1Font.TIMES_ROMAN, fontSize);
        contentStream.showText("Event #" + (position + 1));//show the position of the event (1st being the latest event)
        contentStream.newLineAtOffset(0, -(fontSize + padding));
        wrapText("Date: " + result.getTimelineDate(), contentStream);
        contentStream.newLineAtOffset(0, -(fontSize + padding));
        wrapText("Subjects: " + result.getSubjectsAsString(), contentStream);
        contentStream.newLineAtOffset(0, -(fontSize + padding));
        wrapText("Event: " + result.getEvent(), contentStream);
        contentStream.newLineAtOffset(0, -(fontSize + padding));
        wrapText("From: " + result.getFileData().getFileName() + " (" +
                result.getFileData().getCreationDateFormattedDayMonthYear() + ")", contentStream);
        contentStream.endText();
    }

    /**
     * For the given text, check whether it needs to be wrapped around to fit the constraint of the max number of
     * characters in a sentence. If the text needs to be wrapped, write each sentence with a small space in between (no
     * padding, that's to separate data), if it does not need to be wrapped write the text as it is.
     *
     * @param text          the given text.
     * @param contentStream where the text needs to be written on.
     * @throws IOException due to writing to streams.
     */
    private void wrapText(String text, PDPageContentStream contentStream) throws IOException {
        if (text.length() > maxNoOfCharacters) {//we have text worth wrapping
            String[] wrappedText = WordUtils.wrap(text, maxNoOfCharacters).split("\\r?\\n");
            for (int i = 0; i < wrappedText.length; i++) {
                contentStream.showText(wrappedText[i]);
                if (i < wrappedText.length - 1) {
                    contentStream.newLineAtOffset(0, -(fontSize));
                }
            }
        } else {//no point doing the work of wrapping text, when its below the threshold
            contentStream.showText(text);
        }
    }

    /**
     * For the given bottom left x and y co-ordinates, draw a rectangle (its width and height are given by the static final
     * member values).
     *
     * @param contentStream where the Rectangle is being drawn to.
     * @param bottomLeftX   the x-coordinate of the bottom left of the rectangle.
     * @param bottomLeftY   the y-coordinate of the bottom left of the rectangle
     * @throws IOException
     */
    private void drawRectangle(PDPageContentStream contentStream, int bottomLeftX, int bottomLeftY) throws IOException {
        contentStream.addRect(bottomLeftX, bottomLeftY, widthOfRectangle, heightOfRectangle);
        contentStream.stroke();
    }


}
\end{lstlisting}
\subsection{Process}
\subsubsection{CallbackResults.java}
\begin{lstlisting}
package backend.process;

import java.util.ArrayList;

/**
 * Interface for classes that call backend.process.ProcessFiles, so that when it finishes processing all Files passed in, the class (listener)
 * that called it can be informed.
 */
public interface CallbackResults {
    /**
     * Inform the Listener that all Files have been processed, and return the result of processing the Files.
     *
     * @param results the backend.process.Result objects obtained from processing all Files.
     * @param fileDataList the Data of the Files that Produced these results.
     */
    void gotResults(ArrayList<Result> results, ArrayList<FileData> fileDataList);
}
\end{lstlisting}
\subsubsection{Engine.java}
\begin{lstlisting}
package backend.process;

import backend.system.BackEndSystem;
import edu.stanford.nlp.ling.CoreAnnotations;
import edu.stanford.nlp.ling.Label;
import edu.stanford.nlp.pipeline.Annotation;
import edu.stanford.nlp.pipeline.StanfordCoreNLP;
import edu.stanford.nlp.semgraph.SemanticGraph;
import edu.stanford.nlp.semgraph.SemanticGraphCoreAnnotations;
import edu.stanford.nlp.semgraph.SemanticGraphEdge;
import edu.stanford.nlp.trees.Tree;
import edu.stanford.nlp.trees.TreeCoreAnnotations;
import edu.stanford.nlp.util.CoreMap;
import edu.stanford.nlp.util.Pair;

import java.util.ArrayList;
import java.util.List;

/**
 * backend.process.Engine class that text in text as input and produces a list of backend.process.Result objects, which are events depicted in the text. An event is only picked out, if it has a date
 * associated with it.
 * <p>
 * Uses the algorithm proposed in: Bonnie Dorr, David Zajic and Richard Schwartz. Hedge Trimmer: A Parse-and-Trim Approach to Headline
 * Generation. Proceedings of the HLT-NAACL 03 on Text summarization Workshop-Volume 5. Association for
 * Computational Linguistics, pp. 1–8.
 */
public class Engine {
    private static int threshold;
    private StanfordCoreNLP coreNLP;
    private String baseDate;

    /**
     * Set up the StanfordCoreNLP to analyze text.
     */
    public Engine() {
        coreNLP = BackEndSystem.getInstance().getCoreNLP();//can have it like before as the models will already be  loaded, but this avoids having to check that
        threshold = BackEndSystem.getInstance().getSettings().getThresholdSummary();
        System.out.println("Using a threshold value of: " + threshold);
    }

    /**
     * Produces a list of Results based on the text passed in. (Determine events in the text).
     *
     * @param input The text for which we want to produce events for.
     * @param date  The base date, from which we can determine exact dates from relative dates (eg Yesterday).
     * @return list of Results produced from events depicted in the text passed in, using the base date.
     */
    public ArrayList<Result> getResults(String input, String date) {
        ArrayList<Result> results = new ArrayList<>();
        baseDate = date;
        System.out.println("Base Date: " + baseDate);
        Annotation annotation;
        annotation = new Annotation(input);
        annotation.set(CoreAnnotations.DocDateAnnotation.class, date);//setting a reference so that when it finds a normalazied entity tag that isnt complete will determine it
        coreNLP.annotate(annotation);
        //coreNLP.prettyPrint(annotation, new PrintWriter(System.out));

        for (CoreMap sentence : annotation.get(CoreAnnotations.SentencesAnnotation.class)) {
            System.out.println(sentence);
            Result result = getResult(sentence);
            if (result != null) {
                results.add(result);
            }
        }
        return results;
    }

    /**
     * Returns a backend.process.Result object if the sentence contains a date, else returns null.
     * The backend.process.Result object will include a list of subjects as well as a summary (or the entire text, depending on the
     * length of text) of the sentence.
     *
     * @param sentence the CoreMap that contains the Sentence we want to produce a backend.process.Result for (if it has a date)
     * @return a backend.process.Result object if the sentence contained a Date; null instead.
     */
    private Result getResult(CoreMap sentence) {
        Result result = new Result();
        result.setOriginalString(sentence.toString());
        setDatesAndSubjectsNET(sentence, result);

        if (result.getDates().size() > 0) {//we have found dates, so lets find more subjects and the event of the sentence
            //setGrammaticalSubjects(sentence,result);//setting grammatical subjects in the result object
            setEvent(sentence, result);//set the summarized sentence as the event depicted in the sentence
            return result;
        }
        return null;
    }

    /**
     * Set the Dates and Subjects for the backend.process.Result based on Named-Entity Tags from the CoreMap passed in.
     * Dates have a DATE Named-Entity Tag. Subjects are LOCATIONs, ORGANIZATIONs, PERSONs, and MONEY.
     *
     * @param sentence a CoreMap that holds the sentence we want to extract data from.
     * @param result   the backend.process.Result object that we are determining the data for.
     */
    private void setDatesAndSubjectsNET(CoreMap sentence, Result result) {
        for (CoreMap mention : sentence.get(CoreAnnotations.MentionsAnnotation.class)) {
            String namedEntityTag = mention.get(CoreAnnotations.NamedEntityTagAnnotation.class);
            System.out.println(namedEntityTag + ": " + mention.get(CoreAnnotations.TextAnnotation.class));

            if (namedEntityTag.equals("DATE")) {
                //found a date for the result object
                String date_1 = mention.get(CoreAnnotations.NormalizedNamedEntityTagAnnotation.class);
                System.out.println("About to print time for the sentence: " + sentence);
                System.out.println("Normalized entity tag: " + date_1);
                String date = mention.get(CoreAnnotations.TextAnnotation.class);
                System.out.println("We are storing date: " + date);
                result.addDate(date);
                result.addDate_1(date_1, baseDate);

            } else if (namedEntityTag.equals("LOCATION") || namedEntityTag.equals("ORGANIZATION") ||
                    namedEntityTag.equals("PERSON") || namedEntityTag.equals("MONEY") || namedEntityTag.equals("MISC")) {
                //found a subject for the result object
                String subject = mention.get(CoreAnnotations.TextAnnotation.class);
                result.addSubject(subject);
            }
        }
    }

    /**
     * Finds the subjects of the sentence based on grammatical structure.
     * Gets the Basic Dependencies in the sentence and for each relation that contains the name "subj" (eg nsubj),
     * store its dependent in the backend.process.Result object, this is the grammatical subject of a sentence.
     *
     * @param sentence a CoreMap that holds the sentence we want to extract data from.
     * @param result   the backend.process.Result object that we are determining the grammatical subject for.
     */
    private void setGrammaticalSubjects(CoreMap sentence, Result result) {
        SemanticGraph semanticGraph = sentence.get(SemanticGraphCoreAnnotations.BasicDependenciesAnnotation.class);
        for (SemanticGraphEdge semanticGraphEdge : semanticGraph.edgeIterable()) {
            if (semanticGraphEdge.getRelation().getShortName().contains("subj")) {
                result.addSubject(semanticGraphEdge.getDependent().value());
            }
        }
    }

    /**
     * For the sentence passed in, summarize it and set it as the event of the backend.process.Result object.
     * <p>
     * Following the algorithm proposed in: Bonnie Dorr, David Zajic and Richard Schwartz. Hedge Trimmer: A
     * Parse-and-Trim Approach to Headline Generation. Proceedings of the HLT-NAACL 03 on Text summarization
     * Workshop-Volume 5. Association for Computational Linguistics, pp. 1–8.
     *
     * @param sentence a CoreMap that holds the sentence we want to summarize to get the event for.
     */
    private void setEvent(CoreMap sentence, Result result) {
        Tree tree = sentence.get(TreeCoreAnnotations.TreeAnnotation.class);
        tree = getLeftmostLowestS(tree);//get leftmost-lowest S
        //remove time expressions
        removeTimeExpressions(tree, result);
        removeDeterminers(tree);//remove determiners
        xpOverXP(tree);//apply XP-over-XP rule
        xpBeforeNP(tree);//apply removal of XPs before NP rule
        cleanUp(tree);//remove any punctuation that could be left over
        tree = lastShorten(tree);//shorten the tree with the last two rules, removePPs and removeSBARs
        String event = produceString(tree);
        result.setEvent(event);
    }

    /**
     * Get the LeftmostLowest S subtree or the root of the tree if it doesn't exist.
     *
     * @param tree the Grammatical Structure of the sentence we are summarizing.
     * @return the leftmost-lowest S found in the passed in tree, can be the same tree.
     */
    private Tree getLeftmostLowestS(Tree tree) {
        //post-order to find first S
        Tree firstS = null;
        for (Tree node : tree.postOrderNodeList()) {
            if (!node.isLeaf() && node.value().equals("S")) {
                firstS = node;
                break;
            }
        }

        //if we haven't found a lowest S, then we use the root of the tree; else we just return the lowest S
        return (firstS == null) ? tree : firstS;
    }

    /**
     * Remove determiners like 'a' and 'the' from the tree.
     * Determiners are always the first child of their parent.
     *
     * @param tree the Grammatical Structure of the sentence we are summarizing.
     */
    private void removeDeterminers(Tree tree) {
        for (Tree node : tree.preOrderNodeList()) {//loop over in pre order
            if (!node.isLeaf()) {//if we have children
                Tree firstChild = node.children()[0];
                if (!firstChild.isLeaf() && firstChild.value().equals("DT")) {//if it has children and it has the determiner tag
                    Tree determiner = firstChild.children()[0];//determiners only have one child
                    if (determiner.value().equals("a") || determiner.value().equals("the")) {//if our determiner is of the type
                        node.removeChild(0);    // 'a' or 'the' then we delete it, remove the DT node that contains the 'a' or 'the'
                    }
                }
            }
        }
    }

    /**
     * Go through the tree in pre-order, finding the last appearance of an xp over an xp (so check for np-np,vp-vp,and s-s),
     * then remove all children of the outer xp, except for the first first child.
     * Repeat if we are still above the threshold.
     *
     * @param tree the Grammatical Structure of the sentence we are summmarizing
     */
    private void xpOverXP(Tree tree) {
        if (tree.yield(new ArrayList<Label>()).size() > threshold) {//if we are above threshold, we need to reduce tree
            Tree toRemoveChild = null;
            for (Tree node : tree.preOrderNodeList()) {//need to loop over tree
                if (!node.isLeaf() && node.children().length > 1) {//if found a node
                    Tree possibleXpNode = node.children()[0];//who is of value XP
                    if (node.value().equals(possibleXpNode.value()) && // and its first child is also XP, then need to
                            ((node.value().equals("NP") && possibleXpNode.value().equals("NP")) //record it to delete its
                                    || (node.value().equals("VP") && possibleXpNode.value().equals("VP")) //other children
                                    || (node.value().equals("S") && possibleXpNode.value().equals("S")))) {
                        toRemoveChild = node;
                    }
                }
            }
            if (toRemoveChild != null) {//we have a node to delete all its children except first from
                while (toRemoveChild.children().length > 1) {
                    toRemoveChild.removeChild(1);
                }
                xpOverXP(tree);
            }
        }
    }

    /**
     * If the tree size is greater than the threshold, then delete any XP (PP,NP,VP) before the subject of the sentence
     * which is the NP child of S.
     * Find NP subject of S, then search through the tree in pre order:
     * if we reach the NP of S then stop
     * if we reach an XP that isn't the NP of S (the subject of the sentence), then store its parent and the child index
     * that it is in (in the parents children list), to later remove it from the tree.
     *
     * @param tree Grammatical structure of the sentence we are summarizing.
     */
    private void xpBeforeNP(Tree tree) {
        if (tree.yield(new ArrayList<Label>()).size() > threshold && tree.value().equals("S")) {
            Tree[] childrenOfS = tree.children();
            Tree pointsToFirstNP = null;
            for (int i = 0; i < childrenOfS.length; i++) {
                Tree child = childrenOfS[i];
                if (child.value().equals("NP")) {//found subject of sentence
                    pointsToFirstNP = child;//remember it
                    break;//no point in continue to loop, just wanted to find the NP child of S
                }
            }
            if (pointsToFirstNP != null) {
                Tree toDeleteChild = null;//hold parent of node we have to delete
                int childIndexToDelete = 0;//hold index of child of parent that we have to delete
                for (Tree node : tree.preOrderNodeList()) {
                    if (node == pointsToFirstNP || toDeleteChild != null) {
                        break;
                    } else if (!node.isLeaf()) {
                        Tree[] children = node.children();
                        for (int i = 0; i < children.length; i++) {
                            Tree child = children[i];
                            if (child == pointsToFirstNP) {
                                break;
                            } else if (child.value().equals("VP") || child.value().equals("NP") || child.value().equals("PP")) {
                                toDeleteChild = node;
                                childIndexToDelete = i;
                            }
                        }
                    }
                }
                if (toDeleteChild != null) {
                    toDeleteChild.removeChild(childIndexToDelete);
                    //xpBeforeNP(tree); ? recursively call
                }
            }
        }
    }

    /**
     * Applies the last two rules for iterative shortening. First will try to remove the rightmost-lowest PP until it
     * has reached the threshold. If it has removed all the PPs, or it can't delete anymore and it is still below the
     * threshold, then undo the changes and do the below.
     * Remove the SBARs until we are below the threshold, or there are no more left. If we are still above the threshold,
     * then delete all the PPs until we are below the threshold, or there are no more left.
     * <p>
     * According to the algorithm proposed in: Bonnie Dorr, David Zajic and Richard Schwartz. Hedge Trimmer: A Parse-and-Trim Approach to Headline
     * Generation. Proceedings of the HLT-NAACL 03 on Text summarization Workshop-Volume 5. Association for
     * Computational Linguistics, pp. 1–8.
     *
     * @return a tree that has gone through the shortening of size smaller than threshold, or a tree where nothing can be removed
     */
    private Tree lastShorten(Tree tree) {
        if (tree.yield(new ArrayList<Label>()).size() > threshold) {
            Tree copyOfTree = tree.deepCopy();
            //removePPs(copyOfTree);
            removeXs(tree, "PP");
            if (copyOfTree.yield(new ArrayList<Label>()).size() > threshold) {//we undo what we did previously and do SBAR removal and then PP
                //removeSBARs(tree);//need to remove SBARs
                removeXs(tree, "SBAR");
                //removePPs(tree);//then remove PPs according to rule
                removeXs(tree, "PP");
            } else {//just the PP removal was enough to reduce the size below the threshold so we return the tree which has undergone that removal
                return copyOfTree;
            }
        }
        return tree;
    }

    /**
     * Removes trailing Xs if we are below the threshold.
     *
     * @param tree the Tree that holds the grammatical structure of the text we are summarizing.
     * @param x    the type of X we are removing, eg: PP, SBAR, from the tree.
     */
    private void removeXs(Tree tree, String x) {
        if (tree.yield(new ArrayList<Label>()).size() > threshold) {
            Tree toDeleteChild = null;
            int childIndexToDelete = 0;
            for (Tree node : tree.preOrderNodeList()) {
                if (!node.isLeaf()) {
                    Tree[] children = node.children();
                    for (int i = 0; i < children.length; i++) {
                        if (children[i].value().equals(x)) {
                            toDeleteChild = node;
                            childIndexToDelete = i;
                        }
                    }
                }
            }
            if (toDeleteChild != null) {
                toDeleteChild.removeChild(childIndexToDelete);
                removeXs(tree, x);
            }
        }
    }

    /**
     * Produce a string based on the leaf nodes in the tree.
     *
     * @param tree the Tree that has the grammatical structure of the sentence we want to produce.
     * @return the String that follows the grammatical structure passed in.
     */
    private String produceString(Tree tree) {
        String toReturn = "";
        List<Tree> preOrderList = tree.preOrderNodeList();
        for (int i = 0; i < preOrderList.size(); i++) {
            Tree node = preOrderList.get(i);
            if (node.isLeaf()) {
                toReturn += node.value();
                if (i != preOrderList.size() - 1) {
                    toReturn += " ";
                }
            }
        }
        return toReturn;
    }

    /**
     * Called to remove punctuation that is at the start of the tree
     *
     * @param tree the Tree from which we are removing the punctuation
     */
    private void cleanUp(Tree tree) {//if its the first character, than its the first child of root S
        if (!tree.isLeaf() && tree.value().equals("S") && tree.children()[0].value().equals(",")) {
            tree.removeChild(0);
        }
    }

    /**
     * Remove time expressions in the tree.
     * According to the algorithm proposed in: Bonnie Dorr, David Zajic and Richard Schwartz. Hedge Trimmer: A Parse-and-Trim Approach to Headline
     * Generation. Proceedings of the HLT-NAACL 03 on Text summarization Workshop-Volume 5. Association for
     * Computational Linguistics, pp. 1–8.
     *
     * @param tree   the Tree that contains the Grammatical Structure of the sentence we are summarizing.
     * @param result the backend.process.Result object that contains previously found time expressions in the original sentence.
     */
    private void removeTimeExpressions(Tree tree, Result result) {
        /*go over the tree, for every PP check if it has a NP child, if so check if that NP contains the time expression, if so delete the PP.
        * afterwards delete any NP that contains the time expression
        * */
        //System.out.println("Tree going in:");
        //tree.pennPrint();
        ArrayList<Pair<Tree, Integer>> toDeleteNodes = new ArrayList<>();
        for (Tree node : tree.preOrderNodeList()) {
            if (!node.isLeaf()) {
                Tree[] children = node.children();
                for (int i = 0; i < children.length; i++) {
                    if (children[i].value().equals("PP") && !children[i].isLeaf()) {
                        Tree[] childOfChildren = children[i].children();
                        for (int j = 0; j < childOfChildren.length; j++) {
                            if (childOfChildren[j].value().equals("NP") && !childOfChildren[j].isLeaf()) {
                                String childOfChildrenString = produceString(childOfChildren[j]);
                                if (result.hasDate(childOfChildrenString)) {
                                    System.out.println("Found a match");
                                    toDeleteNodes.add(new Pair<>(node, i));
                                    break;
                                }
                            }
                        }
                    }
                }
            }
        }
        //now delete
        for (Pair<Tree, Integer> pair : toDeleteNodes) {
            if (pair.first.getChild(pair.second).value().equals("PP")) {
                System.out.println("Removing: " + pair.first.getChild(pair.second));
                pair.first.removeChild(pair.second);
            }
        }
        //tree.pennPrint();
        //clear toDeleteNodes
        toDeleteNodes.clear();
        //now find any NP with the time expression
        for (Tree node : tree.preOrderNodeList()) {
            if (!node.isLeaf()) {
                Tree[] children = node.children();
                for (int i = 0; i < children.length; i++) {
                    if (children[i].value().equals("NP") && !children[i].isLeaf()) {
                        String childString = produceString(children[i]);
                        if (result.hasDate(childString)) {
                            System.out.println("Found a match");
                            toDeleteNodes.add(new Pair<>(node, i));
                        }
                    }
                }
            }
        }
        for (Pair<Tree, Integer> pair : toDeleteNodes) {
            if (pair.first.getChild(pair.second).value().equals("NP")) {
                System.out.println("Removing: " + pair.first.getChild(pair.second));
                pair.first.removeChild(pair.second);
            }
        }
        //System.out.println("Tree afterwards:");
        //tree.pennPrint();

    }

}

\end{lstlisting}
\subsubsection{FileData.java}
\begin{lstlisting}
package backend.process;

import java.io.File;
import java.io.IOException;
import java.nio.file.Files;
import java.nio.file.Path;
import java.nio.file.Paths;
import java.nio.file.attribute.BasicFileAttributes;
import java.nio.file.attribute.FileTime;
import java.text.ParseException;
import java.text.SimpleDateFormat;
import java.time.LocalDate;
import java.util.Date;

/**
 * Class that holds the relevant data (for this project) of a File. In this case, the File's name and path in the System.
 */
public class FileData implements Comparable<FileData> {
    private final static SimpleDateFormat inputSimpleDateFormat = new SimpleDateFormat("dd-MM-yyyy");
    private final static SimpleDateFormat outputSimpleDateFormat = new SimpleDateFormat("yyyy-MM-dd");
    private final static String epochDateFormatted = "1970-01-01";
    private String fileName;
    private String filePath;
    private Date creationDate;

    /**
     * Used to create a FileData object, by providing its name and path.
     *
     * @param fileName the name of the File.
     * @param filePath the path of the File.
     */
    public FileData(String fileName, String filePath) {
        this.fileName = fileName;
        this.filePath = filePath;
    }

    public FileData(File file) {
        if (file.exists() && file.isFile()) {
            this.fileName = file.getName();
            this.filePath = file.getAbsolutePath();
            this.creationDate = getCreationDate(file);//set the creation date for this file
        }
    }

    /**
     * Get the name of the File that produced this FileData.
     *
     * @return the name of the File.
     */
    public String getFileName() {
        return fileName;
    }

    /**
     * Set the name of the File that produced this FileData.
     *
     * @param fileName the name of the File.
     */
    public void setFileName(String fileName) {
        this.fileName = fileName;
    }

    /**
     * Get the path of the File that produced this FileData.
     *
     * @return the path of the File.
     */
    public String getFilePath() {
        return filePath;
    }

    /**
     * Set the path of the File that produced this FileData.
     *
     * @param filePath the path of the File.
     */
    public void setFilePath(String filePath) {
        this.filePath = filePath;
    }


    /**
     * Print the information of this FileData.
     *
     * @return a String with the name and path of the File being represented.
     */
    @Override
    public String toString() {
        return String.format("Name: \%s, Path: \%s", fileName, filePath);
    }

    /**
     * Checks whether the given input is equal to this object.
     *
     * @param obj the Object to check equality with this object.
     * @return true if obj is a FileData object with the same data as this object; false otherwise.
     */
    @Override
    public boolean equals(Object obj) {
        try {
            FileData otherFileData = (FileData) obj;
            return fileName.equals(otherFileData.fileName) && filePath.equals(otherFileData.filePath) && creationDate.equals(otherFileData.creationDate);
        } catch (Exception e) {
            return false;
        }
    }

    /**
     * The creationDate (or baseDate) of the File this is representing.
     *
     * @return the creationDate (or baseDate) of the File as a Date object.
     */
    public Date getCreationDate() {
        return creationDate;
    }

    /**
     * Set the creation date as a String for the File this FileData is representing.
     *
     * @param creationDate the creation date as a String input of type: dd-MM-yyyy.
     */
    public void setCreationDate(String creationDate) {
        System.out.println("Set Creation Date: " + creationDate);
        try {
            this.creationDate = inputSimpleDateFormat.parse(creationDate);
        } catch (ParseException e) {
            e.printStackTrace();
        }
    }

    /**
     * Get the creation date in the format of day, followed by month, followed by year.
     *
     * @return the creation date in the format: dd-MM-yyyy
     */
    public String getCreationDateFormattedDayMonthYear() {
        return (creationDate != null) ? inputSimpleDateFormat.format(creationDate) : null;
    }

    /**
     * Get the creation date in the format of year, followed by month, followed by day.
     *
     * @return the creation date in the format: yyyy-MM-dd.
     */
    public String getCreationDateFormattedYearMonthDay() {
        return outputSimpleDateFormat.format(creationDate);
    }

    /**
     * Used to get the creation date of a input File, by looking at the attributes associated with the File.
     *
     * @param file the input File.
     * @return the creation date, or the current date now if no creation date is available, as a Date object.
     */
    private static Date getCreationDate(File file) {
        String creationDate = LocalDate.now().toString();//LocalDate.now gives you the Date for the current moment (default), in the format yyyy-MM-dd
        Date toReturn = null;
        //try and get file creation date
        Path filePath = Paths.get(file.getAbsolutePath());//only way to use BasicFileAttributes is to pass in a Path,
        try {//so get the Path for the passed in File
            BasicFileAttributes basicFileAttributes = Files.readAttributes(filePath, BasicFileAttributes.class);//creation date of a File is a basic file attribute
            FileTime creationTime = basicFileAttributes.creationTime();//can be epoch time if it does not exist, don't set it in that case
            String possibleDate = inputSimpleDateFormat.format(creationTime.toMillis());//need to check its not epoch time, as else the creation is not valid for this file
            if (!possibleDate.equals(epochDateFormatted)) {//epoch time date is given if it cant find a creation date
                creationDate = possibleDate;
            }
        } catch (IOException e) {
            e.printStackTrace();
        }
        try {
            toReturn = inputSimpleDateFormat.parse(creationDate);
        } catch (ParseException e) {
            e.printStackTrace();
        }
        return toReturn;
    }


    /**
     * Used to compare a given FileData to this FileData, by looking at their filename (to sort by their name, in
     * ascending order).
     *
     * @param o the other FileData object we are comparing to.
     * @return the result of comparing the file name of this FileData with the filename of the other FileData.
     */
    @Override
    public int compareTo(FileData o) {
        return fileName.compareTo(o.fileName);
    }
}
\end{lstlisting}
\subsubsection{ProcessFileCallback.java}
\begin{lstlisting}
package backend.process;

import java.util.ArrayList;

/**
 * Interface to be implemented, to inform the Listener of ProcessFile Thread of when a file has been processed.
 */
public interface ProcessFileCallback {
    /**
     * Inform the Listener that the passed in File to ProcessFile has been processed, and pass the Results of processing
     * the given File.
     *
     * @param results the Results from processing the given File passed into the ProcessFile Thread.
     * @param fileData the File Data of the File that produced these Results.
     */
    void callBack(ArrayList<Result> results, FileData fileData);
}
\end{lstlisting}
\subsubsection{ProcessFiles.java}
\begin{lstlisting}
package backend.process;

import backend.system.BackEndSystem;
import backend.system.SystemState;
import org.apache.pdfbox.pdmodel.PDDocument;
import org.apache.pdfbox.text.PDFTextStripper;
import org.apache.poi.xwpf.extractor.XWPFWordExtractor;
import org.apache.poi.xwpf.usermodel.XWPFDocument;

import java.io.File;
import java.io.FileInputStream;
import java.io.FileNotFoundException;
import java.io.IOException;
import java.util.ArrayList;
import java.util.List;
import java.util.Scanner;
import java.util.concurrent.Semaphore;


/**
 * Handles the parsing of files, and the multi-threading of the backend.process.Engine.
 */
public class ProcessFiles implements ProcessFileCallback {
    private static int maxNoOfThreads;//this value is determined by the settings of the System (by default 2)
    private Semaphore semaphore;//will use the maxNoOfThreads to control how many Threads will be running in parallel
    private Semaphore semaphoreFinished = new Semaphore(0);//so that we wait until all threads finish
    private ArrayList<Result> results = new ArrayList<>();
    private int filesToGo;//to notify the listener when it is done

    /**
     * For the list of Files passed in, it will process each in separate threads, while respecting the maximum number
     * of threads that can run at any given time (i.e. maximum 2 extra threads running at a time).
     * Will call gotResults() on the backend.process.CallbackResults object when all files have been processed, and a list of Results has
     * been produced.
     * <p>
     * <b>Should be the only thing called in this class by its Users.</b>
     *
     * @param files     the list of File objects that contain text that needs to be processed (atm only processes .docx/.pdf/.txt files)
     * @param fileDatas the list of FileData objects, where the index corresponds the File in the same index in the files list, that
     *                  contains the needed data for each File.
     */
    public List<Result> processFiles(List<File> files, List<FileData> fileDatas) {
        maxNoOfThreads = BackEndSystem.getInstance().getSettings().getMaxNoOfThreads();//get the Settings value
        semaphore = new Semaphore(maxNoOfThreads);//set the Max number of Threads that can run in parallel
        //should only run if we are not Processing
        //this will also set up the StanfordCoreNLP (when GUI is implemented, it will already by set up, as it will be the first thing ran)
        System.out.println("Will try to run, with maxNoOfThreads: " + maxNoOfThreads + " and available permits: " + semaphore.availablePermits());
        if (BackEndSystem.getInstance().getSystemState() != SystemState.PROCESSING && files.size() == fileDatas.size()) {//if we arent processing, then we can begin to do that
            System.out.println("Is running");
            filesToGo = files.size();//and when we need to call
            BackEndSystem.getInstance().setSystemState(SystemState.PROCESSING);
            System.out.println("In Thread: " + Thread.currentThread().toString());
            for (int i = 0; i < files.size(); i++) {
                File file = files.get(i);
                FileData fileData = fileDatas.get(i);//should be the same
                //check they are the same?
                //acquire from the semaphore
                try {
                    System.out.println("Trying to acquire semaphore for file: " + file);
                    semaphore.acquire();//will wait if there is already maxnoofthreads running, until one finishes: then it gets to run
                    System.out.println("Acquired semaphore for file: " + file);
                    //process file
                    Thread thread = new ProcessFile(file, this, fileData);//pass a reference so that the thread can call this when it finishes processing the file
                    thread.start();//start processing this file(get its text and pass it to the backend.process.Engine)
                } catch (InterruptedException e) {
                    e.printStackTrace();
                    //should release semaphore and reduce filesToGo count if the Thread is interrupted
                }
            }
            try {
                System.out.println("Going to wait until the last files are processed.");
                semaphoreFinished.acquire();
                System.out.println("The last files have been processed.");
            } catch (InterruptedException e) {
                e.printStackTrace();
            }
            return results;
        }
        return null;
    }

    /**
     * Called when the ProcessFile Thread has finished processing a file. It is synchronized to avoid two separate threads
     * trying to add to the list of Results at the same time.
     * Enforces the rule of releasing the semaphore, as this Thread finished running, to then allow waiting Threads to run.
     * Will inform the CallbackResult object when it finishes processing all Files.
     *
     * @param results  the backend.process.Result objects produced by Processing the given file in the Thread.
     * @param fileData the data of the File that produced these Results
     */
    public synchronized void callBack(ArrayList<Result> results, FileData fileData) {
        //we finished processing a file
        filesToGo--;//one less to look at
        System.out.println("Files to go: " + filesToGo);
        //add the results to the list held
        this.results.addAll(results);
        //release semaphore
        System.out.println("Released semaphore from Thread: " + Thread.currentThread().toString());
        semaphore.release();
        //check if we have processed everything, if so release the finished semaphore
        if (filesToGo == 0) {
            //has processed
            BackEndSystem.getInstance().setSystemState(SystemState.PROCESSED);
            //has returned the results so we finished
            BackEndSystem.getInstance().setSystemState(SystemState.FINISHED);
            semaphoreFinished.release();//value is now 1, so the thread that was acquiring can continue
        }
    }

    /**
     * Called to return the text of a given File.
     *
     * @param file the given File.
     * @return the text of the given File.
     */
    public String getTextInFile(File file) {
        return new ProcessFile(file, null, null).getText(file);
    }

    /**
     * In charge of Processing just one File in a separate Thread.
     */
    private static class ProcessFile extends Thread {
        //Need to release even if it messes up
        File file;
        ProcessFileCallback processFileCallback;
        FileData fileData;

        /**
         * Create a ProcessFile object that holds the data needed: the File to process and who to callback when the
         * backend.process.Engine finishes processing.
         *
         * @param file                the File to process.
         * @param processFileCallback who to inform when the backend.process.Engine finished processing the given file.
         */
        ProcessFile(File file, ProcessFileCallback processFileCallback, FileData fileData) {//hold sempahore
            this.file = file;
            this.processFileCallback = processFileCallback;
            this.fileData = fileData;
        }

        /**
         * What starts running on a separate Thread. Will first get the Text for the given file, and then pass it to the
         * backend.process.Engine, which will return a list of Results that is passed to the backend.process.ProcessFileCallback.
         */
        @Override
        public void run() {
            super.run();
            System.out.println("For: " + file + " in Thread: " + Thread.currentThread().toString());//for logging purposes
            ArrayList<Result> toReturnResults = new ArrayList<>();//initially no results
            //check file exists in system
            if (fileExists(file)) {
                //fileData = new FileData(file.getName(), file.getAbsolutePath());
                //get the text for that file
                String toProcess = getText(file);//will get the text for the file considering its extension
                //run engine on this
                if (!toProcess.equals("")) {//if we actually have text to process, don't waste time attempting to process else
                    String baseDate = fileData.getCreationDateFormattedYearMonthDay();//since we will need to process, get a base date to use
                    System.out.println("Base Date for " + file.getName() + " is: " + baseDate);
                    toReturnResults = new Engine().getResults(toProcess, baseDate);//pass in file data, so each result holds it
                    addFileData(fileData, toReturnResults);
                }
            }
            //call the backend.process.ProcessFileCallback that we finished processing and return the results of processing that one File.
            processFileCallback.callBack(toReturnResults, fileData);
        }

        /**
         * For the list passed in, set for all of them the given FileData.
         *
         * @param fileData the FileData to be set to all the Results passed in.
         * @param results  the Results for which the FileData needs to be set.
         */
        private void addFileData(FileData fileData, ArrayList<Result> results) {
            for (Result result : results) {
                result.setFileData(fileData);
            }
        }

        /**
         * Check that we can actually use the given File.
         *
         * @param file the File that we can check that we can read it.
         * @return true if  it is indeed a File, that it still exists in the system, and that we can read data from it, have access for that.
         */
        private boolean fileExists(File file) {
            return file.exists() && file.isFile() && file.canRead();
        }

        /**
         * For the given File, perform the operations needed to obtain its text, as different file extension have different
         * encodings.
         *
         * @param file the File for which we need to get text from.
         * @return the text from the File, or empty text if it was not possible to get text from the File.
         */
        private String getText(File file) {
            String fileName = file.getName();//file should have format: name.extension
            int lastDotPosition = fileName.lastIndexOf(".");//to separate name from the extension
            if (lastDotPosition != -1) {//-1 if . was no where in the file name
                String fileExtension = fileName.substring(lastDotPosition);//everything after the last dot, the extension
                System.out.println("File name: " + fileName + " file extension: " + fileExtension);
                switch (fileExtension) {
                    case ".pdf"://file has a .pdf extension
                        //call method to get string from pdf
                        return getTextPDF(file);
                    case ".txt"://file has a .txt extension
                        //call method to get string from txt
                        return getTextTXT(file);
                    case ".docx":
                        return getTextDoc(file);
                    default://could hold list of files that could not be processed
                        //get string from default process (open file and use buffered reader)
                        break;
                }
            }
            return "";
        }


        /**
         * Get the text for a PDF File using the appropriate library (Apache PDFBox).
         *
         * @param file PDF File to get text from.
         * @return the text in the PDF File, or an empty String if it was not possible.
         */
        private String getTextPDF(File file) {
            String toReturn = "";//base text, if it fails we just return empty text
            try {
                PDDocument pdDocument = PDDocument.load(file);//create Document that has processed the bytes in the pdf file
                PDFTextStripper pdfTextStripper = new PDFTextStripper();
                //pdfTextStripper.setSortByPosition(true);//in the case the program that created the page, didnt place the text in the order it is shown (so could read text in wrong order)
                toReturn = pdfTextStripper.getText(pdDocument).replaceAll(pdfTextStripper
                	.getLineSeparator(), "");//to then get its text
                System.out.println("Line Seperator: " + pdfTextStripper.getLineSeparator());
                pdDocument.close();//always remember to close the stream, or document in this case
            } catch (IOException e) {
                e.printStackTrace();
            }
            //System.out.println("For file:"+file.getName()+" has text: "+toReturn);
            return toReturn;
        }

        /**
         * Get the text for a TXT File, assuming the enconding is UTF-8(later determine encoding and pass it to the Scanner
         * that processes the File).
         *
         * @param file a TXT File to get the text from.
         * @return the text in the TXT File, or an empty String if it was not possible.
         */
        private String getTextTXT(File file) {
            //should determine character set, can look at https://code.google.com/archive/p/juniversalchardet/
            String toReturn = "";//base text to return, empty String
            try {
                Scanner scanner = new Scanner(file, "UTF-8");//uses the default character set UTF-8
                while (scanner.hasNextLine()) {//while we can read more
                    toReturn += scanner.nextLine();
                }
            } catch (FileNotFoundException e) {
                e.printStackTrace();
            }
            return toReturn;
        }

        /**
         * Process a DOCX Document using the appropriate library (Apache POI).
         *
         * @param file the DOCX File to get the text from.
         * @return the text in the File, or an empty String if it was not possible to process the File.
         */
        private String getTextDoc(File file) {
            String toReturn = "";//base text to return, empty String
            try {
                XWPFDocument xwpfDocument = new XWPFDocument(new FileInputStream(file));//used to process docx documents
                XWPFWordExtractor xwpfWordExtractor = new XWPFWordExtractor(xwpfDocument);
                toReturn = xwpfWordExtractor.getText();
            } catch (IOException e) {
                e.printStackTrace();
            }
            //System.out.println("For file:"+file.getName()+" has text: "+toReturn);
            return toReturn;
        }
    }
}
\end{lstlisting}
\subsubsection{Result.java}
\begin{lstlisting}
package backend.process;

import java.util.ArrayList;
import java.util.Date;
import java.util.HashSet;
import java.util.Set;

/**
 * Holds the data for one event in the Timeline.
 */
public class Result implements Comparable<Result>, Cloneable {
    private ArrayList<String> dates;
    private String event;
    private Set<String> subjects;
    private TimelineDate timelineDate;
    private FileData fileData;
    private String originalString;

    /**
     * Initialises variables.
     */
    public Result() {
        this("");
    }

    /**
     * Initialises variables and sets the event.
     *
     * @param event the event text of this event.
     */
    private Result(String event) {
        dates = new ArrayList<>();
        this.event = event;
        subjects = new HashSet<>();
        timelineDate = new TimelineDate();
        originalString = "";
    }

    /**
     * Checks whether this backend.process.Result already has a given date (in String format)
     *
     * @param date the String date being checked.
     * @return true if the list of String dates has date, false otherwise.
     */
    public boolean hasDate(String date) {
        for (String dDate : dates) {
            if (dDate.equals(date)) {
                return true;
            }
        }
        return false;
    }

    /**
     * The event string of this event.
     *
     * @return the string event of this event.
     */
    public String getEvent() {
        return event;
    }

    /**
     * Sets the event of this Timeline event.
     *
     * @param event what we are setting as an event.
     */
    public void setEvent(String event) {
        this.event = event;
    }

    /**
     * Add a string date to the list of held string dates, dates relevant to this event.
     *
     * @param date the String date that is being added.
     */
    public void addDate(String date) {
        dates.add(date);
    }

    /**
     * Get the list of String dates held.
     *
     * @return the list of String dates held.
     */
    public ArrayList<String> getDates() {
        return dates;
    }

    /**
     * Get the list of Subjects relevant to this event.
     *
     * @return the list of Subjects held.
     */
    public Set<String> getSubjects() {
        return subjects;
    }

    /**
     * Add a subject to the list of Subjects held.
     *
     * @param subject a subject to add.
     */
    public void addSubject(String subject) {
        subjects.add(subject);
    }

    /**
     * @return a String showing all the information held by this backend.process.Result object.
     */
    @Override
    public String toString() {
        return String.format("Date: %s, Subject: %s, Event: %s", timelineDate, subjects, event);
    }

    /**
     * For the given passed in String date, it will update the backend.process.TimelineDate of this backend.process.Result object.
     *
     * @param date a String that contains date information that needs to be passed into the timeline date of this backend.process.Result object.
     */
    public void addDate_1(String date, String baseDate) {
        System.out.println("About to parse: " + date);
        timelineDate.parse(date, baseDate);
    }


    /**
     * Used to compare two backend.process.Result objects by their backend.process.TimelineDate.
     *
     * @param o the other backend.process.Result object.
     * @return the result of comparing the backend.process.TimelineDate of this object, with the backend.process.TimelineDate of the other (input) backend.process.Result object.
     */
    @Override
    public int compareTo(Result o) {
        return this.timelineDate.compareTo(o.timelineDate);
    }

    /**
     * Get the TimelineDate for this Result.
     *
     * @return the date data used for this Result in the Timeline.
     */
    public TimelineDate getTimelineDate() {
        return timelineDate;
    }

    /**
     * Set the TimelineDate for this Result.
     *
     * @param timelineDate the date data used for this Result in the Timeline.
     */
    public void setTimelineDate(TimelineDate timelineDate) {
        this.timelineDate = timelineDate;
    }

    /**
     * Get the FileData for this Result.
     *
     * @return the File data of the File that produced this Result.
     */
    public FileData getFileData() {
        return fileData;
    }

    /**
     * Set the FileData for this Result.
     *
     * @param fileData the File data of the File that produced this Result.
     */
    public void setFileData(FileData fileData) {
        this.fileData = fileData;
    }

    /**
     * Used to produce a String of the Subjects of this Result. The String is every item in the Subjects set, separated
     * by commas.
     *
     * @return a String of the Subjects of this Result, separated by commas.
     */
    public String getSubjectsAsString() {
        String toReturn = "";
        for (int i = 0; i < subjects.size(); i++) {
            toReturn += subjects.toArray()[i];
            if (i < subjects.size() - 1) {
                toReturn += ", ";
            }
        }
        return toReturn;
    }

    /**
     * Makes a new Result object, with the same data as this one but not referencing to the same place in memory. So
     * changes in the new Result object do not change the data of this Result object.
     *
     * @return a new Result object with the same data but not pointing to the data in memory.
     * @throws CloneNotSupportedException
     */
    @Override
    public Object clone() throws CloneNotSupportedException {
        super.clone();
        Result copyResult = new Result();
        //copying the dates
        TimelineDate copyTimelineDate = new TimelineDate();
        copyTimelineDate.setDate1((Date) timelineDate.getDate1().clone());
        if (timelineDate.getDate2() != null) {
            copyTimelineDate.setDate2((Date) timelineDate.getDate2().clone());
        } else {
            copyTimelineDate.setDate2(null);
        }
        copyResult.setTimelineDate(copyTimelineDate);
        //copying the event
        copyResult.setEvent(event);
        //copying the subjects
        copyResult.subjects = new HashSet<>(subjects);
        //set the filedata
        copyResult.setFileData(fileData);//all results of the same file point to the same filedata (not a unique one)
        //set the original sentence
        copyResult.setOriginalString(originalString);
        return copyResult;
    }

    @Override
    public boolean equals(Object obj) {
        if(obj instanceof Result){
            Result other = (Result) obj;
            boolean areFileDataEqual = ((fileData == null && other.fileData == null) ||
                    (fileData != null && other.fileData != null && fileData.equals(other.fileData)));//as FileData is the only one that could be null

            return subjects.equals(other.getSubjects()) && event.equals(other.event) && originalString.equals(other.originalString)
                    && areFileDataEqual && timelineDate.equals(other.getTimelineDate());
        }
        return false;
    }

    /**
     * Get the original sentence that produced this Result object.
     *
     * @return the original sentence that produced this.
     */
    public String getOriginalString() {
        return originalString;
    }

    /**
     * Set the original sentence that produced this Result object.
     *
     * @param originalString the original sentence that produced this.
     */
    public void setOriginalString(String originalString) {
        this.originalString = originalString;
    }
}
\end{lstlisting}
\subsubsection{TimelineDate.java}
\begin{lstlisting}
package backend.process;

import edu.stanford.nlp.util.Pair;
import org.joda.time.DateTime;
import org.joda.time.Days;

import java.text.ParseException;
import java.text.SimpleDateFormat;
import java.util.*;
import java.util.regex.Pattern;

/**
 * Attempts to generate an exact date for an event, to then order the events.
 * Holds the start and end date (appropriately) for each event in the timeline. It updates as new dates, relevant to the
 */
//If just year-month should create range?
public class TimelineDate implements Comparable<TimelineDate> {
    private static final String year = "0001";
    private static final String month = "01";
    private static final String day = "01";
    private static final Map<String, Pair<String, String>> seasonMap;
    private static final Map<Character, String> durationMap;
    private static final Map<Character, String> timeMap;
    private final SimpleDateFormat simpleDateFormat = new SimpleDateFormat("yyyy-MM-dd G");
    private final SimpleDateFormat dayMonthYearFormat = new SimpleDateFormat("dd-MM-yyyy G");

    static {
        /*
            Seasons are given according to educationuk.org, by:
                Winter December , January, February (12,1,2)
                Spring is March, April and May (3,4,5)
                Summer is June, July, August (6,7,8)
                Autumn September, October, November (9,10,11)
         */
        Map<String, Pair<String, String>> seasonM = new HashMap<>();
        seasonM.put("WI", new Pair<>("12", "02"));
        seasonM.put("SP", new Pair<>("03", "05"));
        seasonM.put("SU", new Pair<>("06", "08"));
        seasonM.put("FA", new Pair<>("09", "11"));
        seasonMap = seasonM;

        //durationMap
        Map<Character, String> durationM = new HashMap<>();
        durationM.put('P', "Period:");
        durationM.put('Y', "Year(s)");
        durationM.put('M', "Month(s)");
        durationM.put('W', "Week(s)");
        durationM.put('D', "Day(s)");
        durationMap = durationM;

        //timeMap
        Map<Character, String> timeM = new HashMap<>();
        timeM.put('T', "Time:");
        timeM.put('H', "Hour(s)");
        timeM.put('M', "Minute(s)");
        timeM.put('S', "Second(s)");
        timeMap = timeM;
    }

    //patterns are thread safe
    private final static Pattern onlyYearPattern = Pattern.compile("(\\d{4})|(\\d{3}X)|(\\d{2}XX)|(\\dXXX)|(XXXX)");
    private final static Pattern onlyMonthPattern = Pattern.compile("\\d{2}");
    private final static Pattern onlyWeekNumberPattern = Pattern.compile("W\\d{2}");
    private final static Pattern onlySeasonPattern = Pattern.compile("[A-Z]{2}");
    private final static Pattern onlyDayPattern = Pattern.compile("\\d{2}.*");
    private final static Pattern onlyPresentRefPattern = Pattern.compile(".*PRESENT_REF.*");
    private final static Pattern onlyPastRefPattern = Pattern.compile(".*PAST_REF.*");
    private final static Pattern onlyFutureRefPattern = Pattern.compile(".*FUTURE_REF.*");
    private final static Pattern onlyBeforeYearPattern = Pattern.compile("(\\-\\d{4})|(\\-\\d{3}X)|(\\-\\d{2}XX)|(\\-\\dXXX)|(\\-XXXX)");
    private final static Pattern onlyWeekendPattern = Pattern.compile("WE");
    private final static Pattern yearMonthDayPattern = Pattern.compile("\\d{4}\\-\\d{2}\\-\\d{2}");
    private Calendar calendar;
    private Date date1;//first (min, start) date
    private Date date2;//second (max, end) date
    private String dateStr;
    private String baseDate;
    private String durationData;//holds the latest duration data (additional info to show with event)
    //have a pair of list dates and duration string, if you use the dates pass in the string as additional info
    private int range = -1;

    /**
     * Initialises the Calendar used to determine dates based on week number.
     */
    public TimelineDate() {
        calendar = Calendar.getInstance();
        simpleDateFormat.setLenient(false);//can only create correct dates
    }

    /**
     * Update the dates hold by this, based on the input text.
     *
     * @param date a date provided by the StanfordCoreNLP library: it is a normalized entity
     */
    public void parse(String date, String baseDate) {
        System.out.println("Input: " + date);
        Pair<ArrayList<Date>, String> dateDurationPair = null;
        ArrayList<Date> dates = new ArrayList<>();
        String durationData = null;
        this.baseDate = baseDate;
        //splitting INTERSECT
        String[] splitDate = date.split("INTERSECT");
        if (splitDate.length > 0) {// on the first part of the date, which is just a date, get its specfic date
            String possibleDates = splitDate[0];//this date could also be a range, ie include /
            String[] splitRange = possibleDates.split("/");
            for (String possibleDate : splitRange) {
                dates.addAll(getDate(possibleDate.trim()));// from processing the individual date, add it to dates
            }
            //process INTERSECT data
            if (splitDate.length > 1) {
                String trimmedDuration = splitDate[1].trim();
                System.out.println("Processing: " + trimmedDuration);
                //process trimmed part which contains the duration data
                durationData = processINTERSECT(trimmedDuration);
            }
            //resulting list of dates and duration data should be put in a pair that is processed, where durationData
            //is only set if one of our dates are set
            dateDurationPair = new Pair<>(dates, durationData);
        }
        //if we had INTERSECT then we should process it for additional info to show
        //if we have more than 2 dates in the list, then keep the minimum date and max date and remove all the others
        enforceRule(dateDurationPair, date);
    }

    /**
     * Processes the INTERSECT data (duration of an event) provided sometimes with Date normalized entity tags.
     * Based on the ISO Standard 8601
     *
     * @param intersectData the data after "INTERSECT" in the normalized Date entity tags. Should start with 'P'.
     * @return a technical String representation of that data based on the ISO Standard 8601, i.e. Period X Year(s) Y Day(s)...
     */
    private String processINTERSECT(String intersectData) {
        String toReturn = "";
        //input starts with P
        char[] dataSplit = intersectData.toCharArray();
        if (dataSplit.length > 0 && dataSplit[0] == 'P') {
            System.out.println("Entered Period");
            for (int i = 0; i < dataSplit.length; i++) {
                char info = dataSplit[i];
                System.out.println(info);
                //when find T (go to time process method, once returned from that break)
                if (info == 'T') {
                    //go time method, and what it returns add toReturn
                    //make string from here till end of data
                    String hereToEnd = intersectData.substring(i, intersectData.length());
                    System.out.println("End Data: " + hereToEnd);
                    String toAdd = processTime(hereToEnd);
                    System.out.println("Adding: " + toAdd);
                    toReturn += toAdd;
                    break;//not processing the rest of the array as its character data
                }
                String fullText = ((durationMap.get(info) != null) ? " " + durationMap.get(info) + " " : info + "");
                System.out.println("Found: " + fullText);

                toReturn += fullText;
            }
            System.out.println("Every: " + toReturn.trim());
        }
        return toReturn.trim();
    }

    /**
     * Processes the Time duration part of the INTERSECT data after a Data has been given a normalized entity tag.
     * Based on the ISO Standard 8601.
     *
     * @param timeData time data that starts with 'T', based on ISO Standard 8601.
     * @return a technical String representation of the data passed in, i.e. Time X Hour(s) Y Minute(s)...
     */
    private String processTime(String timeData) {
        String toReturn = "";
        char[] splitTimeData = timeData.toCharArray();
        if (splitTimeData.length > 0 && splitTimeData[0] == 'T') {
            System.out.println("Processing Time data");
            for (char charTime : splitTimeData) {
                System.out.println("Processing: " + charTime);
                String fullText = ((timeMap.get(charTime) != null) ? " " + timeMap.get(charTime) + " " : charTime + "");
                System.out.println("Found: " + fullText);
                toReturn += fullText;
            }
            System.out.println(toReturn.trim());
        }
        return toReturn.trim();
    }

    /**
     * For the given input, produce a list of dates based on it.
     * Based on the ISO Standard 8601.
     *
     * @param date an input text that contains date information (can be exact or relative).
     * @return a list of exact Dates formed from the input.
     */
    private ArrayList<Date> getDate(String date) {
        calendar.clear();
        ArrayList<Date> toReturn = new ArrayList<>();
        String year1 = year;
        String month1 = month;
        String day1 = day;
        String year2 = null;
        String month2 = null;
        String day2 = null;
        boolean isBC = false;
        System.out.println("Trying to match: " + date);
        //need to check if its a PRESENT_REF, FUTURE_REF or PAST_REF
        if (onlyPastRefPattern.matcher(date).matches()) {
            System.out.println("PAST REF");
            //past so make range from 0001-01-01 -> base date (range)
            if (yearMonthDayPattern.matcher(baseDate).matches()) {
                //base date has the format yyyy-MM-dd
                String[] splitBaseDate = baseDate.split("-");
                //year1, month1, day1 values stay with default value
                year2 = splitBaseDate[0];
                month2 = splitBaseDate[1];
                day2 = splitBaseDate[2];//safe as pattern matched
            }
        } else if (onlyPresentRefPattern.matcher(date).matches()) {
            System.out.println("PRESENT REF");
            //use the base date (single date)
            if (yearMonthDayPattern.matcher(baseDate).matches()) {
                System.out.println("BaseDate matches: " + baseDate);
                String[] splitDate = baseDate.split("-");//so its safe to split it into 3 parts as pattern matched above
                year1 = splitDate[0];
                month1 = splitDate[1];
                day1 = splitDate[2];
            }
        } else if (onlyFutureRefPattern.matcher(date).matches()) {
            System.out.println("FUTURE REF");
            //future, from now til the last date we allow 9999-12-31 (range)
            if (yearMonthDayPattern.matcher(baseDate).matches()) {
                System.out.println("BaseDate matches: " + baseDate);
                String[] splitDate = baseDate.split("-");//so its safe to split it into 3 parts as pattern matched above
                year1 = splitDate[0];
                month1 = splitDate[1];
                day1 = splitDate[2];
            }
            year2 = "9999";//future point
            month2 = "12";
            day2 = "31";
        } else {
            System.out.println("ELSE");
            //else, need to check whether it is BC or AD
            if (date.length() >= 5 && onlyBeforeYearPattern.matcher(date.substring(0, 5)).matches()) {//got a negative date
                System.out.println("Past Date");
                isBC = true;
                date = date.substring(1, date.length());//removed - sign infront of year
            }
            //then calculate date
            boolean isWeekNumber = false;
            String[] dateInfo = date.split("-");
            for (int i = 0; i < dateInfo.length; i++) {
                if (i == 0) {//this can only be a year
                    //check year format
                    if (onlyYearPattern.matcher(dateInfo[i]).matches()) {
                        if (isBC) {
                            year1 = dateInfo[i].replace("X", "9");
                        } else {
                            year1 = dateInfo[i].replace("X", "0");
                            //check if its all 0000, then its BC
                            if (year1.equals("0000")) {
                                year1 = "0001";//we will start from year 1 (AD)
                            }
                        }
                        if (dateInfo[i].contains("X")) {//if we do have a range then we need to set the values for the second date
                            if (isBC) {
                                year2 = dateInfo[i].replace("X", "0");
                            } else {
                                year2 = dateInfo[i].replace("X", "9");//could be the case that we dont want it to point to the last day in the year

                            }
                            month2 = "12";//last day of the second year
                            day2 = "31";//assuming a range for 1980s means 1980 to the last day of 1989 (maximum possible range)
                            //could have every day in january, would want it to end in january?
                        }
                    }
                } else if (i == 1) {//this can be a week number, a month number or a season
                    //checking if its a month
                    System.out.println("Checking: " + dateInfo[i]);
                    if (onlyMonthPattern.matcher(dateInfo[i]).matches()) {
                        System.out.println("In onlyMonthPattern");
                        month1 = dateInfo[i];
                    } else if (onlyWeekNumberPattern.matcher(dateInfo[i]).matches()) {//checking if its a week number
                        isWeekNumber = true;
                        //calculate month and start day-end
                        //split W from actual week number
                        if (isBC) {
                            calendar.set(Calendar.ERA, GregorianCalendar.BC);
                        } else {
                            calendar.set(Calendar.ERA, GregorianCalendar.AD);
                        }
                        String weekNumber = dateInfo[i].substring(1);//W is the first part of the string, after it is the week number
                        calendar.set(Calendar.YEAR, getInt(year1));//should be set from previously
                        calendar.set(Calendar.WEEK_OF_YEAR, getInt(weekNumber));
                        calendar.set(Calendar.DAY_OF_WEEK, Calendar.MONDAY);//assuming we start on monday and end on sunday
                        year1 = new SimpleDateFormat("yyyy").format(calendar.getTime());//in the case the week goes into the next year, update our year1
                        month1 = new SimpleDateFormat("MM").format(calendar.getTime());
                        day1 = new SimpleDateFormat("dd").format(calendar.getTime());
                        //now for the end of the week
                        calendar.set(Calendar.DAY_OF_WEEK, Calendar.SUNDAY);
                        month2 = new SimpleDateFormat("MM").format(calendar.getTime());
                        day2 = new SimpleDateFormat("dd").format(calendar.getTime());
                        year2 = new SimpleDateFormat("yyyy").format(calendar.getTime());
                    } else if (onlySeasonPattern.matcher(dateInfo[i]).matches()) {//checking if its a season
                        String season = dateInfo[i];
                        Pair<String, String> seasonPair = seasonMap.get(season);//get the start and end month for the season
                        if (seasonPair != null) {//year1 should be set previously
                            //not quite sure what to do with BC here
                            month1 = seasonPair.first;
                            month2 = seasonPair.second;
                            day2 = "31";//set the second day as we are using a range
                            if (seasonPair.first.equals("12")) {//move onto the next year, so year2 should be updated
                                year2 = (getInt(year1) + 1) + "";
                            } else {//we dont need to increment the year as the season is within the same year
                                year2 = year1;
                            }
                        }
                    }
                } else if (i == 2) {//can be a day, or previously had week this could be weekend
                    System.out.println("Checking: " + dateInfo[i] + "size: " + dateInfo[i].length());
                    if (isWeekNumber) {//then if it has another digit, its the day of the week (so its not a week, but a specific day of the week)
                        System.out.println("Day of the week: " + dateInfo[i]);
                        int day = getInt(dateInfo[i]);
                        //as the calendar here starts with sunday and in iso it starts with monday, we need to increase by one and mod
                        day = (day % 7) + 1;
                        calendar.set(Calendar.DAY_OF_WEEK, day);
                        year1 = new SimpleDateFormat("yyyy").format(calendar.getTime());//in the case the week goes into the next year, update our year1
                        month1 = new SimpleDateFormat("MM").format(calendar.getTime());
                        day1 = new SimpleDateFormat("dd").format(calendar.getTime());
                        day2 = null;
                        month2 = null;
                        year2 = null;
                    }
                    if (onlyDayPattern.matcher(dateInfo[i]).matches()) {//got the day
                        System.out.println("Got day: " + dateInfo[i]);
                        day1 = dateInfo[i].substring(0, 2);
                    } else if (onlyWeekendPattern.matcher(dateInfo[i]).matches()) {//previously should have had week number so its a range
                        //checking its range has been set before
                        if (day1 != null && day2 != null && month1 != null && month2 != null) {
                            //then lets refine the dates that were just before a week long, now down to a weekend
                            if (isBC) {
                                calendar.set(Calendar.ERA, GregorianCalendar.BC);
                            } else {
                                calendar.set(Calendar.ERA, GregorianCalendar.AD);
                            }
                            calendar.set(Calendar.DAY_OF_WEEK, Calendar.SATURDAY);
                            year1 = new SimpleDateFormat("yyyy").format(calendar.getTime());//in the case the week goes into the next year, update our year1
                            month1 = new SimpleDateFormat("MM").format(calendar.getTime());
                            day1 = new SimpleDateFormat("dd").format(calendar.getTime());
                            calendar.set(Calendar.DAY_OF_WEEK, Calendar.SUNDAY);
                            month2 = new SimpleDateFormat("MM").format(calendar.getTime());
                            day2 = new SimpleDateFormat("dd").format(calendar.getTime());
                            year2 = new SimpleDateFormat("yyyy").format(calendar.getTime());
                        }
                    }
                }
            }
        }
        //make the dates
        Date date1;
        date1 = createDates(year1, month1, day1, isBC);
        toReturn.add(date1);
        System.out.println(date1);
        if (year2 != null && month2 != null && day2 != null) {
            Date date2 = createDates(year2, month2, day2, isBC);
            toReturn.add(date2);
            System.out.println(date2);
        }
        return toReturn;
    }

    /**
     * Creates a Data with the given input data. In the worst case we overestimated the day value (i.e. 31 when for
     * that month it can be 30), so we check that it is a legal date. If it isn't, then we reduce the day value,until we
     * get a correct day value. This is assuming that only the date values are wrong not month and years(these are given
     * by normalized entity tags, which should be correct according to SUTime annotator)
     *
     * @param year  the year value for this Date
     * @param month the month value for this Date
     * @param day   the day value for this Date (which can be wrong, so we reduce it until its right, assuming we always overestimate.
     * @param isBC  whether or not this is a BC or AD date; true if it is a BC Date.
     * @return a correct Date object based on the data passed in.
     */
    private Date createDates(String year, String month, String day, boolean isBC) {
        Date toReturn;
        String date = returnDate(year, month, day, isBC);
        System.out.println("Trying to create date for: " + date);
        try {
            toReturn = simpleDateFormat.parse(date);
        } catch (ParseException e) {
            //e.printStackTrace();//could comment this out
            System.out.println("Couldnt create date, so trying for a lower value");
            //couldnt create date, so most likely day value is to high, so reduce it
            int dayUpdate = Integer.parseInt(day);
            int monthUpdate;
            dayUpdate--;
            //should update month and year if fall below a certain threshold?
            if (dayUpdate <= 0) {//should decrease month val and set dayUpdate to 31
                dayUpdate = 31;
                System.out.println("Our day value fell below threshold, so go to month before this");
                monthUpdate = Integer.parseInt(month);
                monthUpdate--;//assuming we will never produce a month  below 1
                //which makes sense, as we only ever over estimate day values
                toReturn = createDates(year, Integer.toString(monthUpdate), Integer.toString(dayUpdate), isBC);
            } else {
                //update day value, so try again with this value
                toReturn = createDates(year, month, Integer.toString(dayUpdate), isBC);
            }
        }
        System.out.println("Created date for: " + toReturn);
        return toReturn;
    }

    /**
     * Used to find update the min/max dates held. Will update the dates held if a new min/max has been found.
     * MinMax Algorithm.
     *
     * @param dateDurationPair a Pair that has a list of possible new min/max dates, and their corresponding duration data (which can be null).
     * @param date             the string that produced these dates.
     */
    private void enforceRule(Pair<ArrayList<Date>, String> dateDurationPair, String date) {
        if (dateDurationPair != null) {
            ArrayList<Date> newDates = dateDurationPair.first();
            for (Date newDate : newDates) {
                if (this.date1 == null || this.date1.compareTo(newDate) > 0) {//if we dont have a date1, or we have a smaller one
                    this.date1 = newDate;
                    dateStr = date;
                    //else, we found a newDate that we havent set that is bigger than date1, look at date2
                    durationData = dateDurationPair.second();
                } else if ((this.date2 == null || this.date2.compareTo(newDate) < 0) && !this.date1.equals(newDate)) {//if we dont have date2, or we found a bigger date
                    this.date2 = newDate;                       //and the new date is not the first date
                    dateStr = date;
                    durationData = dateDurationPair.second();
                }
            }
            updateRange();
        }
    }

    /**
     * Updates the int value of the range based on the date1 and date2 values. If date2 is null then the range is 0, ie
     * the size of the range of dates is 0.
     */
    private void updateRange() {
        if (date1 != null && date2 != null) {
            DateTime dateTime1 = new DateTime(date1);
            DateTime dateTime2 = new DateTime(date2);
            range = Days.daysBetween(dateTime1, dateTime2).getDays();
        } else {
            range = 0;
        }
    }

    /**
     * Get the Range, ie the number of days between date1 and date2.
     *
     * @return the number of days between date1 and date2 (0 if date2 == null)
     */
    public int getRange() {
        updateRange();
        return range;
    }

    /**
     * Produces a date String of the format yyyy-MM-dd, for the given input.
     *
     * @param year  the year of the date
     * @param month the month of the date
     * @param day   the day of the date
     * @return a String of the format yyyy-MM-dd
     */
    private String returnDate(String year, String month, String day, boolean isBC) {
        if (isBC) {
            return String.format("%s-%s-%s BC", year, month, day);
        }
        return String.format("%s-%s-%s AD", year, month, day);
    }

    /**
     * On the given input, produce an integer.
     *
     * @param number a number in String form.
     * @return the integer value corresponding to the string, or 1 if it is not possible to produce an int.
     */
    private int getInt(String number) {
        try {
            return Integer.parseInt(number);
        } catch (Exception e) {
            return 1;
        }
    }

    /**
     * @return date1 -> date2, or just date1 if we don't have a date2
     */
    @Override
    public String toString() {
        String toReturn = "";
        SimpleDateFormat simpleDateFormat = new SimpleDateFormat("yyyy-MM-dd G");
        if (date1 != null) {
            toReturn += simpleDateFormat.format(date1);
        }
        if (date2 != null) {
            toReturn += " -> " + simpleDateFormat.format(date2);
        }
        if (durationData != null) {
            toReturn += String.format(" (%s)", durationData);
        }
        return toReturn;
    }

    /**
     * Compares two TimelineDates based on their start date.
     *
     * @param o the other backend.process.TimelineDate that is being compared to.
     * @return the comparison of this date1 to the other's date1, or -1 if the other result does not have a date1, or 1 if this does not have a date1.
     */
    @Override
    public int compareTo(TimelineDate o) {
        if (range == -1) {
            updateRange();
        }
/*        if (o.date1 != null && this.date1 != null) {
            return o.date1.compareTo(this.date1);
        }
        if (this.date1 == null) {
            return -1;
        }
        return 1;*/
        if (range > o.range) {
            return 1;
        } else if (range < o.range) {
            return -1;
        } else {
            return 0;
        }
    }

    /**
     * Checks whether a given backend.process.TimelineDate object is equal to this object.
     *
     * @param obj the other object we are comparing to.
     * @return true if the other backend.process.TimelineDate object has the same dates as this.
     */
    @Override
    public boolean equals(Object obj) {
        boolean toReturn;
        try {
            TimelineDate other = (TimelineDate) obj;
            toReturn = date1.equals(other.date1);
            if (toReturn) {
                if (date2 != null || other.date2 != null) {//if we have a scond date then so should other, if we dont then they shouldnt either
                    toReturn = date2.equals(other.date2);
                }
            }
        } catch (Exception e) {
            return false;
        }
        return toReturn;
    }

    /**
     * Set the Date for date1.
     *
     * @param date1 Date for date1.
     */
    public void setDate1(Date date1) {
        this.date1 = date1;
        updateRange();
    }

    /**
     * Set the Date for date2.
     *
     * @param date2 Date for date2.
     */
    public void setDate2(Date date2) {
        this.date2 = date2;
        updateRange();
    }

    /**
     * Get the Date for date1.
     *
     * @return Date for date1.
     */
    public Date getDate1() {
        return date1;
    }

    /**
     * Get the Date for date2.
     *
     * @return Date for date2.
     */
    public Date getDate2() {
        return date2;
    }

    /**
     * Get the additional duration data for a date.
     *
     * @return String representing the duration an event occurred for (every when it repeated).
     */
    public String getDurationData() {
        return durationData;
    }

    /**
     * Get date1 as a String in the format dd-MM-yyyy G, or null if that is not possible.
     *
     * @return a formatted String of date1.
     */
    public String getDate1FormattedDayMonthYear() {
        try {
            return (date1 != null) ? dayMonthYearFormat.format(date1) : null;
        } catch (Exception e) {
            return null;
        }
    }

    /**
     * Get date2 as a String in the format dd-MM-yyyy G, or null if that is not possible.
     *
     * @return a formatted String of date2.
     */
    public String getDate2FormattedDayMonthYear() {
        try {
            return (date2 != null) ? dayMonthYearFormat.format(date2) : null;
        } catch (Exception e) {
            return null;
        }
    }
}
\end{lstlisting}

\subsection{Ranges}
\subsubsection{ProduceRanges.java}
\begin{lstlisting}
package backend.ranges;

import backend.process.Result;

import java.util.*;

/**
 * Produces a List of Range (that contain ranges - like a Tree, so it Returns a Forest) to use in showing the Timeline.
 */
public class ProduceRanges {
    private List<Range> trees = new ArrayList<>();

    /**
     * Get the list of Range Trees formed by processing the list of the Results.
     *
     * @return a list of Range Trees.
     */
    public List<Range> getTrees() {
        return trees;
    }

    /**
     * For the given Result list, sort them by size of range of dates held by their TimelineDate, then make a Forest of
     * Range Trees using the list of Results, and then finally sort the Forest by their date1 value (descending)
     *
     * @param resultList the given Result list.
     */
    public void produceRanges(List<Result> resultList) {
        System.out.println("Number of Results: " + resultList.size());
        sortByRange(resultList);//sort the list by their Range
        makeForest(resultList);//from the sorted list attempt to add it to existing Range Trees or create a new Tree if not possible.
        sortForest();//sort the Range (recursively) Forest by their date1 value
    }

    /**
     * For the given list of Results, sort them by the number of days in between their first range of dates and second.
     *
     * @param inputResults the given list of Results.
     * @return the list of Results sorted by their Range (descending)
     */
    private List<Result> sortByRange(List<Result> inputResults) {
        Collections.sort(inputResults);//java +7 mergesort with O(nlogn) but if its almost sorted its closer to O(n)
        Collections.reverse(inputResults);//as the list has been sorted in ascending order and we want descending
        System.out.println("Sorted by Range");
        for(Result result: inputResults){
            System.out.println(result+" "+result.getTimelineDate().getRange());
        }
        return inputResults;
    }

    /**
     * For the given sorted list of Results (has to be sorted for this algorithm to work properly as we are making the
     * Range trees starting from the highest Range) attempt to add each Result to existing Range trees, if that is not
     * possible then create a new Range tree for the Result.
     *
     * @param sortedResults the list of sorted Result objects.
     */
    private void makeForest(List<Result> sortedResults) {
        //if have to update parent, make copy of current, then reset, then set copy as child
        for (Result result : sortedResults) {//should be sorted from largest range to lowest
            //try to add to one of the trees if not possible make new range
            if (!addToRange(result)) {
                //we couldnt add it so we make a new Range
                Range range = new Range(result.getTimelineDate().getDate1(), result.getTimelineDate().getDate2());
                range.add(result);
                trees.add(range);
            }
        }
    }

    /**
     * Attempt to add the given Result to the the different Range trees.
     *
     * @param result the given Result.
     * @return true if the Result was added to one of the existing Range trees; false otherwise.
     */
    private boolean addToRange(Result result) {
        for (Range range : trees) {
            if (range.add(result)) {
                return true;
            }
        }
        return false;
    }

    /**
     * For the forest of Range trees recursively sort the trees (and their subtrees) by their date1 in ascending order.
     */
    private void sortForest() {
        //sort the roots first
        Collections.sort(trees);
        //for all the roots call sortChildren (which will recursively sort)
        for (Range root : trees) {
            root.sortChildren();
        }
    }


}

\end{lstlisting}
\subsubsection{Range.java}
\begin{lstlisting}
package backend.ranges;

import backend.process.Result;
import backend.process.TimelineDate;

import java.text.SimpleDateFormat;
import java.util.*;

/**
 * A Range of Dates (which can be only one date) which can be represented as a Node in a Tree, due to it holding a list of
 * its sub-ranges.
 */
public class Range implements Cloneable, Comparable<Range> {
    private List<Result> results;//results that are exactly in this range
    private Date date1;//first part of range
    private Date date2;//only if we have a range of size > 0
    private List<Range> children;//the children in this range within constraint children.date1 >= date1 children.date1 < date2 && children.date2 <= date2

    public Range(Date date1, Date date2) {
        this.date1 = date1;
        this.date2 = date2;
        results = new ArrayList<>();
        children = new ArrayList<>();
    }

    /**
     * For the given Result, check that the dates of the Result are partially in the scope of this Range, if that is not
     * possible then return false.
     * If the Result is within this Rage, then attempt to find between the current Range and its children a Range where
     * the Results range of dates fit in perfectly, if it can be added then add it else attempt to create a new Range
     * child of this Range and find its location to place it.
     * <p>
     * Algorithm given by:
     * 1.check constraints, if they are false then return false else move onto step 2.
     * 2.attempt to add the Result to one of the children (or children-children, or etc) of this Range where they both
     * have the same Range; move onto step 3 if that is unsuccessful.
     * 3.create a new Range and find the location of where to add it in the Range (extending the Range where necessary)
     *
     * @param result the given Result.
     * @return true if the Result was added to this Range (i.e. this Tree) or its children or its children-children etc;
     * false otherwise.
     */
    public boolean add(Result result) {
        TimelineDate timelineDate = result.getTimelineDate();
        //check constraints
        if (!shouldAdd(timelineDate)) {//check constraints if we can even add to this range (ie are we in the root constraints)
            return false;
        }
        //attempt to add through an existing range
        Range toAdd = checkCanAdd(result);
        if (toAdd != null) {
            toAdd.results.add(result);//add to the results of the given range where we could add
            return true;
        }
        //now we try to extend the range
        return createRangeAndAdd(result);
    }

    /**
     * Whether or not the given TimelineDate is partially or fully in the range of Dates of this Range. Whether or not
     * we should try to add the Result of this TimelineDate to this Range.
     *
     * @param timelineDate the given TimelineDate.
     * @return true if the TimelineDate is within/partially (ie if date1 is in the Range but not date2, or vice versa,
     * or both dates are within the range)in this Range; false otherwise.
     */
    private boolean shouldAdd(TimelineDate timelineDate) {
        //check the date1 is in constraint to t.date1 then if date2 then check with date2
        return isWithinConstraints(date1, date2, timelineDate.getDate1())
                || (timelineDate.getDate2() != null && isWithinConstraints(date1, date2, timelineDate.getDate2()));
    }

    /**
     * Attempt to find a Range (either this or one of its children, or its children-children, etc) that has the date
     * values that match the given Result.
     *
     * @param result the given Result.
     * @return a Range that holds the same exact date values as the TimelineDate of the Result or null if that is not
     * possible.
     */
    private Range checkCanAdd(Result result) {
        TimelineDate timelineDate = result.getTimelineDate();
        if (timelineDate.getDate1().equals(date1) && ((date2 == null && timelineDate.getDate2() == null)
                || (date2 != null && timelineDate.getDate2() != null && date2.equals(timelineDate.getDate2())))) {
            return this;
        }
        Range toReturn = null;
        for (Range child : children) {
            toReturn = child.checkCanAdd(result);
            if (toReturn != null) {
                return toReturn;
            }
        }
        return toReturn;
    }

    /**
     * For the given Result, create a Range using its TimelineDate and attempt to find the location where this Range
     * should be added in the Tree structure (ie find the lowest/best possible Range that should contain the created
     * Range).
     * Normally called if we couldn't find any Range that has the same exact Dates so we create a Range to contain it
     * and attempt to find its spot (extending Ranges where necessary to self contain Ranges in bigger Ranges).
     *
     * @param result the given Result.
     * @return true if the created Range was added to this (sub-)tree, false otherwise.
     */
    private boolean createRangeAndAdd(Result result) {
        //we couldn't add it to any of the childrens so make new child and try to add it
        TimelineDate timelineDate = result.getTimelineDate();
        Range newChild = new Range(timelineDate.getDate1(), timelineDate.getDate2());
        newChild.add(result);//the new potential range child holds its result
        return addChild(newChild);// try and add this child (need to check constraints
    }


    /**
     * For the given Range (r') attempt to add it to this Range (r). For this to occur it must be within the range of dates of
     * this Range (i.e. r'.date1 >= r.date1 && r'.date1 <r.date2 && r'.date2 < r.date2) or it must overlap (i.e. date 1
     * is within the range, but date2 isn't or vice versa). In the latter case we extend this Range, such that it holds
     * the old range and the given range.
     *
     * @param potentialChild the given Range (r').
     * @return true if the potentialChild could be added to the children of this Range, or we could extend the Range;
     * false otherwise.
     */
    public boolean addChild(Range potentialChild) {
        /*
            if you are within the constraints
                then attempt to add it to the children,if it cant because there are no
                    children, or it doesnt fit the childrens constraints then add it to this Range's children
            check for partial constraints
                if it matches then extend the ranges and add it to the new extended range (that holds as a child the old Range of this)
            otherwise return false (it doesn't belong in this area)
         */
        if (isWithinConstraints(date1, date2, potentialChild.date1)
                && (potentialChild.date2 == null || isWithinConstraints(date1, date2, potentialChild.date2))) {
            //we are within the constraints, try to add to children Ranges, if cant because there are none, or we cant then add directly to this Ranges children
            for (Range child : children) {
                if (child.addChild(potentialChild)) {//successfully added to one of them
                    return true;
                }
            }
            //we couldnt add to any of them, then add directly to my children
            children.add(potentialChild);
            return true;
            //partial constraints
        } else if (isWithinConstraints(date1, date2, potentialChild.date1) && (potentialChild.date2 != null && !isWithinConstraints(date1, date2, potentialChild.date2))) {
            //checking if we need to extend the range (this case the first date is within the constraints but date2 isnt
            //date1 is in the range but date2 isnt so extend and us to the extensions children as we didnt fit in previous range (which is now its child)
            extendRange(date1, potentialChild.date2);//extend the range and add the child to the now extended range (as we fit in the now extended
            children.add(potentialChild);//range, but not in its other child which was the previous range held before extending, else we wouldnt be here)
            return true;//if we called addChild it would just loop on the child that was the previous range, and thats the only child so there is no point checking just add to this
        } else if (!isWithinConstraints(date1, date2, potentialChild.date1) && (potentialChild.date2 != null && isWithinConstraints(date1, date2, potentialChild.date2))) {
            //date2 is in the range but date1 isnt so extend and add to the children of the extended range.
            extendRange(potentialChild.date1, date2);
            children.add(potentialChild);
            return true;
        }
        //We werent able to add the potentialRange to this Range,
        return false;//so return false (this method is recursive, so it will try on the next child of the parent Range that called it on this child)
    }

    /**
     * Used to extend the current Range to the given Dates, and hold the previous data in a new Range as a child.
     *
     * @param date1 the start of the extended Range.
     * @param date2 the end of the extended Range.
     */
    private void extendRange(Date date1, Date date2) {
        try {
            Range cloneRange = (Range) clone();//the date of this will be a new child, we are setting its parent
            clear();//clear all the values and set the new ones
            this.date1 = date1;//extended range value
            this.date2 = date2;//extended range value
            children.add(cloneRange);//add what used to be this (but now we cloned it) as out child
        } catch (CloneNotSupportedException e) {
            e.printStackTrace();
        }
    }

    /**
     * For the given Date, return a String of the format: dd-MM-yyyy G (days-months-years ERA)
     *
     * @param date the given Date.
     * @return a String of the format dd-MM-yyyy G or null if the given date can't be formatted (i.e. is null)
     */
    private String printDate(Date date) {
        SimpleDateFormat simpleDateFormat = new SimpleDateFormat("dd-MM-yyyy G");
        try {
            return simpleDateFormat.format(date);
        } catch (Exception e) {
        }
        return null;
    }

    /**
     * Clear all the data held by the current Range.
     */
    private void clear() {
        date1 = null;
        date2 = null;
        children = new ArrayList<>();
        results = new ArrayList<>();
    }

    /**
     * Check that the given Date toCheck is within the other two given Dates constraint1 and constraint2 (both inclusive)
     *
     * @param constraint1 the given start constraint Date (inclusive).
     * @param constraint2 the given end constraint Date (inclusive).
     * @param toCheck     the given Date we are checking.
     * @return true if the given Date to check is within the two constraints, or if any of them are null equal to the other;
     * false otherwise.
     */
    private boolean isWithinConstraints(Date constraint1, Date constraint2, Date toCheck) {
        if (constraint1 != null && constraint2 != null) {
            return toCheck.compareTo(constraint1) >= 0 && toCheck.compareTo(constraint2) <= 0;
        }
        if (constraint1 == null && constraint2 != null) {
            return toCheck.equals(constraint2);
        }
        if (constraint1 != null) {
            return toCheck.equals(constraint1);
        }
        return false;
    }

    /**
     * Return a clone of this Range. Holds the same data.
     *
     * @return an Object that can be cast to a Range object, which holds the current Data of the Range that called it
     * (when clone() was called).
     * @throws CloneNotSupportedException
     */
    @Override
    protected Object clone() throws CloneNotSupportedException {
        super.clone();
        Range cloneRange = new Range(date1, date2);
        cloneRange.children = new ArrayList<>(children);
        cloneRange.results = new ArrayList<>(results);
        return cloneRange;
    }

    /**
     * Returns a String of the given Range.
     *
     * @return a String that shows the range of Dates held by this Range, its Result and the its child Ranges (recursive).
     */
    @Override
    public String toString() {
        String toReturn = printDate(date1);
        if (date2 != null) {
            toReturn += " -> " + printDate(date2);
        }
        toReturn += " has results: " + results;
        toReturn += " and has children: \n" + children;
        return toReturn;
    }

    /**
     * Overwrote the equals method to check whether two Ranges are equal to each other: that is they hold the same Dates,
     * results and their children are equal (recursive).
     *
     * @param obj the other Range we are checking if its equal to this Range.
     * @return true if the other Range has the same Dates as this Range, holds the same Results as this Range, and their
     * children are equal.
     */
    @Override
    public boolean equals(Object obj) {
        try {
            Range other = (Range) obj;
            if (date1.equals(other.date1) && ((date2 == null & other.date2 == null) || (date2 != null & other.date2 != null && date2.equals(other.date2)))) {
                if (children.isEmpty() && other.children.isEmpty()) {
                    return results.equals(other.results);
                } else {//they have children
                    return results.equals(other.results) && children.equals(other.children);
                }
            }
        } catch (Exception e) {
            return false;
        }
        return false;
    }


    /**
     * Sort the children held by this Range, and calls sortChildren on its children.
     * Sorts them such that the resulting children list is in descending order based on date1 (ie at index 0 you have
     * the newest range and at the last index the range with date1 furthest away.
     */
    public void sortChildren() {
        if (children.size() > 0) {
            Collections.sort(children);
            for (Range range : children) {
                range.sortChildren();
            }
        }
    }


    /**
     * Implemented the Comparable interface to sort Ranges.
     *
     * @param o the other Range we are comparing to.
     * @return 1 if the date1 of this Range is greater than the other, 0 if they are equal and -1 if this Range's date1
     * is less than the other.
     */
    @Override
    public int compareTo(Range o) {
        if (o.date1 != null && this.date1 != null) {
            return date1.compareTo(o.date1);
        }
        if (this.date1 == null) {
            return -1;
        }
        return 1;
    }

    /**
     * Get the list of Range children held by this Range.
     *
     * @return the list of Range children.
     */
    public List<Range> getChildren() {
        return children;
    }

    /**
     * Get a String representation of date1 and date2 (if set) of this Range (i.e. a String of the range of dates held
     * by this Range).
     *
     * @return a String representation of the range of dates held.
     */
    public String getDateRange() {
        String toReturn = printDate(date1);
        if (date2 != null) {
            toReturn += " -> " + printDate(date2);
        }
        return toReturn;
    }

    /**
     * A list of the Results held by this Range. These are Results that have the same exact TimelineDate range of dates
     * as this Range (ie date1 and date2 of the Results and this Range are equal).
     *
     * @return the list of Results held by this Range.
     */
    public List<Result> getResults() {
        return results;
    }
}
\end{lstlisting}
\subsection{System}
\subsubsection{BackEndSystem.java}
\begin{lstlisting}
package backend.system;

import edu.stanford.nlp.pipeline.StanfordCoreNLP;
import edu.stanford.nlp.util.PropertiesUtils;

/**
 * Holds all the data needed by the entire Back-end: the StanfordCoreNLP used to process text, the System state, etc.
 * Follows a Singleton design pattern, as there should be just one backend.system.BackEndSystem during the entire lifetime of the
 * system, and allows to hold data needed in different areas to be held at one place accessible to all areas.
 * Before the Object is created the backend.system.SystemState is NOT_STARTED, it moves to STARTED when it is created.
 */
public class BackEndSystem {
    private static BackEndSystem ourInstance = new BackEndSystem();
    private StanfordCoreNLP coreNLP;
    private SystemState systemState = SystemState.NOT_STARTED;
    private Settings settings;

    /**
     * Will make a new backend.system.BackEndSystem if one has not been created yet, or else return the one that was previously
     * created.
     *
     * @return a backend.system.BackEndSystem that was just created or was previously created.
     */
    public static BackEndSystem getInstance() {
        return ourInstance;
    }

    /**
     * Initialises the StanfordCoreNLP and sets the backend.system.SystemState to STARTED.
     */
    private BackEndSystem() {
        coreNLP = new StanfordCoreNLP(PropertiesUtils.asProperties(
                "annotators",
                "tokenize,ssplit,pos,lemma,ner,entitymentions,parse,dcoref",
                "tokenize.language", "en"
        ));
        System.out.println("Finished running");
        systemState = SystemState.STARTED;
        settings = new Settings(true);//load the settings of the file, or use default Settings
    }

    /**
     * Get a reference to the StandfordCoreNLP that has already been loaded with all the models.
     *
     * @return a StanfordCoreNLP with all the models loaded.
     */
    public StanfordCoreNLP getCoreNLP() {
        return coreNLP;
    }

    /**
     * Get the current state of the System according to the backend.system.SystemState enum.
     *
     * @return the current state of the System.
     */
    public SystemState getSystemState() {
        return systemState;
    }

    /**
     * Update the current state of the System, i.e. when the we are beginning to process Files the backend.system.SystemState should be
     * set to PROCESSING.
     *
     * @param systemState the state to which the System should be set to.
     */
    public void setSystemState(SystemState systemState) {
        this.systemState = systemState;
        System.out.println("System is now in state: " + systemState);
    }

    /**
     * Get the Settings that this System uses (i.e preferred width/height on startup, how many threads to run in parallel, etc).
     *
     * @return the Settings used by this System.
     */
    public Settings getSettings() {
        return settings;
    }

    /**
     * Set the Settings that this System uses.
     *
     * @param settings the Settings used by this System.
     */
    public void setSettings(Settings settings) {
        this.settings = settings;
        this.settings.saveSettingsFile();
    }
}
\end{lstlisting}
\subsubsection{Settings.java}
\begin{lstlisting}
package backend.system;


import com.google.gson.*;
import com.google.gson.annotations.Expose;

import java.io.*;
import java.net.URISyntaxException;
import java.util.Scanner;

/**
 * Class that represents the possible Settings of the System.
 * Settings are attempted to be retrieved and/or stored in a Settings file, where the data is encoded in JSON.
 */
public class Settings implements Cloneable {
    //default values (to reset to default, or if settings file not found)
    public final static int defaultMaxNoThreads = 2;
    public final static int defaultThresholdSummary = 10;
    public final static int defaultWidth = 1024;
    public final static int defaultHeight = 800;
    private final static String threadTag = "maxNoOfThreads";
    private final static String thresholdTag = "thresholdSummary";
    private final static String widthTag = "width";
    private final static String heightTag = "height";
    //actual values
    @Expose(deserialize = false)
    private int maxNoOfThreads;
    @Expose(deserialize = false)
    private int thresholdSummary;
    @Expose(deserialize = false)
    private int width;
    @Expose(deserialize = false)
    private int height;

    /**
     * Used to create an instance of the Settings of this System, by attempting to load the Settings file.
     * If this can't be done the default values of the Settings are used to populate the fields.
     *
     * @param loadFromFile whether or not the Settings File should be loaded.
     */
    public Settings(boolean loadFromFile) {
        //try to read from settings file, if not possible, then use default values
        if (!loadFromFile || !loadSettingsFile()) {
            reset();
        }
    }

    /**
     * Used to create an instance of the Settings of this System, where the fields are populated by the default values.
     */
    public Settings() {
        this(false);
    }

    /**
     * Attempt to load the Settings from the Settings file.
     *
     * @return whether or not the Settings file could be retrieved and the fields of the Settings be populated.
     */
    private boolean loadSettingsFile() {
        System.out.println("Trying to load from file");
        //load the file, return true if it could
        try {
            File settingsFile = new File("settings.ini");
            //now need to read settings file
            System.out.println("Could find settings.ini file");
            String text = getJSONString(settingsFile);
            System.out.println("String in file: " + text);
            JsonElement jsonElement = new JsonParser().parse(text);
            JsonObject jsonObject = jsonElement.getAsJsonObject();
            System.out.println(jsonObject);
            setValues(jsonObject);
        } catch (Exception e) {
            return false;
        }
        return true;
    }

    /**
     * For the given File attempt to retrieve the JSON stored in it.
     *
     * @param file the given File (which should have the JSON settings data).
     * @return a String of the JSON encoding of the Settings in the File.
     * @throws FileNotFoundException thrown if the File could not be found.
     */
    private String getJSONString(File file) throws FileNotFoundException {
        Scanner scanner = new Scanner(file);
        StringBuilder stringBuilder = new StringBuilder();
        while (scanner.hasNext()) {
            stringBuilder.append(scanner.nextLine());
        }
        scanner.close();
        return stringBuilder.toString();
    }

    /**
     * For the given JsonObject, populate the Settings field with its data, if its present, else use the default values.
     *
     * @param jsonObject the given JsonObject that should have the Settings data.
     */
    private void setValues(JsonObject jsonObject) {
        maxNoOfThreads = (jsonObject.get(threadTag) != null) ? jsonObject.get(threadTag).getAsInt() : defaultMaxNoThreads;
        thresholdSummary = (jsonObject.get(thresholdTag) != null) ? jsonObject.get(thresholdTag).getAsInt() : defaultThresholdSummary;
        width = (jsonObject.get(widthTag) != null) ? jsonObject.get(widthTag).getAsInt() : defaultWidth;
        height = (jsonObject.get(heightTag) != null) ? jsonObject.get(heightTag).getAsInt() : defaultHeight;
        if (!isConstrained()) {//if constraints have not been set then reset
            reset();
        }
    }


    /**
     * For the Settings fields in this case, store them in JSON format in the Settings File.
     *
     * @return whether or not the Settings of this File could be saved in a File.
     */
    public boolean saveSettingsFile() {
        final GsonBuilder gsonBuilder = new GsonBuilder();
        gsonBuilder.excludeFieldsWithoutExposeAnnotation();
        final Gson gson = gsonBuilder.create();
        try {
            String json = gson.toJson(this);
            System.out.println(json);
            saveToSettingsFile(json);
        } catch (Exception e) {
            e.printStackTrace();
            return false;
        }
        return true;
    }

    /**
     * Save the given text in the "settings.ini" File.
     *
     * @param text the given Text
     * @throws IOException if an error occurs while accessing the File
     */
    private void saveToSettingsFile(String text) throws IOException {
        PrintWriter writer = new PrintWriter("settings.ini");
        writer.println(text);
        writer.close();
    }

    /**
     * For all the Fields in the Settings class, set their values to the default values.
     */
    public void reset() {
        //resets all the values to default
        maxNoOfThreads = defaultMaxNoThreads;
        thresholdSummary = defaultThresholdSummary;
        width = defaultWidth;
        height = defaultHeight;
    }

    /**
     * Getter for the Max Number of Threads that are allowed to run in parallel in this System.
     *
     * @return the Max Number of Threads that are allowed to run in parallel in this System.
     */
    public int getMaxNoOfThreads() {
        return maxNoOfThreads;
    }

    /**
     * Setter for the Max Number of Threads that are allowed to run in parallel in this System.
     *
     * @param maxNoOfThreads the Max Number of Threads that are allowed to run in parallel in this System.
     */
    public void setMaxNoOfThreads(int maxNoOfThreads) {
        if (isMaxNoOfThreadsConstrained(maxNoOfThreads)) {
            this.maxNoOfThreads = maxNoOfThreads;
        }
    }

    /**
     * Getter for the Threshold used when attempting to produce a Summary of a Sentence.
     *
     * @return the Threshold used when attempting to produce a Summary of a Sentence.
     */
    public int getThresholdSummary() {
        return thresholdSummary;
    }

    /**
     * Setter for the Threshold used when attempting to produce a Summary of a Sentence.
     *
     * @param thresholdSummary the Threshold used when attempting to produce a Summary of a Sentence.
     */
    public void setThresholdSummary(int thresholdSummary) {
        if (isThresholdConstrained(thresholdSummary)) {
            this.thresholdSummary = thresholdSummary;
        }
    }

    /**
     * Getter for the preferred start up width of the Program window.
     *
     * @return the preferred start up width of the Program window.
     */
    public int getWidth() {
        return width;
    }

    /**
     * Setter for the preferred start up width of the Program window.
     *
     * @param width the preferred start up width of the Program window.
     */
    public void setWidth(int width) {
        if (isWidthConstrained(width)) {
            this.width = width;
        }
    }

    /**
     * Getter for the preferred start up height of the Program window.
     *
     * @return the preferred start up height of the Program window.
     */
    public int getHeight() {
        return height;
    }

    /**
     * Setter for the preferred start up height of the Program window.
     *
     * @param height the preferred start up height of the Program window.
     */
    public void setHeight(int height) {
        if (isHeightConstrained(height)) {
            this.height = height;
        }
    }

    /**
     * Produce a clone of this Object, such that the populated data in the clone is a clone of the data in this Object.
     *
     * @return a Clone of this Object (data values are the same but stored in a different area in memory)
     * @throws CloneNotSupportedException
     */
    @Override
    public Object clone() throws CloneNotSupportedException {
        super.clone();
        Settings clonedSettings = new Settings();
        clonedSettings.setHeight(height);
        clonedSettings.setMaxNoOfThreads(maxNoOfThreads);
        clonedSettings.setThresholdSummary(thresholdSummary);
        clonedSettings.setWidth(width);
        return clonedSettings;
    }

    /**
     * Whether or not the given threshold summary is within the constraints.
     *
     * @param thresholdSummary the given threshold summary
     * @return true if the threshold summary is within the constraints; false otherwise.
     */
    private boolean isThresholdConstrained(int thresholdSummary) {
        return thresholdSummary > 0 && thresholdSummary <= 20;
    }

    /**
     * Whether or not the given max number of threads is within the constraints.
     *
     * @param maxNoOfThreads the given max number of threads.
     * @return true if the max number of threads is within the constraints; false otherwise.
     */
    private boolean isMaxNoOfThreadsConstrained(int maxNoOfThreads) {
        return maxNoOfThreads > 0 && maxNoOfThreads <= 20;
    }

    /**
     * Whether or not the given width is within the constraints.
     *
     * @param width the given width.
     * @return true if the width is within the constraints; false otherwise.
     */
    private boolean isWidthConstrained(int width) {
        return width > 0 && width <= 1920;
    }

    /**
     * Whether or not the given height is within the constraints.
     *
     * @param height the given height.
     * @return true if the height is within the constraints; false otherwise.
     */
    private boolean isHeightConstrained(int height) {
        return height > 0 && height <= 1080;
    }

    /**
     * Whether or not all the fields in this Settings object are all within their constraints.
     *
     * @return true if all the fields are within their constraints; false otherwise.
     */
    private boolean isConstrained() {
        return isThresholdConstrained(thresholdSummary) && isMaxNoOfThreadsConstrained(maxNoOfThreads)
                && isWidthConstrained(width) && isHeightConstrained(height);
    }

}
\end{lstlisting}
\subsubsection{SystemState.java}
\begin{lstlisting}
package backend.system;

/**
 * Represents all the different states the Back-end (specifically the backend.process.Engine) can be in.
 */
public enum SystemState {
    NOT_STARTED,//initially when the StanfordCoreNLP has not been started
    STARTED,//when the StanfordCoreNLP has been started
    PROCESSING,//when the back-end is processing files
    PROCESSED,//when it has processed all files
    FINISHED//when it has returned the Results of processing the Files
}
\end{lstlisting}
\section{src/main/java/frontend}
\subsection{Controllers}
\subsubsection{CustomResultRowController.java}
\begin{lstlisting}
package frontend.controllers;

import backend.process.FileData;
import backend.process.Result;
import frontend.dialogs.EditEventDialog;
import frontend.observers.DocumentReaderObserver;
import frontend.observers.TimelineRowObserver;
import javafx.event.ActionEvent;
import javafx.event.EventHandler;
import javafx.fxml.FXML;
import javafx.fxml.FXMLLoader;
import javafx.scene.Scene;
import javafx.scene.control.Button;
import javafx.scene.control.Dialog;
import javafx.scene.control.Label;
import javafx.scene.layout.Pane;
import javafx.scene.layout.VBox;
import javafx.stage.Stage;

import java.io.IOException;
import java.util.Optional;
import java.util.function.Consumer;

/**
 * Controller for the custom Result rows in the ListView of the RangeData layout.
 */
public class CustomResultRowController {
    @FXML
    private VBox rootVBox;
    @FXML
    private Button viewButton;
    @FXML
    private Button editButton;
    @FXML
    private Label subjectsLabel;
    @FXML
    private Label eventLabel;
    @FXML
    private Label fromLabel;
    private Result result;
    private TimelineRowObserver timelineRowObserver;
    private int rangePosition;
    /**
     * For the given Result, set up the data to be shown for this row, and the onclick events.
     * Builds a layout that shows the subjects, and the event of the given Result, along with buttons to view the
     * original document that produced this Result and the edit dialog to edit the Result.
     *
     * @param result the given Result.
     */
    public CustomResultRowController(Result result, TimelineRowObserver timelineRowObserver, int rangePosition) {
        this.result = result;
        this.timelineRowObserver = timelineRowObserver;
        this.rangePosition = rangePosition;
        FXMLLoader fxmlLoader = new FXMLLoader(getClass().getResource("customResultRow.fxml"));
        fxmlLoader.setController(this);
        try {
            rootVBox = fxmlLoader.load();//set the root layout
            setUpData();//set the data
            setUpOnClicks();//set the onclicks
        } catch (IOException e) {
            e.printStackTrace();
        }
    }

    /**
     * Set the data held by the Result in the appropriate labels.
     */
    private void setUpData() {
        if (result != null) {
            subjectsLabel.setText(result.getSubjectsAsString());
            eventLabel.setText(result.getEvent());
            if(result.getFileData() != null){
                FileData fileData = result.getFileData();
                fromLabel.setText(fileData.getFileName()+" ("+fileData.getCreationDateFormattedDayMonthYear()+")");
            }
        }
    }

    /**
     * Set the EventHandlers for the onClicks of the buttons of this Result row.
     */
    private void setUpOnClicks() {
        viewButton.setOnAction(new EventHandler<ActionEvent>() {
            @Override
            public void handle(ActionEvent event) {
                System.out.println("Go to Document button for timeline event: " + result.getTimelineDate() + " has been pressed");
                Stage stage = new Stage();
                DocumentReaderController documentReaderController = new DocumentReaderController(result, new DocumentReaderObserver() {
                    @Override
                    public void close() {
                        System.out.println("Closing document reader window");
                        stage.close();
                    }
                });
                Pane rootLayout = documentReaderController.getRootBorderPane();
                if (rootLayout != null) {
                    stage.setScene(new Scene(documentReaderController.getRootBorderPane(), 1024, 800));
                    stage.setTitle("Document Reader - " + result.getFileData().getFileName());
                    stage.show();
                }
            }
        });

        editButton.setOnAction(new EventHandler<ActionEvent>() {
            @Override
            public void handle(ActionEvent event) {
                System.out.println("Clicked Event on: " + result);
                Dialog dialog = EditEventDialog.getEditEventDialog(result, (rangePosition + 1));
                Optional<EditEventDialog.DialogResult> response = dialog.showAndWait();
                response.ifPresent(new Consumer<EditEventDialog.DialogResult>() {
                    @Override
                    public void accept(EditEventDialog.DialogResult dialogResult) {
                        if (dialogResult.getResultType() == EditEventDialog.DialogResult.ResultType.DELETE) {
                            System.out.println("Delete the event");
                            timelineRowObserver.delete(result);
                        } else if (dialogResult.getResultType() == EditEventDialog.DialogResult.ResultType.SAVE) {
                            System.out.println("Update the timeline");
                            Result copy = dialogResult.getResult();
                            timelineRowObserver.update(result, copy);
                        } else if (dialogResult.getResultType() == EditEventDialog.DialogResult.ResultType.CANCEL) {
                            System.out.println("Dont do anything");
                        }
                    }
                });
            }
        });
    }

    /**
     * Get the root layout represented by this Controller (i.e. labels showing the subjects and events held by this
     * Result, and buttons to view the document that produced the Result and edit the Result.
     *
     * @return the root layout represented by this Controller.
     */
    public Pane getRootLayout() {
        return rootVBox;
    }
}
\end{lstlisting}
\subsubsection{CustomTimelineRow.java}
\begin{lstlisting}
package frontend.controllers;

import backend.ranges.Range;
import frontend.observers.TimelineRowObserver;
import javafx.geometry.HPos;
import javafx.geometry.Orientation;
import javafx.geometry.VPos;
import javafx.scene.control.Separator;
import javafx.scene.layout.GridPane;
import javafx.scene.layout.Pane;

import java.util.ArrayList;
import java.util.List;

/**
 * Class for the Custom Rows in the Timeline. Builds a row made of a GridPane with the root Range (given in the
 * constructor) and recursively adds to the other coloumn the Ranges below this Range.
 */
public class CustomTimelineRow {
    private Range root;
    private GridPane rootLayout;
    private TimelineRowObserver timelineRowObserver;

    /**
     * Constructor used to set up the layout of the Timeline row. Afterwards a call to getPane() will return the row
     * to be added to the Timeline.
     *
     * @param root the root Range (which can hold Ranges as children which are added Recursively to the 2nd column of
     *             the GridPane.
     */
    public CustomTimelineRow(Range root, TimelineRowObserver timelineRowObserver) {
        this.timelineRowObserver = timelineRowObserver;
        this.root = root;
        setUpLayout();//set up the layout of this Timeline Row
    }

    /**
     * Called to create the Row based on the given Range (given to the constructor). The RootLayout (a root GridPane)
     * will have then been set.
     */
    private void setUpLayout() {
        ArrayList<Range> list = new ArrayList<>();
        list.add(root);
        rootLayout = getGridPane(list, false);//get the root GridPane
        rootLayout.setGridLinesVisible(true);//add grid lines to the root only (for optical purposes)
    }

    /**
     * For the given list of Range objects, add them to the first column of the GridPane (that will be returned), if we
     * need to add the vertical Separator (given by addVerticalSep) add it to the second column of the GridPane. If we
     * have multiple Ranges in the RangeList then they are separated by a horizontal Separator.
     * Recursively adds to the children of each Range in the list, in the third column of the row of their Parent, this
     * way the Parent contains all of its children in its row (and the children do the same for their children, etc).
     * <p>
     * Range data in column 1 and row i (given by their index in the list).
     * Children Range's (if they have) in column 2  and row i (same as parent).
     *
     * @param rangeList      the given list of Range objects.
     * @param addVerticalSep whether or not we need to add a Vertical Separator
     * @return the GridPane built based on the given list of Range objects.
     */
    private GridPane getGridPane(List<Range> rangeList, boolean addVerticalSep) {
        GridPane gridPane = loadGridPane();//get the root GridPane layout that we are adding to
        if (gridPane != null) {
            for (int i = 0; i < 2 * rangeList.size() - 1; i++) {
                //add a separator between them
                if (i % 2 == 1) {
                    Separator separator = new Separator(Orientation.HORIZONTAL);
                    gridPane.add(separator, 0, i, GridPane.REMAINING, 1);
                } else {
                    //layout for this Range
                    Range range = rangeList.get(i / 2);
                    Pane toAdd = rangeDataLayout(range, (i / 2));//add the layout for this given Range
                    gridPane.add(toAdd, 0, i);
                    GridPane.setValignment(toAdd, VPos.TOP);//set the items added to top left of the layout
                    GridPane.setHalignment(toAdd, HPos.LEFT);
                    if (range.getChildren().size() > 0) {//if this Range has children, then build their layout (recursively) and add them
                        if (addVerticalSep) {//we need to add a separator
                            Separator separator = new Separator(Orientation.VERTICAL);
                            gridPane.add(separator, 1, i);//add the separator
                        }
                        GridPane gridPane1 = getGridPane(range.getChildren(), true);
                        gridPane.add(gridPane1, 2, i);
                    }
                }
            }
        }
        return gridPane;//return the gridpane layout we added to
    }

    /**
     * Get the layout to show the data for the given Range (ie its Date range and its Results)
     *
     * @param range the given Range.
     * @return the layout representing the data of the given Range.
     */
    private Pane rangeDataLayout(Range range, int position) {
        RangeDataController rangeDataController = new RangeDataController(range, timelineRowObserver, position);
        return rangeDataController.getRootBorderPane();
    }

    /**
     * Get the base GridPane used to add a Range's data and its Children data (recursively). The GridPane is of the size
     * of its content, but it can be increased to any size.
     *
     * @return a base GridPane used to add a Range's data and its Children data.
     */
    private GridPane loadGridPane() {
        GridPane gridPane = new GridPane();
        gridPane.setPrefSize(GridPane.USE_COMPUTED_SIZE, GridPane.USE_COMPUTED_SIZE);
        gridPane.setMinSize(GridPane.USE_PREF_SIZE, GridPane.USE_PREF_SIZE);
        gridPane.setMaxSize(Double.MAX_VALUE, Double.MAX_VALUE);
        return gridPane;
    }


    /**
     * Get the root layout of this row in the Timeline (i.e. a GridPane that holds the information of the Root range,
     * passed in the timeline, in the first column with the Roots childrens information in the second column, and them
     * holding their children in their columns and so on).
     *
     * @return the root layout of a row in the Timeline given by the Range passed into the constructor.
     */
    public Pane getPane() {
        return rootLayout;
    }

}
\end{lstlisting}
\subsubsection{DocumentLoadedRowController.java}
\begin{lstlisting}
package frontend.controllers;

import backend.process.FileData;
import frontend.observers.DocumentsLoadedObserver;
import javafx.event.EventHandler;
import javafx.fxml.FXML;
import javafx.fxml.FXMLLoader;
import javafx.scene.control.Label;
import javafx.scene.image.ImageView;
import javafx.scene.input.MouseEvent;
import javafx.scene.layout.GridPane;

import java.io.IOException;

/**
 * Controller for each row in the Documents Loaded listview.
 */
public class DocumentLoadedRowController {
    @FXML
    private ImageView removingImageView;
    @FXML
    private Label documentLabel;
    @FXML
    private GridPane gridPane;
    private DocumentsLoadedObserver documentsLoadedObserver;

    /**
     * Creates a Controller and loads the layout for a row in the Loaded Documents listview. Observer is informed when
     * that file needs to be removed.
     *
     * @param fileData                FileData for which this row is made for.
     * @param documentsLoadedObserver observer that gets notified when the Results for the given FileData need to be removed.
     */
    public DocumentLoadedRowController(FileData fileData, DocumentsLoadedObserver documentsLoadedObserver) {
        this.documentsLoadedObserver = documentsLoadedObserver;
        FXMLLoader fxmlLoader = new FXMLLoader(getClass().getResource("documentLoadedRow.fxml"));
        fxmlLoader.setController(this);
        try {
            fxmlLoader.load();
        } catch (IOException e) {
            e.printStackTrace();
        }
        setData(fileData);
        removingImageView.setOnMouseClicked(new EventHandler<MouseEvent>() {
            @Override
            public void handle(MouseEvent event) {
                System.out.println("Remove this Document: " + fileData);
                documentsLoadedObserver.remove(fileData);
            }
        });
    }

    /**
     * For the label that is supposed to show the File name, set it.
     *
     * @param fileData label for which we set the File name.
     */
    private void setData(FileData fileData) {
        documentLabel.setText(fileData.getFileName()+" ("+fileData.getCreationDateFormattedDayMonthYear()+")");
    }

    /**
     * Get the root layout for this row (i.e. its GridPane, under which everything is set).
     *
     * @return the root GridPane layout for this row.
     */
    public GridPane getGridPane() {
        return gridPane;
    }
}
\end{lstlisting}
\subsubsection{DocumentReaderController.java}
\begin{lstlisting}
package frontend.controllers;

import backend.process.FileData;
import backend.process.ProcessFiles;
import backend.process.Result;
import frontend.observers.DocumentReaderObserver;
import javafx.beans.value.ChangeListener;
import javafx.beans.value.ObservableValue;
import javafx.event.ActionEvent;
import javafx.event.EventHandler;
import javafx.fxml.FXML;
import javafx.fxml.FXMLLoader;
import javafx.scene.control.Alert;
import javafx.scene.control.ButtonType;
import javafx.scene.control.ContextMenu;
import javafx.scene.control.MenuItem;
import javafx.scene.layout.BorderPane;
import org.fxmisc.richtext.InlineCssTextArea;

import java.awt.*;
import java.awt.datatransfer.StringSelection;
import java.io.File;
import java.io.IOException;

/**
 * Controller for the layout of the Document Reader
 */
public class DocumentReaderController {
    @FXML
    private InlineCssTextArea documentInlineCssTextArea;
    @FXML
    private MenuItem closeMenuItem;
    @FXML
    private MenuItem copyMenuItem;
    @FXML
    private BorderPane rootBorderPane;
    private DocumentReaderObserver documentReaderObserver;

    /**
     * Called to create the layout for the Document Reader. It creates a text area with the text of the File where the
     * given Result originates from, and it highlights the specific sentence that produced the given Result.
     * The Observer is used to inform the creator of the window that uses this layout, to inform them when the Close
     * menu item was pressed (to close the window).
     *
     * @param result                 the given Result.
     * @param documentReaderObserver the Observer that holds this layout, to inform when to close the window.
     */
    public DocumentReaderController(Result result, DocumentReaderObserver documentReaderObserver) {
        this.documentReaderObserver = documentReaderObserver;
        FXMLLoader fxmlLoader = new FXMLLoader(getClass().getResource("documentReader.fxml"));
        fxmlLoader.setController(this);
        try {
            fxmlLoader.load();
        } catch (IOException e) {
            e.printStackTrace();
        }
        setData(result);
        setOnActionListeners();
    }


    /**
     * Called to set the data displayed in the layout, which in this case is the text of the File set in the text area
     * with the relevant sentence that produced the given result highlighted.
     *
     * @param result the given Result (used to determine the File and sentence that produced this Result).
     */
    private void setData(Result result) {
        String stringInFile;
        if ((stringInFile = getStringInFile(result.getFileData())) != null) {
            documentInlineCssTextArea.clear();
            documentInlineCssTextArea.replaceText(stringInFile);
            ContextMenu contextMenu = new ContextMenu();
            MenuItem menuItemCopy = new MenuItem("Copy");
            menuItemCopy.setDisable(true);//initially cant copy as no text is selected
            copyMenuItem.setDisable(true);
            menuItemCopy.setOnAction(new EventHandler<ActionEvent>() {
                @Override
                public void handle(ActionEvent event) {
                    copy();//copy the selected text
                }
            });
            contextMenu.getItems().setAll(menuItemCopy);
            documentInlineCssTextArea.selectedTextProperty().addListener(new ChangeListener<String>() {
                @Override
                public void changed(ObservableValue<? extends String> observable, String oldValue, String newValue) {
                    if (newValue.equals("")) {
                        menuItemCopy.setDisable(true);
                        copyMenuItem.setDisable(true);
                    } else {
                        menuItemCopy.setDisable(false);
                        copyMenuItem.setDisable(false);
                    }
                }
            });
            documentInlineCssTextArea.setContextMenu(contextMenu);
            highlightText(result.getOriginalString(), stringInFile);
        } else {
            System.out.println("File is unavailable, cant be read");
            documentInlineCssTextArea.replaceText("");
            Alert documentUnvailable = documentUnavailableDialog(result.getFileData());
            documentUnvailable.showAndWait();
            //should close this window, as the file is unavailable
            rootBorderPane = null;
        }
    }

    /**
     * For the given FileData (holds the Files name, and represents it), produce an Alert Dialog to show to the User,
     * to inform them that the File which they wish to read is unavailable.
     *
     * @param fileData the given FileData
     * @return the Alert Dialog the informs the User the File is unavailable to read.
     */
    private Alert documentUnavailableDialog(FileData fileData) {
        Alert alert = new Alert(Alert.AlertType.INFORMATION);
        alert.setTitle("Document Unavailable");
        alert.setHeaderText(null);
        alert.setContentText("The File: " + fileData.getFileName() + " is Unavailable");
        alert.getDialogPane().getButtonTypes().setAll(ButtonType.OK);
        return alert;
    }

    /**
     * Called to set the action listeners for the Menu Items: Close and Copy
     */
    private void setOnActionListeners() {
        copyMenuItem.setOnAction(new EventHandler<ActionEvent>() {
            @Override
            public void handle(ActionEvent event) {
                copy();
            }
        });
        closeMenuItem.setOnAction(new EventHandler<ActionEvent>() {
            @Override
            public void handle(ActionEvent event) {
                close();
            }
        });
    }

    /**
     * Called to highlight the text "textToHighlight" in the "from" text, in the text area of this layout.
     *
     * @param textToHighlight the given text that needs to be highlighted. Part of the "from" text
     * @param from            the full text, that contains the text that needs to be highlighted.
     */
    private void highlightText(String textToHighlight, String from) {//highlight first occurrence of the text to highlight
        int startHighlight = from.indexOf(textToHighlight);
        int endHighlight = startHighlight + textToHighlight.length();//as we highlight the original sentence, which we know its length, and its start index
        if (endHighlight > from.length()) {//if we are over the limit of the text, then we go up to that point
            endHighlight = documentInlineCssTextArea.getText().length();
        }
        System.out.println("EndHighlight: " + endHighlight);
        if (startHighlight >= 0 && startHighlight < from.length() && endHighlight >= 0) {
            documentInlineCssTextArea.setStyle(startHighlight, endHighlight, "-fx-fill: blue; -fx-font-weight: bold");
        }
    }

    /**
     * Called to extract the text in the File which the given FileData represents.
     *
     * @param fileData the FileData that represents the File from which we are extracting text from.
     * @return null, if the File has been moved/deleted or we don't have read rights; otherwise the text of the File is
     * returned.
     */
    private String getStringInFile(FileData fileData) {
        if (fileData != null) {
            File file = new File(fileData.getFilePath());
            if (file.exists() && file.isFile() && file.canRead()) {
                ProcessFiles processFiles = new ProcessFiles();
                return processFiles.getTextInFile(file);
            }
        }
        return null;
    }

    /**
     * Called when the Close Menu Item is pressed. Inform the observer that they need to close the window in which this
     * layout resides in.
     */
    private void close() {
        System.out.println("Close Window");
        if (documentReaderObserver != null) {
            documentReaderObserver.close();
        }
    }

    /**
     * Called when the Copy Menu Item is pressed. Used to copy the text in the file to the clipboard.
     */
    private void copy() {
        if (documentInlineCssTextArea != null) {
            System.out.println("Copy Text to Clipboard: " + documentInlineCssTextArea.getSelectedText());
            Toolkit toolkit = Toolkit.getDefaultToolkit();
            if (toolkit != null && documentInlineCssTextArea.getSelectedText() != null) {
                toolkit.getSystemClipboard().setContents(new StringSelection(documentInlineCssTextArea.getSelectedText()), null);
            }
        }
    }

    /**
     * Get the BorderPane layout of the layout produced when a new object is made (that was populated with the given
     * Result object). This is the root layout.
     *
     * @return the BorderPane (root) layout, produced when the constructor is invoked.
     */
    public BorderPane getRootBorderPane() {
        return rootBorderPane;
    }
}
\end{lstlisting}
\subsubsection{EditEventController.java}
\begin{lstlisting}
package frontend.controllers;

import backend.process.Result;
import frontend.helpers.TextFieldState;
import frontend.observers.EditEventDialogObserver;
import javafx.beans.value.ChangeListener;
import javafx.beans.value.ObservableValue;
import javafx.css.PseudoClass;
import javafx.event.ActionEvent;
import javafx.event.EventHandler;
import javafx.fxml.FXML;
import javafx.fxml.FXMLLoader;
import javafx.scene.control.*;
import javafx.scene.image.ImageView;
import javafx.scene.input.MouseEvent;
import javafx.scene.layout.GridPane;
import javafx.scene.layout.HBox;

import java.io.IOException;
import java.text.ParseException;
import java.text.SimpleDateFormat;
import java.util.*;

/**
 * Controller of the Content in the EditEventDialog
 */
public class EditEventController {
    private final static int charLimit = 100;//100 characters should be enough for the user to write a summary of a sentence (might reduce it to less)
    private final static String maxChar = "Maximum 100 characters ";
    @FXML
    private TextField firstDateTextField;
    @FXML
    private TextField secondDateTextField;
    @FXML
    private TextField subjectTextField;
    @FXML
    private ImageView addImageView;
    @FXML
    private TextArea eventTextArea;
    @FXML
    private HBox subjectsHBox;
    @FXML
    private GridPane rootGridPane;
    @FXML
    private Label maxCharacterLabel;
    private Result result;
    private SimpleDateFormat simpleDateFormat = new SimpleDateFormat("dd-MM-yyyy");
    private SimpleDateFormat validInputFormat = new SimpleDateFormat("dd-MM-yyyy");
    private final PseudoClass errorClass = PseudoClass.getPseudoClass("error");
    private EditEventDialogObserver editEventDialogObserver;
    private ArrayList<TextFieldState> textFieldStates;
    private final static int subjectLimit = 5;
    private int subjectCount = 0;


    /**
     * Constructor to create the layout for the Content of the EditEventDialog.
     *
     * @param result                  the Result object used to populate the fields of the layout of the Dialog content.
     * @param editEventDialogObserver the Observer of this content (the holder of the Dialog), to inform when to enable
     *                                and disable the buttons (which aren't part of the content of the dialog, but
     *                                separate).
     */
    public EditEventController(Result result, EditEventDialogObserver editEventDialogObserver) {
        this.result = result;
        this.editEventDialogObserver = editEventDialogObserver;
        textFieldStates = new ArrayList<>(3);
        FXMLLoader fxmlLoader = new FXMLLoader(getClass().getResource("editEventDialog.fxml"));
        fxmlLoader.setController(this);
        try {
            fxmlLoader.load();
        } catch (IOException e) {
            e.printStackTrace();
        }
        setData();
        addListeners();
    }

    /**
     * For the Result object that has been set, populate the fields.
     */
    private void setData() {
        //for the set result object populate the fields
        if (result != null) {
            firstDateTextField.setText(getStringFromDate(result.getTimelineDate()
            		.getDate1()));
            firstDateTextField.getStylesheets().add(getClass()
            		.getResource("customErrorFields.css").toExternalForm());
            textFieldStates.add(0, TextFieldState.CORRECT);
            secondDateTextField.setText(getStringFromDate(result.getTimelineDate()
            		.getDate2()));
            secondDateTextField.getStylesheets().add(getClass()
            		.getResource("customErrorFields.css").toExternalForm());
            textFieldStates.add(1, TextFieldState.CORRECT);
            subjectsHBox.getChildren().setAll(getSubjectLabels(result.getSubjects()));
            eventTextArea.setText(result.getEvent());
            eventTextArea.getStylesheets().add(getClass()
            		.getResource("customErrorFields.css").toExternalForm());
            textFieldStates.add(2, TextFieldState.CORRECT);
        }
    }

    /**
     * For the set of Strings (which should be Subjects) produce a list of Labels (to add to the HBox Subjects layout).
     *
     * @param subjects the set of Strings.
     * @return a list of Labels (corresponding to each each item in the set and a listener to show a dialog to remove the
     * subject from the Result).
     */
    private List<Label> getSubjectLabels(Set<String> subjects) {
        List<Label> labels = new ArrayList<>();
        subjectCount = subjects.size();
        for (String subject : subjects) {
            Label label = new Label(subject);
            label.setUnderline(true);
            label.setOnMouseClicked(new EventHandler<MouseEvent>() {
                @Override
                public void handle(MouseEvent event) {
                    System.out.println("Remove this label and subject from the Result: " + subject);
                    if (shouldDeleteSubjectLabel(subject)) {
                        System.out.println("Removed: " + subject);
                        result.getSubjects().remove(subject);//remove the given subject
                        System.out.println("Remaining: " + result.getSubjects());//need to update the subjects shown
                        subjectsHBox.getChildren()
                        	.setAll(getSubjectLabels(result.getSubjects()));//so get the new labels and set them

                    }
                }
            });
            labels.add(label);
        }
        return labels;
    }

    /**
     * For a given Date (which can be null), return it in string format-
     *
     * @param date the given Date.
     * @return an empty String if the given date is null; otherwise the date as a String in the format dd-MM-yyyy.
     */
    private String getStringFromDate(Date date) {
        return (date != null) ? simpleDateFormat.format(date) : "";
    }

    /**
     * Add the listeners (event handlers) to the input fields (text fields, and the image view).
     */
    private void addListeners() {
        //for the fields add their event listeners (eg iamgeview add to hbox of label)
        addImageView.setOnMouseClicked(new EventHandler<MouseEvent>() {
            @Override
            public void handle(MouseEvent event) {
                System.out.println("Add subject in textfield to subject list: " + subjectTextField.getText());
                addSubjectAndClear(subjectTextField);
            }
        });
        subjectTextField.setOnAction(new EventHandler<ActionEvent>() {
            @Override
            public void handle(ActionEvent event) {
                System.out.println("Enter is pressed, add subject to list: " + subjectTextField.getText());
                addSubjectAndClear(subjectTextField);
            }
        });
        firstDateTextField.focusedProperty().addListener(new ChangeListener<Boolean>() {
            @Override
            public void changed(ObservableValue<? extends Boolean> observable, Boolean oldValue, Boolean newValue) {
                //validate date before setting it
                if (isValidDateInput(firstDateTextField.getText())) {
                    firstDateTextField.pseudoClassStateChanged(errorClass, false);
                    result.getTimelineDate().setDate1(getDate(firstDateTextField.getText().trim()));
                    textFieldStates.set(0, TextFieldState.CORRECT);
                    enableSave();
                } else {
                    firstDateTextField.pseudoClassStateChanged(errorClass, true);
                    textFieldStates.set(0, TextFieldState.WRONG);
                    disableSave();
                }
            }
        });
        secondDateTextField.focusedProperty().addListener(new ChangeListener<Boolean>() {
            @Override
            public void changed(ObservableValue<? extends Boolean> observable, Boolean oldValue, Boolean newValue) {
                //validate date before setting it
                if (secondDateTextField.getText().trim().length() == 0) {
                    //no second date
                    result.getTimelineDate().setDate2(null);
                    secondDateTextField.pseudoClassStateChanged(errorClass, false);
                    textFieldStates.set(1, TextFieldState.CORRECT);
                    enableSave();
                } else if (isValidDateInput(secondDateTextField.getText()) && getDate(secondDateTextField.getText().trim())
                		.compareTo(result.getTimelineDate().getDate1()) > 0) {
                    //valid input
                    secondDateTextField.pseudoClassStateChanged(errorClass, false);
                    result.getTimelineDate()
                    	.setDate2(getDate(secondDateTextField.getText().trim()));
                    textFieldStates.set(1, TextFieldState.CORRECT);
                    enableSave();
                } else {
                    secondDateTextField.pseudoClassStateChanged(errorClass, true);//disable button (need observer of this to tell)
                    textFieldStates.set(1, TextFieldState.WRONG);
                    disableSave();
                }
            }
        });
        eventTextArea.lengthProperty().addListener(new ChangeListener<Number>() {
            @Override
            public void changed(ObservableValue<? extends Number> observable, Number oldValue, Number newValue) {
                if (updateLabel(newValue.intValue())) {
                    textFieldStates.set(2, TextFieldState.WRONG);
                    eventTextArea.pseudoClassStateChanged(errorClass, true);
                    disableSave();
                } else {
                    textFieldStates.set(2, TextFieldState.CORRECT);
                    eventTextArea.pseudoClassStateChanged(errorClass, false);
                    enableSave();
                }
            }
        });
        eventTextArea.focusedProperty().addListener(new ChangeListener<Boolean>() {
            @Override
            public void changed(ObservableValue<? extends Boolean> observable, Boolean oldValue, Boolean newValue) {
                //save the event once we leave it
                if (textFieldStates.size() == 3 && textFieldStates.get(2) == TextFieldState.CORRECT) {
                    System.out.println("Saving result event: " + eventTextArea.getText());
                    result.setEvent(eventTextArea.getText());
                }
            }
        });
    }

    /**
     * Update the label that shows how many characters the users has left to type into the text area, or by how many
     * characters they have exceeded the limit. Returns whether or not the user has exceeded the character limit.
     *
     * @param charAmount the amount of characters the user has written in the textarea.
     * @return true if the charAmount has exceeded the maximum number of characters.
     */
    private boolean updateLabel(int charAmount) {
        String textToSet = maxChar;
        if (charAmount > charLimit) {
            textToSet += "(" + (charAmount - charLimit) + " too many)";
        } else {
            textToSet += "(" + (charLimit - charAmount) + " remaining)";
        }
        maxCharacterLabel.setText(textToSet);
        return charAmount > charLimit;
    }

    /**
     * For the given TextField, add its text to the list of Subjects of the Result (including showing its label), and
     * clear the field (as the subjects has already been added).
     *
     * @param textFieldSubject the givenTextField.
     */
    private void addSubjectAndClear(TextField textFieldSubject) {
        if (canAdd()) {
            addTextToSubjects(textFieldSubject.getText());
            textFieldSubject.setText("");
        } else {
            //inform the user they cant add more subjects
            Alert maxNumberSubjectDialog = getInformMaxSubject();
            maxNumberSubjectDialog.show();
        }
    }

    /**
     * For the given String (which should be a Subject), add it to Subjects of the Result (and update the list of
     * Subjects that is being shown).
     *
     * @param newSubject the given String.
     */
    private void addTextToSubjects(String newSubject) {
        if (newSubject.length() >= 1) {
            result.addSubject(newSubject);
            List<Label> labels = getSubjectLabels(result.getSubjects());
            subjectsHBox.getChildren().setAll(labels);
        }
    }

    /**
     * For the given String, produce a Date (assuming the String is in the format: dd-MM-yyyy), else null will be
     * returned.
     *
     * @param date the given String.
     * @return a non-null Date object if the given String is of the format: dd-MM-yyyy; otherwise null.
     */
    private Date getDate(String date) {
        Date toReturn = null;
        try {
            toReturn = validInputFormat.parse(date);
        } catch (ParseException e) {
            e.printStackTrace();
        }
        return toReturn;
    }

    /**
     * Called to disable the save button in the Dialog.
     */
    private void disableSave() {
        //tell the dialog to disable save
        editEventDialogObserver.disableSave(true);
    }

    /**
     * Called to enable the save button in the Dialog.
     */
    private void enableSave() {
        boolean canSave = true;
        for (TextFieldState textFieldState : textFieldStates) {
            if (textFieldState == TextFieldState.WRONG) {
                canSave = false;
                break;
            }
        }
        if (canSave) {
            //tell the dialog to enable save
            editEventDialogObserver.disableSave(false);
        }
    }

    /**
     * Checks whether the given input is of the valid format (dd-MM-yyyy).
     *
     * @param input a String date
     * @return true if the input is of the format dd-MM-yyyy, false otherwise.
     */
    private boolean isValidDateInput(String input) {
        try {
            validInputFormat.setLenient(false);
            validInputFormat.parse(input);
            return input.length() == 10;
        } catch (ParseException e) {
            return false;
        }
    }

    /**
     * Getter method for the root layout that this Controller is representing, which is a GridPane.
     *
     * @return the GridPane of the layout this controller is representing, which is the root layout.
     */
    public GridPane getRootGridPane() {
        return rootGridPane;
    }

    /**
     * Creates a Confirmation Alert Dialog, to allow the user to choose whether the given Subject String should be
     * removed from the Result. The users response is returned.
     *
     * @param subject the given Subject String.
     * @return true if the YES option is picked in the Dialog (that is the subject should be deleted); false otherwise
     */
    private boolean shouldDeleteSubjectLabel(String subject) {
        Alert confirmationDeleteDialog = new Alert(Alert.AlertType.CONFIRMATION);
        confirmationDeleteDialog.setTitle("Deleting a Subject");
        confirmationDeleteDialog.setContentText("Are you sure you want to delete: \"" + subject + "\" from the subjects of this event?");
        confirmationDeleteDialog.getButtonTypes().setAll(ButtonType.YES, ButtonType.NO, ButtonType.CANCEL);
        Optional<ButtonType> response = confirmationDeleteDialog.showAndWait();
        return response.get() == ButtonType.YES;
    }

    /**
     * Whether or not a Subject can be added to the list of Subjects of the Result object.
     *
     * @return true if the current subject count is below the subjectLimit (5); false otherwise.
     */
    private boolean canAdd() {
        return subjectCount < subjectLimit;
    }

    /**
     * Get a Information Alert Dialog that informs the user that the subject limit has been reached, and that they can
     * remove subjects by pressing them in the list (to make space for the new subject). This is so that the user doesn't
     * add too many subjects for an event which was given by just 1 sentence.
     *
     * @return a Information Dialog to inform the user the subject limit has been reached.
     */
    private Alert getInformMaxSubject() {
        Alert alert = new Alert(Alert.AlertType.INFORMATION);
        alert.setTitle("Add Subject Information");
        alert.setHeaderText(null);
        alert.setContentText("The maximum number of Subjects allowed has been reached. To remove a Subject just click on it.");
        return alert;
    }

}
\end{lstlisting}
\subsubsection{ListViewController.java}
\begin{lstlisting}
package frontend.controllers;

import backend.helpers.Sort;
import backend.process.FileData;
import backend.process.Result;
import backend.ranges.ProduceRanges;
import backend.ranges.Range;
import frontend.dialogs.LoadingDialog;
import frontend.dialogs.RemoveConfirmationDialog;
import frontend.observers.DocumentsLoadedObserver;
import frontend.observers.TimelineObserver;
import frontend.observers.TimelineRowObserver;
import javafx.collections.FXCollections;
import javafx.collections.ObservableList;
import javafx.event.ActionEvent;
import javafx.event.EventHandler;
import javafx.fxml.FXML;
import javafx.fxml.Initializable;
import javafx.geometry.Insets;
import javafx.scene.Node;
import javafx.scene.control.*;
import javafx.scene.input.ScrollEvent;
import javafx.scene.layout.StackPane;
import javafx.scene.layout.VBox;
import javafx.util.Callback;

import java.net.URL;
import java.util.*;


/**
 * Controller for the layout where the ListView is shown. Allows the listview to be populated with Result data.
 */
public class ListViewController implements Initializable, MenuBarControllerInter, DocumentsLoadedObserver, TimelineRowObserver {
    private enum ViewType {
        RANGE, DATE
    }

    private final static double MAX_ZOOM = 1.0d;
    private final static double MIN_ZOOM = 0.75d;

    @FXML
    private StackPane stackPane;
    @FXML
    private VBox vBox;
    @FXML
    private ListView<Object> timelineListView;
    @FXML
    private Button loadDocumentsButton;
    @FXML
    private Button saveToButton;
    @FXML
    private ListView<FileData> documentListView;
    @FXML
    private RadioMenuItem dateView;
    @FXML
    private RadioMenuItem rangeView;
    @FXML
    private ScrollPane scrollPane;
    private List<Result> results;
    private List<FileData> fileDatas;
    private ObservableList<Object> timelineObservableList = FXCollections.observableArrayList();
    private ObservableList<FileData> documentsLoadedObservableList = FXCollections.observableArrayList();
    private TimelineObserver timelineObserver;
    private LoadingDialog loadingDialog;
    private ToggleGroup radioMenuItemGroup;
    private ViewType viewType = ViewType.DATE;


    /**
     * Called when the layout is created.
     */
    @Override
    public void initialize(URL location, ResourceBundle resources) {
        System.out.println("Initialised timelineListView");
        System.out.println("documentListView: " + documentListView);
        System.out.println("loadDocumentsButton: " + loadDocumentsButton);
        System.out.println("saveToButton: " + saveToButton);
        loadingDialog = new LoadingDialog(stackPane, vBox);//pass the root layout and main content layout to know where
        //to show the loading dialog, and what to disable.
        //set up the group of the radio buttons in the menu
        //by default the dateView is shown
        dateView.setSelected(true);
        radioMenuItemGroup = new ToggleGroup();
        rangeView.setToggleGroup(radioMenuItemGroup);
        dateView.setToggleGroup(radioMenuItemGroup);
        rangeView.setOnAction(new EventHandler<ActionEvent>() {
            @Override
            public void handle(ActionEvent event) {
                showRangeTimeline();
            }
        });
        dateView.setOnAction(new EventHandler<ActionEvent>() {
            @Override
            public void handle(ActionEvent event) {
                showDateTimeline();
            }
        });

    }

    /**
     * Called to set the Observer for this Scene.
     *
     * @param timelineObserver the Observer for this Timeline Scene.
     */
    public void setTimelineObserver(TimelineObserver timelineObserver) {
        this.timelineObserver = timelineObserver;
        loadDocumentsButton.setOnAction(new EventHandler<ActionEvent>() {
            @Override
            public void handle(ActionEvent event) {
                System.out.println("Load in Documents");
                if (timelineObserver != null) {
                    timelineObserver.loadDocuments();
                }
            }
        });
        saveToButton.setText("Save \nTo...");
        saveToButton.setOnAction(new EventHandler<ActionEvent>() {
            @Override
            public void handle(ActionEvent event) {
                if (timelineObserver != null) {
                    timelineObserver.saveTo(results);
                }
            }
        });
    }

    /**
     * For the input List, sort and reverse it.
     * O(n)
     *
     * @param results the input List.
     */
    private void sortAndReverse(List<Result> results) {
        Sort.sortByDate1(results);//sort the results by their date1 value
        Collections.reverse(results);
    }

    /**
     * For the given Node make it visible and enable it depending on shouldShow. If shouldShow is true then make the
     * node visible and enable it; if it false make the node invisible and disable it.
     *
     * @param node       the given Node.
     * @param shouldShow whether or not the Node should be shown and enabled or not.
     */
    private void show(Node node, boolean shouldShow) {
        node.setVisible(shouldShow);
        node.setDisable(!shouldShow);
    }

    /**
     * For the input List of Results, set it as the items of the TimelineList. The type of timeline shown in the
     * ListView is given by the ViewType (member of the class that is changed with the RadioMenuItems).
     * <p>
     * The ListView and its ObservableList both set their generic type to Object, this way we don't have to create a
     * separate ListView when we want to swap from Results to Ranges (and vice versa). As when a ListView is set to
     * use a list of Ranges (for example) it will set its type to that, so then when it attempts to set a list of
     * Results it will throw an exception of the Casting failing (because its type is Ranges but its attempting to cast
     * it to Result).
     * <p>
     * When a row needs to be created from an item in that list, its type is looked at to determine what layout to use.
     * <p>
     * In order to change between Timeline Views use a ViewType that sets which View to use, where the RadioMenuItems
     * change its value according, with the default Timeline view being the Date Timeline (showing Results
     * individually). To allow the ListView to show different views of lists that each hold different kind of data, the
     * list and the observable had to be made of type Object as otherwise swapping from a list of Results to a list of
     * Ranges and setting them to the ListView would throw an exception of the Casting failing of a Range to a Result
     * (this also happens vice versa).
     * <p>
     * The timeline for the Range view is a ScrollPane that can be zoomed in/out.
     *
     * @param results the input List.
     */
    private void setTimelineList(List<Result> results) {
        //should show loading dialog while its setting the timeline
        timelineObservableList.clear();
        //check what kind of view we need to show
        if (viewType == ViewType.RANGE) {
            //we need to show the range view, which supports zooming in and out (so we need to use a vbox and scrollpane,
            //to be able to zoom)
            ProduceRanges produceRanges = new ProduceRanges();
            produceRanges.produceRanges(results);//produce the results

            VBox listVBox = new VBox();//the vbox that holds the Ranges
            listVBox.setPadding(new Insets(10));
            for (Range range : produceRanges.getTrees()) {//for each tree build its layout and add it to the vbox (row by row)
                CustomTimelineRow customTimelineRow = new CustomTimelineRow(range, this);
                listVBox.getChildren().add(customTimelineRow.getPane());
            }
            scrollPane.setContent(listVBox);//set the content of the scrollpane
            scrollPane.addEventFilter(ScrollEvent.ANY, new EventHandler<ScrollEvent>() {//add the event handler for the zooming
                @Override
                public void handle(ScrollEvent event) {
                    if (event.isControlDown()) {//if we are holding the control button
                        double scale = getScale(event, listVBox);//get the value by which we scale
                        listVBox.setScaleX(scale);//and set it
                        listVBox.setScaleY(scale);
                        event.consume();
                    }
                }

                private double getScale(ScrollEvent scrollEvent, Node node) {
                    double scale = node.getScaleX() + scrollEvent.getDeltaY() / 100;
                    if (scale <= MIN_ZOOM) {//we only want the user to zoom out (not in, hence the scale is never over 1)
                        scale = MIN_ZOOM;
                    } else if (scale >= MAX_ZOOM) {
                        scale = MAX_ZOOM;
                    }
                    return scale;
                }
            });
            show(scrollPane, true);//show the scrollpane
            show(timelineListView, false);//and hide the timeline list view
            return;//no need to follow the rest
        } else if (viewType == ViewType.DATE) {//if we have to show a date timeline
            show(scrollPane, false);//hide the scrollpane
            show(timelineListView, true);//show the timeline
            timelineListView.getStylesheets().setAll(getClass().getResource("listViewTheme.css")
            		.toExternalForm());
            timelineObservableList.setAll(results);//add the results
        }
        //assuming the observable list items have been set
        timelineListView.setItems(timelineObservableList);
        timelineListView.setCellFactory(new Callback<ListView<Object>, ListCell<Object>>() {
            @Override
            public ListCell<Object> call(ListView<Object> param) {
                return new ListCell<Object>() {
                    /**
                     * Called whenever a row needs to be shown/created on the screen.
                     * @param item the Range object for which this row has to be displayed for.
                     * @param empty whether or nor the row is empty.
                     */
                    @Override
                    protected void updateItem(Object item, boolean empty) {
                        super.updateItem(item, empty);
                        if (item != null && !empty) {
                            if (item instanceof Range) {
                                Range range = (Range) item;
                                CustomTimelineRow customTimelineRow = new CustomTimelineRow(range, ListViewController.this);
                                setGraphic(customTimelineRow.getPane());
                            } else if (item instanceof Result) {
                                Result result = (Result) item;
                                TimelineRowController timelineRowController = new TimelineRowController(getIndex(), ListViewController.this);
                                timelineRowController.setData(result);
                                setGraphic(timelineRowController.getGroup());
                            } else {
                                setGraphic(null);
                            }
                        } else {
                            setGraphic(null);
                        }
                    }
                };
            }
        });
    }

    /**
     * For the input List of FileData, set it as the items of the Documents Loaded List.
     *
     * @param fileDatas the input List.
     */
    private void setDocumentListView(List<FileData> fileDatas) {
        documentsLoadedObservableList.clear();
        documentsLoadedObservableList.addAll(fileDatas);
        documentListView.setItems(documentsLoadedObservableList);
        documentListView.setCellFactory(new Callback<ListView<FileData>, ListCell<FileData>>() {
            @Override
            public ListCell<FileData> call(ListView<FileData> param) {
                return new ListCell<FileData>() {
                    @Override
                    protected void updateItem(FileData item, boolean empty) {
                        super.updateItem(item, empty);
                        if (item != null && !empty) {
                            DocumentLoadedRowController documentLoadedRowController = new DocumentLoadedRowController(item, ListViewController.this);
                            setGraphic(documentLoadedRowController.getGridPane());
                        } else {
                            setGraphic(null);
                        }
                    }
                };
            }
        });
    }

    /**
     * For the given Results and FileData, set the data of the timelineListView and documentListView with them.
     *
     * @param results   a list of Result objects which contain data to populate the rows of the timelineListView with.
     * @param fileDatas a list of FileData objects which needs to populate the rows in the documentListView.
     */
    public void setTimelineListView(List<Result> results, List<FileData> fileDatas) {
        this.results = results;
        if (viewType == ViewType.DATE) {//to not waste time sorting Results that will be sorted by their Ranges later anyways
            sortAndReverse(this.results);
        }
        setTimelineList(this.results);

        this.fileDatas = fileDatas;
        Collections.sort(this.fileDatas);
        setDocumentListView(this.fileDatas);
    }

    /**
     * For the given input, add it to their appropriate lists.
     *
     * @param results   a list of Result objects which contain data to add to the timelineListView.
     * @param fileDatas a list of FileData objects which needs to be added to the documentListView.
     */
    public void addToTimelineListView(List<Result> results, List<FileData> fileDatas) {
        cleanRepeatedResults(results, this.fileDatas);
        cleanRepeatedFileData(this.fileDatas, fileDatas);
        this.results.addAll(results);
        sortAndReverse(this.results);
        setTimelineList(this.results);

        this.fileDatas.addAll(fileDatas);
        Collections.sort(this.fileDatas);
        setDocumentListView(this.fileDatas);
    }

    /**
     * Called to remove from the newFileData list, all the FileData objects that are already present in the oldFileData.
     *
     * @param oldFileData the given oldFileData.
     * @param newFileData the given newFileData.
     */
    private void cleanRepeatedFileData(List<FileData> oldFileData, List<FileData> newFileData) {
        Iterator<FileData> newFileDataIterator = newFileData.iterator();
        while (newFileDataIterator.hasNext()) {//for each new result
            FileData fileData = newFileDataIterator.next();
            if (oldFileData.contains(fileData)) {//if its in the old list
                newFileDataIterator.remove();//then dont add it to it (so remove it)
            }
        }
    }

    /**
     * Called to remove all the Result objects in results list that have their FileData object present in the fileDatas
     * List.
     *
     * @param results   the given Result list.
     * @param fileDatas the given FileData list.
     */
    private void cleanRepeatedResults(List<Result> results, List<FileData> fileDatas) {
        Iterator<Result> resultIterator = results.iterator();
        while (resultIterator.hasNext()) {
            //for each result, check it with the filedata, if its filedata is already there, then remove it
            Result result = resultIterator.next();
            if (fileDatas.contains(result.getFileData())) {
                //this file data already exists, so remove the result
                resultIterator.remove();
                System.out.println("Removed: " + result + " from results to be added");
            }
        }
    }

    /**
     * Called when the Close menu item is pressed.
     */
    @Override
    public void close() {
        System.out.println("Close pressed");
        if (timelineObserver != null) {
            timelineObserver.close();
        }
    }

    /**
     * Called when the About menu item is pressed.
     */
    @Override
    public void about() {
        System.out.println("About pressed");
        if (timelineObserver != null) {
            timelineObserver.showAbout();
        }
    }

    /**
     * Called when the Timeline menu item is pressed.
     */
    @Override
    public void timeline() {
        System.out.println("Timeline pressed");
        if (timelineObserver != null) {
            timelineObserver.timeline();
        }
    }

    /**
     * When the preferences menu item is pressed.
     */
    @Override
    public void preferences() {
        if (timelineObserver != null) {
            timelineObserver.preferences();
        }
    }

    /**
     * Called when a given row in the Loaded Documents listview is removed.
     *
     * @param fileData the FileData for the File where we want to remove its results from the Timeline.
     */
    @Override
    public void remove(FileData fileData) {
        System.out.println("Need to remove: " + fileData);
        Alert removeConfirmationDialog = RemoveConfirmationDialog.getRemoveConfirmationDialog(fileData.getFileName());
        Optional<ButtonType> response = removeConfirmationDialog.showAndWait();
        if (response.isPresent() && response.get() == ButtonType.YES) {
            fileDatas.remove(fileData);
            removeResults(results, fileData);
            setTimelineListView(results, fileDatas);
        }//else dont remove the events related to that file (as we keep it)
    }

    /**
     * Removes the given FileData from the FileData list and all the Results linked to it in the Results list.
     *
     * @param results  list of Results for which we we need to delete the Results linked to the given FileData.
     * @param fileData FileData for which in the given Results list we need to remove the linked Results.
     */
    private void removeResults(List<Result> results, FileData fileData) {
        Iterator<Result> resultIterator = results.iterator();
        while (resultIterator.hasNext()) {
            Result result = resultIterator.next();
            if (result.getFileData().equals(fileData)) {
                System.out.println("Need to remove: " + result);
                resultIterator.remove();
            }
        }
    }

    /**
     * Called by a member of the ListView, to inform the ListView to update, as it updated its values.
     * The updated list can have this new Result row somewhere else as the user could have changed its date (and the list
     * is sorted by dates).
     *
     * @param updatedResult the Result object of the row that was edited.
     * @param position      the position this cell was in the ListView.
     */
    @Override
    public void update(Result updatedResult, int position) {
        System.out.println("Update the list");
        if (results.size() > position) {
            results.set(position, updatedResult);
            setTimelineListView(results, fileDatas);
        }

    }

    /**
     * Called by a member of the ListView, to inform the ListView to update, as it updated its values.
     * The updated list can have this new Result row somewhere else as the user could have changed its date (and the list
     * is sorted by dates).
     *
     * @param previous      the Result object that was previously in this Row.
     * @param updatedResult the updated Result object of the TimelineRow.
     */
    @Override
    public void update(Result previous, Result updatedResult) {
        int pos = results.indexOf(previous);
        System.out.println("Previous: " + pos);
        if (pos != -1) {
            update(updatedResult, pos);
        }
    }

    /**
     * Called by a member of the ListView, to inform the ListView that it needs to be removed from the List.
     * Thereby the list needs to be updated (i.e. set again).
     *
     * @param position the position of the event that needs to be deleted.
     */
    @Override
    public void delete(int position) {
        System.out.println("Deleting the event");
        if (results.size() > position) {
            results.remove(position);
            setTimelineListView(results, fileDatas);
        }
    }

    /**
     * Called by a member of the ListView, to inform the ListView that it needs to be removed from the List.
     * Thereby the list needs to be updated (i.e. set again).
     *
     * @param result the given event (Result) to be deleted.
     */
    @Override
    public void delete(Result result) {
        int pos = results.indexOf(result);
        System.out.println("pos: " + pos);
        if (pos != -1) {
            delete(pos);
        }
    }

    /**
     * Called to show the loading dialog. (Only if the layouts have been passed to LoadingDialog)
     */
    public void showLoadingDialog() {
        loadingDialog.showLoadingDialog();
    }

    /**
     * Called to remove the loading dialog.
     */
    public void removeLoadingDialog() {
        loadingDialog.removeLoadingDialog();
    }

    /**
     * Called to show the Timeline with the individual dates and events in separate rows.
     */
    private void showDateTimeline() {
        if (viewType != ViewType.DATE) {
            viewType = ViewType.DATE;
            //show what is shown
            setTimelineListView(results, fileDatas);
        }
    }

    /**
     * Called to show the Timeline with the Results grouped into their ranges and then shown.
     */
    private void showRangeTimeline() {
        if (viewType != ViewType.RANGE) {
            viewType = ViewType.RANGE;
            setTimelineListView(results, fileDatas);
        }
    }
}
\end{lstlisting}
\subsubsection{MenuBarControllerInter.java}
\begin{lstlisting}
package frontend.controllers;

/**
 * Implemented for Observers of a Controller that has a menu bar.
 */
public interface MenuBarControllerInter {
    /**
     * Called when the close menu item is selected.
     */
    void close();

    /**
     * Called when the about menu item is selected.
     */
    void about();

    /**
     * Called when the timeline menu item is selected.
     */
    void timeline();

    /**
     * Called when the preferences menu item is selected.
     */
    void preferences();
}
\end{lstlisting}
\subsubsection{RangeDataController.java}
\begin{lstlisting}
package frontend.controllers;

import backend.process.Result;
import backend.ranges.Range;
import frontend.observers.TimelineRowObserver;
import javafx.collections.FXCollections;
import javafx.collections.ObservableList;
import javafx.fxml.FXML;
import javafx.fxml.FXMLLoader;
import javafx.scene.control.Label;
import javafx.scene.control.ListCell;
import javafx.scene.control.ListView;
import javafx.scene.layout.Pane;
import javafx.scene.layout.VBox;
import javafx.util.Callback;

import java.io.IOException;

/**
 * Controller for the layout that represents the data held by a Range (i.e. its Data and its Results).
 */
public class RangeDataController {
    private Range range;
    @FXML
    private VBox rootVBox;
    @FXML
    private ListView<Result> resultsListView;
    @FXML
    private Label dateLabel;
    private ObservableList<Result> results;
    private TimelineRowObserver timelineRowObserver;
    private int rangePosition;

    /**
     * Constructor that sets up the layout to represent the data held by the given Range.
     *
     * @param range the given Range.
     */
    public RangeDataController(Range range, TimelineRowObserver timelineRowObserver, int rangePosition) {
        this.range = range;
        this.timelineRowObserver = timelineRowObserver;
        this.rangePosition = rangePosition;
        FXMLLoader fxmlLoader = new FXMLLoader(getClass().getResource("rangeDataLayout.fxml"));//load the base layout
        fxmlLoader.setController(this);
        try {
            rootVBox = fxmlLoader.load();
            setUp();//set the data in the layout
        } catch (IOException e) {
            e.printStackTrace();
        }
    }

    /**
     * Set up the data shown in this layout. This means setting the date that this Range layout is representing, and the
     * list of Result objects held by the given Range (in the constructor).
     */
    private void setUp() {
        if (range != null) {
            dateLabel.setText(range.getDateRange());//set the date text
            if (range.getResults().size() > 0) {//if we have results then show the list
                results = FXCollections.observableArrayList();
                results.addAll(range.getResults());//set up the observable list of Results used to add to the listView
                resultsListView.setItems(results);
                resultsListView.setCellFactory(new Callback<ListView<Result>, ListCell<Result>>() {
                    @Override
                    public ListCell<Result> call(ListView<Result> param) {
                        return new ListCell<Result>() {//set up the listcell used in this listview
                            @Override
                            protected void updateItem(Result item, boolean empty) {
                                super.updateItem(item, empty);
                                if (item != null) {//if we have a valid item, build its layout and set it
                                    CustomResultRowController customResultRowController = new CustomResultRowController(item, timelineRowObserver, rangePosition);
                                    setGraphic(customResultRowController.getRootLayout());//set the layout built with the result
                                } else {//dont have a valid result so dont show anything
                                    setGraphic(null);
                                }
                            }
                        };
                    }
                });
            } else {//dont have results then dont show the list
                rootVBox.getChildren().remove(resultsListView);
            }
        }
    }

    /**
     * Get the root layout that this Controller controls (ie a VBox with the Date of the given Range - passed in the
     * constructor; and a list of the Results held by the Range).
     *
     * @return the root layout that this Controller represents.
     */
    public Pane getRootBorderPane() {
        return rootVBox;
    }


}
\end{lstlisting}
\subsubsection{StartUpController.java}
\begin{lstlisting}
package frontend.controllers;

import frontend.dialogs.LoadingDialog;
import frontend.observers.StartUpObserver;
import javafx.fxml.FXML;
import javafx.fxml.Initializable;
import javafx.scene.control.Button;
import javafx.scene.layout.StackPane;
import javafx.scene.layout.VBox;

import java.net.URL;
import java.util.ResourceBundle;

/**
 * Controller for the Startup Scene. Holds a StartUpObserver, where it calls the relevant methods, when the Controller is
 * called. As the scene has a menu bar, it implements the MenuBarControllerInterface, by implementing the relevant methods
 * for the menu bar.
 */
public class StartUpController implements Initializable, MenuBarControllerInter {
    private final static String TAG = "STARTUPCONTROLLER: ";
    private StartUpObserver observer;
    @FXML
    private Button loadDocumentsButton;
    @FXML
    private StackPane stackPane;
    @FXML
    private VBox vBox;//main layout
    private LoadingDialog loadingDialog;

    /**
     * Called on creation of the Scene.
     */
    @Override
    public void initialize(URL location, ResourceBundle resources) {
        System.out.println(TAG + "StartUp fxml is starting to run");
        loadingDialog = new LoadingDialog(stackPane, vBox);
    }

    /**
     * Set the StartUpObserver to be used when the Controller has to process an action.
     *
     * @param observer the StartUpObserver linked to this Controller.
     */
    public void setObserver(StartUpObserver observer) {
        System.out.println("Observer has been set");
        this.observer = observer;
    }

    /**
     * Called when the button to load documents is pressed.
     */
    public void loadDocuments() {
        System.out.println(TAG + "load documents");
        if (observer != null) {
            observer.loadFiles();//should disable button to not load in more files while doing this
        }
    }

    /**
     * When the close menu item is pressed.
     */
    @Override
    public void close() {
        if (observer != null) {
            observer.close();
        }
    }

    /**
     * When the about menu item is pressed.
     */
    @Override
    public void about() {
        if (observer != null) {
            observer.showAbout();
        }
    }

    /**
     * When the timeline menu item is pressed.
     */
    @Override
    public void timeline() {
        if (observer != null) {
            observer.timeline();
        }
    }

    /**
     * When the preferences menu item is pressed.
     */
    @Override
    public void preferences() {
        if (observer != null) {
            observer.preferences();
        }
    }

    /**
     * Set whether or not the Load Documents Button should be disabled (so that it cannot be pressed)
     *
     * @param disable whether or not the Load Documents Button should be disabled.
     */
    public void setDisableLoadDocumentsButton(boolean disable) {
        loadDocumentsButton.setDisable(disable);
    }

    /**
     * Called to show the loading dialog. (Only if the layouts have been passed to LoadingDialog)
     */
    public void showLoadingDialog() {
        loadingDialog.showLoadingDialog();
    }

    /**
     * Called to remove the loading dialog.
     */
    public void removeLoadingDialog() {
        loadingDialog.removeLoadingDialog();
    }
}
\end{lstlisting}
\subsubsection{TimelineRowController.java}
\begin{lstlisting}
package frontend.controllers;

import backend.process.Result;
import frontend.dialogs.EditEventDialog;
import frontend.observers.DocumentReaderObserver;
import frontend.observers.TimelineRowObserver;
import javafx.event.ActionEvent;
import javafx.event.EventHandler;
import javafx.fxml.FXML;
import javafx.fxml.FXMLLoader;
import javafx.scene.Group;
import javafx.scene.Scene;
import javafx.scene.control.Button;
import javafx.scene.control.Dialog;
import javafx.scene.control.Label;
import javafx.scene.layout.BorderPane;
import javafx.scene.layout.Pane;
import javafx.stage.Stage;

import java.io.IOException;
import java.util.Optional;
import java.util.function.Consumer;

/**
 * Controller class for each row in the timeline listview.
 */
public class TimelineRowController {
    @FXML
    private Group group;
    @FXML
    private Label eventNumberLabel;
    @FXML
    private Label dateLabel;
    @FXML
    private Label subjectsLabel;
    @FXML
    private Label eventLabel;
    @FXML
    private BorderPane borderPane;
    @FXML
    private Button editButton;
    @FXML
    private Button viewButton;
    @FXML
    private Label fromLabel;
    private int position;
    private Result result;
    private TimelineRowObserver timelineRowObserver;

    /**
     * Loads the layout for the row (appropriate layout picked depending on whether row is even or not).
     *
     * @param position the position this row is in the timeline (to determine if its odd or even and to display its index).
     */
    public TimelineRowController(int position, TimelineRowObserver timelineRowObserver) {
        this.position = position;
        this.timelineRowObserver = timelineRowObserver;
        FXMLLoader fxmlLoader;
        boolean isEven = (position % 2) == 0;
        if (isEven) {
            fxmlLoader = new FXMLLoader(getClass().getResource("timelineRowEven.fxml"));
        } else {
            fxmlLoader = new FXMLLoader(getClass().getResource("timelineRowOdd.fxml"));
        }
        fxmlLoader.setController(this);
        try {
            fxmlLoader.load();
        } catch (IOException e) {
            e.printStackTrace();
        }
    }

    /**
     * For the row, where the layout has already been loaded, set the values at the labels, and the event handlers for
     * the buttons.
     *
     * @param result the Result object with wich the values of the labels will be populated (e.g. date, subjects, event,
     *               etc.).
     */
    public void setData(Result result) {
        this.result = result;
        eventNumberLabel.setText("Event #" + (position + 1));
        dateLabel.setText("Date: " + result.getTimelineDate().toString());
        subjectsLabel.setText("Subjects: " + result.getSubjectsAsString());
        eventLabel.setText("Event: " + result.getEvent());
        fromLabel.setText("From: " + result.getFileData().getFileName() + " (" + result.getFileData().getCreationDateFormattedDayMonthYear() + ")");
        borderPane.setStyle("-fx-border-color: black; -fx-border-width: 4; -fx-border-style: solid inside;");
        editButton.setOnAction(new EventHandler<ActionEvent>() {
            @Override
            public void handle(ActionEvent event) {
                System.out.println("Edit button for timeline event: " + result.getTimelineDate() + " has been pressed");
                Dialog dialog = EditEventDialog.getEditEventDialog(result, (position + 1));
                Optional<EditEventDialog.DialogResult> response = dialog.showAndWait();
                response.ifPresent(new Consumer<EditEventDialog.DialogResult>() {
                    @Override
                    public void accept(EditEventDialog.DialogResult dialogResult) {
                        if (dialogResult.getResultType() == EditEventDialog.DialogResult.ResultType.DELETE) {
                            System.out.println("Delete the event");
                            timelineRowObserver.delete(position);
                        } else if (dialogResult.getResultType() == EditEventDialog.DialogResult.ResultType.SAVE) {
                            System.out.println("Update the timeline");
                            Result copy = dialogResult.getResult();
                            timelineRowObserver.update(copy, position);
                        } else if (dialogResult.getResultType() == EditEventDialog.DialogResult.ResultType.CANCEL) {
                            System.out.println("Dont do anything");
                        }
                    }
                });
            }
        });
        viewButton.setOnAction(new EventHandler<ActionEvent>() {
            @Override
            public void handle(ActionEvent event) {
                System.out.println("Go to Document button for timeline event: " + result.getTimelineDate() + " has been pressed");
                Stage stage = new Stage();
                DocumentReaderController documentReaderController = new DocumentReaderController(result, new DocumentReaderObserver() {
                    @Override
                    public void close() {
                        System.out.println("Closing document reader window");
                        stage.close();
                    }
                });
                Pane rootLayout = documentReaderController.getRootBorderPane();
                if (rootLayout != null) {
                    stage.setScene(new Scene(documentReaderController.getRootBorderPane(), 1024, 800));
                    stage.setTitle("Document Reader - " + result.getFileData().getFileName());
                    stage.show();
                }
            }
        });
    }

    /**
     * After the layout has been loaded (and the data has been populated), return the layout to use it in the timeline.
     *
     * @return the parent layout of the row in the timeline.
     */
    public Group getGroup() {
        return group;
    }
}
\end{lstlisting}
\subsection{Dialogs}
\subsubsection{AboutDialog.java}
\begin{lstlisting}
package frontend.dialogs;

import javafx.scene.control.ButtonType;
import javafx.scene.control.Dialog;
import javafx.scene.control.Label;
import javafx.scene.layout.GridPane;

/**
 * Class to make the About Dialog that shows information of the Application.
 */
public class AboutDialog {
    private final static String AUTHOR = "Oliver Philip Hohn";
    private final static String APP_NAME = "Automated Timeline Extraction";
    private final static String DEFAULT_VERSION = "DEVELOPMENT";

    /**
     * Get a Dialog that shows the information of the System (the version number, who created it, etc.)
     *
     * @return a Dialog that shows the information of the System.
     */
    public Dialog getAboutDialog() {
        Dialog aboutDialog = new Dialog();
        aboutDialog.setTitle("About");
        aboutDialog.setHeaderText(null);

        //layout of body gridpanes
        GridPane gridPane = new GridPane();
        String version = (getClass().getPackage().getImplementationVersion() != null) ? getClass().getPackage().getImplementationVersion() : DEFAULT_VERSION;
        Label appNameLabel = new Label("App Name: ");
        Label appLabel = new Label(APP_NAME + " v" + version);
        Label authorLabel = new Label(AUTHOR);
        Label authorNameLabel = new Label("Created By ");
        gridPane.add(appNameLabel, 0, 0);
        gridPane.add(appLabel, 1, 0);
        gridPane.add(authorNameLabel, 0, 1);
        gridPane.add(authorLabel, 1, 1);
        aboutDialog.getDialogPane().setContent(gridPane);
        aboutDialog.getDialogPane().getButtonTypes().setAll(ButtonType.OK);
        return aboutDialog;
    }

}
\end{lstlisting}
\subsubsection{EditEventDialog.java}
\begin{lstlisting}
package frontend.dialogs;

import backend.process.Result;
import frontend.controllers.EditEventController;
import frontend.observers.EditEventDialogObserver;
import javafx.event.EventHandler;
import javafx.scene.Node;
import javafx.scene.control.Alert;
import javafx.scene.control.ButtonBar;
import javafx.scene.control.ButtonType;
import javafx.scene.control.Dialog;
import javafx.scene.input.MouseEvent;
import javafx.util.Callback;

import java.util.Optional;

/**
 * Class that creates an EditEvent Dialog and returns it.
 */
public class EditEventDialog {
    /**
     * Called to produce a edit Dialog for the given Result object. Any changes made do not change the given Result
     * object (as it is cloned).
     *
     * @param result   the Result object for which this Edit Dialog needs to be produced (i.e. to populate the fields).
     * @param position the position the Result is in the timeline (to show at the top of the Edit Dialog)
     * @return a Dialog  where if the user decides to Save the changes made then its result will be a new Result object with the changes made by the
     * user (which can be the same as the given Result object if no changes are made, or a new Result object with the data of the previous plus the changes made).
     * This Result object is encapsulated in a DialogResult, that holds the option selected by the user (SAVE, DELETE, CANCEL),
     * to determine what to do with the copy Result object and the original Result object in the ListView.
     */
    public static Dialog<DialogResult> getEditEventDialog(Result result, int position) {
        //make a copy of the passed in Result object, use it to change values, and pass that to return.
        Dialog<DialogResult> dialog = new Dialog<>();
        try {
            Result copyResult = (Result) result.clone();
            DialogResult dialogResult = new DialogResult(copyResult);
            dialog.setTitle("Editing Event");
            dialog.setHeaderText("Event #" + position);
            ButtonType buttonTypeSave = new ButtonType("Save", ButtonBar.ButtonData.OK_DONE);
            ButtonType buttonTypeCancel = new ButtonType("Cancel", ButtonBar.ButtonData.CANCEL_CLOSE);
            ButtonType buttonTypeDelete = new ButtonType("Delete", ButtonBar.ButtonData.LEFT);

            dialog.getDialogPane().getButtonTypes().addAll(buttonTypeDelete, buttonTypeSave, buttonTypeCancel);
            Node deleteButton = dialog.getDialogPane().lookupButton(buttonTypeDelete);
            deleteButton.setOnMousePressed(new EventHandler<MouseEvent>() {
                @Override
                public void handle(MouseEvent event) {
                    System.out.println("Delete Pressed");
                    //show dialog to delete event, to confirm etc, tell listener they need to delete this event
                    Alert alert = new Alert(Alert.AlertType.CONFIRMATION);
                    alert.setTitle("Delete");
                    alert.setHeaderText(null);
                    alert.setContentText("Are you sure you want to delete?");
                    alert.getButtonTypes().setAll(ButtonType.YES, ButtonType.NO);
                    Optional<ButtonType> result = alert.showAndWait();
                    if (result.isPresent() && result.get() == ButtonType.YES) {
                        //delete event
                        dialogResult.setResultType(DialogResult.ResultType.DELETE);
                        //then stop showing
                        if (dialog.isShowing()) {
                            dialog.close();
                        }
                    }
                }
            });
            dialog.getDialogPane().setContent(new EditEventController(copyResult, new EditEventDialogObserver() {
                @Override
                public void disableSave(boolean disableSave) {
                    Node saveButton = dialog.getDialogPane().lookupButton(buttonTypeSave);
                    saveButton.setDisable(disableSave);
                }
            }).getRootGridPane());
            dialog.setResultConverter(new Callback<ButtonType, DialogResult>() {
                @Override
                public DialogResult call(ButtonType param) {
                    if (param.equals(buttonTypeSave)) {
                        dialogResult.setResultType(DialogResult.ResultType.SAVE);
                    }
                    return dialogResult;
                }
            });
        } catch (CloneNotSupportedException e) {
            e.printStackTrace();
        }
        return dialog;
    }

    /**
     * Class that holds the option selected by the user in the Edit Event Dialog, and the Result object with the changes
     * made it to it through the input fields in the Dialog.
     */
    public static class DialogResult {
        /**
         * The different options the user has available in the Edit Event Dialog.
         */
        public enum ResultType {
            SAVE, CANCEL, DELETE
        }

        private Result result;
        private ResultType resultType;

        /**
         * To create a new DialogResult object.
         *
         * @param result the copy of the Result object that the user will modify with in the input fields.
         */
        DialogResult(Result result) {
            this.result = result;
            this.resultType = ResultType.CANCEL;
        }

        /**
         * Getter for the copy of the Result object that the user modified with the input fields.
         *
         * @return the copy of the Result object.
         */
        public Result getResult() {
            return result;
        }

        /**
         * Getter for the option selected by the User when they leave the Edit Event Dialog.
         *
         * @return the option selected.
         */
        public ResultType getResultType() {
            return resultType;
        }

        /**
         * Setter for the option selected by the User when they leave the Edit Event Dialog.
         *
         * @param resultType the option selected.
         */
        void setResultType(ResultType resultType) {
            this.resultType = resultType;
        }
    }
}
\end{lstlisting}
\subsubsection{FileConfirmationDialog.java}
\begin{lstlisting}
package frontend.dialogs;

import backend.process.FileData;
import frontend.helpers.TextFieldState;
import javafx.beans.value.ChangeListener;
import javafx.beans.value.ObservableValue;
import javafx.css.PseudoClass;
import javafx.geometry.Insets;
import javafx.geometry.Pos;
import javafx.scene.control.*;
import javafx.scene.layout.GridPane;
import javafx.scene.layout.Priority;

import java.text.ParseException;
import java.text.SimpleDateFormat;
import java.util.ArrayList;
import java.util.List;

/**
 * Used to create and the confirmation dialog that will contain the filenames, and editable base dates for the File.
 */
public class FileConfirmationDialog {
    private final PseudoClass errorClass = PseudoClass.getPseudoClass("error");
    private final SimpleDateFormat validInputFormat = new SimpleDateFormat("dd-MM-yyyy");
    private ArrayList<TextFieldState> textFieldStates;

    /**
     * For the List of FileData, produce an Alert Confirmation Dialog, where the filename and file creation date are
     * shown, so that the user can determine whether or not to use it as a base date.
     *
     * @param fileDatas the List of FileData for which we need to create the list of labels and text fields in the dialog.
     * @return an Alert object that can be shown for the user to confirm to use these dates for the given files.
     */
    public Alert getConfirmationFileDialog(List<FileData> fileDatas) {
        Alert alert = new Alert(Alert.AlertType.CONFIRMATION);
        alert.setTitle("Confirm the Base Dates");
        alert.setContentText("The following Base Dates will be used for the given Files. Change them appropriately: ");

        ScrollPane scrollPane = new ScrollPane();
        scrollPane.setPadding(new Insets(10));
        scrollPane.setFitToWidth(true);//its content should resize to fit the width of the scroll pane (if it gets resized)
        GridPane fileDateGridPane = new GridPane();
        fileDateGridPane.setMaxWidth(Double.MAX_VALUE);
        fileDateGridPane.setPadding(new Insets(10));
        fileDateGridPane.setHgap(10);//add spacing between the cells in the gridpane
        fileDateGridPane.setVgap(10);
        scrollPane.setContent(fileDateGridPane);
        textFieldStates = new ArrayList<>(fileDatas.size());
        for (int i = 0; i < fileDatas.size(); i++) {
            FileData fileData = fileDatas.get(i);
            Label fileNameLabel = new Label(fileData.getFileName());
            TextField dateTextField = new TextField(fileData.getCreationDateFormattedDayMonthYear());
            dateTextField.setPromptText("dd-MM-yyyy");//give the user a prompt in case they dont know what the input format is
            dateTextField.getStylesheets().add(getClass().getResource("customErrorFields.css")
            		.toExternalForm());
            dateTextField.setMinWidth(150);//150 pixels is enough for 12 characters
            dateTextField.setMaxHeight(35);//enough for one line, if text size is 15
            dateTextField.setAlignment(Pos.CENTER);
            textFieldStates.add(i, TextFieldState.CORRECT);
            GridPane.setVgrow(dateTextField, Priority.NEVER);//1 line always
            GridPane.setHgrow(dateTextField, Priority.ALWAYS);
            int finalI = i;//to input in the array list at this index
            dateTextField.focusedProperty().addListener(new ChangeListener<Boolean>() {
                @Override
                public void changed(ObservableValue<? extends Boolean> observable, Boolean oldValue, Boolean newValue) {
                    if (!newValue) {
                        //check validate text field input
                        String input = dateTextField.getText().trim();
                        System.out.println("Inputted: " + input);
                        if (isValidDateInput(input)) {
                            dateTextField.pseudoClassStateChanged(errorClass, false);
                            fileData.setCreationDate(input);
                            textFieldStates.set(finalI, TextFieldState.CORRECT);
                            if (isAllCorrect(textFieldStates)) {
                                alert.getDialogPane().lookupButton(ButtonType.OK)
                                	.setDisable(false);
                            }
                        } else {
                            dateTextField.pseudoClassStateChanged(errorClass, true);
                            textFieldStates.set(finalI, TextFieldState.WRONG);
                            alert.getDialogPane().lookupButton(ButtonType.OK)
                            		.setDisable(true);
                        }
                    }
                }
            });

            fileDateGridPane.add(fileNameLabel, 0, i);
            fileDateGridPane.add(dateTextField, 1, i);
        }
        alert.getDialogPane().setExpandableContent(scrollPane);
        //alert.getDialogPane().setPrefSize(480,250);
        return alert;
    }

    /**
     * Whether or not all the states in the given list of TextFieldStates are CORRECT (true) or not (false).
     *
     * @param textFieldStates a list of TextFieldStates to check
     * @return true if all TextFieldStates in the given input are CORRECT, false otherwise.
     */
    private boolean isAllCorrect(ArrayList<TextFieldState> textFieldStates) {
        boolean toReturn = true;
        for (TextFieldState textFieldState : textFieldStates) {
            if (textFieldState == TextFieldState.WRONG) {
                toReturn = false;
                break;
            }
        }
        System.out.println("Checked: " + textFieldStates + " returned: " + toReturn);
        return toReturn;
    }

    /**
     * Checks whether the given input is of the valid format (dd-MM-yyyy).
     *
     * @param input a String date
     * @return true if the input is of the format dd-MM-yyyy, false otherwise.
     */
    private boolean isValidDateInput(String input) {
        try {
            validInputFormat.setLenient(false);
            validInputFormat.parse(input);
            return input.length() == 10;
        } catch (ParseException e) {
            return false;
        }
    }
}
\end{lstlisting}
\subsubsection{LoadingDialog.java}
\begin{lstlisting}
package frontend.dialogs;

import javafx.geometry.Insets;
import javafx.geometry.Pos;
import javafx.scene.control.ProgressIndicator;
import javafx.scene.layout.*;
import javafx.scene.text.Text;

/**
 * Class used to show/hide the loading dialog in the current shown screen
 */
public class LoadingDialog {
    private StackPane stackPane;
    private VBox loadingDialog;
    private Pane mainLayout;

    /**
     * Creates a constructor used to show a loading dialog, and remove it.
     *
     * @param stackPane  the Root layout of the window. (Has to be stack pane, as we are adding a loading dialog on top
     *                   of the main layout of the window).
     * @param mainLayout the main content layout, which is disabled while the loading dialog is being shown (to avoid
     *                   the user from pressing buttons while the system is "Loading".
     */
    public LoadingDialog(StackPane stackPane, Pane mainLayout) {
        this.stackPane = stackPane;
        this.mainLayout = mainLayout;//the layout on which we are adding on top the loading dialog
    }

    /**
     * If the loading dialog (the VBox layout) has not been initialized, ie its null, then set it up: creates a layout
     * with text and a progress wheel, indicating the System is loading.
     */
    private void initializeLoadingDialog() {
        if (this.loadingDialog == null) {
            VBox loadingDialog = new VBox();//main layout to hold the dialog
            loadingDialog.getStyleClass().add("loading-dialog");//for the css
            loadingDialog.setPadding(new Insets(20, 15, 20, 15));//add padding
            loadingDialog.getStylesheets().add(getClass().getResource("loadingDialog.css")
            		.toExternalForm());//load the css styles
            //add the "title" of the dialog
            Text pleaseWaitText = new Text("Please Wait... ");
            pleaseWaitText.getStyleClass().add("please-wait-text");
            loadingDialog.getChildren().add(pleaseWaitText);
            loadingDialog.setMargin(pleaseWaitText, new Insets(0, 0, 10, 0));
            //the "body" of the dialog: the progress indicator and text
            GridPane gridPane = new GridPane();
            gridPane.setAlignment(Pos.CENTER);
            gridPane.setVgap(10);
            gridPane.setHgap(28);
            //add progress circle
            ProgressIndicator progressIndicator = new ProgressIndicator();
            progressIndicator.setProgress(-1.0f);
            gridPane.add(progressIndicator, 0, 0, 1, 1);
            //add the text
            Text processingText = new Text("Processing Files... ");
            processingText.getStyleClass().add("processing-text");
            gridPane.add(processingText, 1, 0, 2, 1);//want the text to be wider than the progress indicator
            //add the body
            loadingDialog.getChildren().add(gridPane);
            StackPane.setAlignment(loadingDialog, Pos.CENTER);//to the center of the given stack pane
            loadingDialog.setPrefSize(VBox.USE_COMPUTED_SIZE, VBox.USE_COMPUTED_SIZE);//so that the dialog is never squashed
            loadingDialog.setMaxWidth(VBox.USE_PREF_SIZE);//and never stretched, but instead fit to the size of its content
            loadingDialog.setMaxHeight(VBox.USE_PREF_SIZE);
            this.loadingDialog = loadingDialog;
        }
    }

    /**
     * Will initialize the loading layout if it hasn't been. If the provided StackPane and main layout Pane aren't null,
     * the loading dialog will be shown on top of the layout in the Stack Pane, and the main layout will be disabled
     * until the System finishes "loading".
     */
    public void showLoadingDialog() {
        initializeLoadingDialog();
        if (stackPane != null && mainLayout != null) {
            mainLayout.setDisable(true);
            stackPane.getChildren().add(loadingDialog);
        }
    }

    /**
     * Used to indicate the System has finished "loading". If the Pane's aren't null then the main layout Pane will be
     * enabled again, and the loading dialog will be removed from on top of the Stack Pane.
     * If the loadingDialog has not been shown, then nothing happens (the main layout is enabled - but it should already
     * be enabled).
     */
    public void removeLoadingDialog() {
        if (stackPane != null && mainLayout != null) {
            mainLayout.setDisable(false);
            if (loadingDialog != null) {
                stackPane.getChildren().remove(loadingDialog);
            }
        }
    }
}
\end{lstlisting}
\subsubsection{RemoveConfirmationDialog.java}
\begin{lstlisting}
package frontend.dialogs;

import javafx.scene.control.Alert;
import javafx.scene.control.ButtonType;

/**
 * Used to produce a Confirmation Dialog to ask whether the events for a specified File should be removed from the
 * Timeline.
 */
public class RemoveConfirmationDialog {
    /**
     * Produces a Dialog to allow the User to confirm whether or not the events related to the selected File should be
     * removed or not.
     *
     * @param fileName the name of the File which links to the events to be removed.
     * @return a Confirmation Alert Dialog with options to remove or not the events linked to the given File name.
     */
    public static Alert getRemoveConfirmationDialog(String fileName) {
        Alert removeConfirmationDialog = new Alert(Alert.AlertType.CONFIRMATION);
        removeConfirmationDialog.setTitle("Removing File");
        removeConfirmationDialog.setContentText("Are you sure you want to remove the events for " + fileName + "?");
        removeConfirmationDialog.getButtonTypes().setAll(ButtonType.YES, ButtonType.NO, ButtonType.CANCEL);
        return removeConfirmationDialog;
    }
}
\end{lstlisting}
\subsubsection{SettingsDialog.java}
\begin{lstlisting}
package frontend.dialogs;

import backend.system.BackEndSystem;
import backend.system.Settings;
import javafx.event.EventHandler;
import javafx.geometry.Insets;
import javafx.scene.Node;
import javafx.scene.control.*;
import javafx.scene.input.MouseEvent;
import javafx.scene.layout.GridPane;
import javafx.scene.text.Text;
import javafx.util.Callback;
import javafx.util.StringConverter;

import java.util.Optional;

/**
 * Class to create the Settings Dialog.
 */
public class SettingsDialog {
    private Spinner<Integer> threadCountSpinner;
    private Spinner<Integer> thresholdSpinner;
    private Spinner<Integer> widthSpinner;
    private Spinner<Integer> heightSpinner;

    /**
     * Get a Settings Dialog, that will allow the user to change the Settings of the System and either Save them or
     * discard them. The option to return to the Default Settings is also added. The changed Settings would be applied
     * to the cloned Settings held in the BackEndSystem, which should be then set to the Settings of the System.
     *
     * @return the Settings Dialog.
     * @throws CloneNotSupportedException when cloning the BackEndSystem Settings.
     */
    public Dialog<Settings> settingsDialog() throws CloneNotSupportedException {
        Settings copy = (Settings) BackEndSystem.getInstance().getSettings().clone();
        Dialog<Settings> settingsDialog = new Dialog<>();
        //set up dialog layout
        settingsDialog.setTitle("Preferences");
        settingsDialog.setHeaderText(null);
        settingsDialog.getDialogPane().setContent(getDialogLayout(copy));

        //buttons
        ButtonType buttonTypeSave = new ButtonType("Save", ButtonBar.ButtonData.OK_DONE);
        ButtonType buttonTypeCancel = new ButtonType("Cancel", ButtonBar.ButtonData.CANCEL_CLOSE);
        ButtonType buttonTypeDefault = new ButtonType("Default", ButtonBar.ButtonData.LEFT);
        settingsDialog.getDialogPane().getButtonTypes().setAll(buttonTypeDefault, buttonTypeSave, buttonTypeCancel);

        //set default onaction
        Node defaultButton = settingsDialog.getDialogPane().lookupButton(buttonTypeDefault);
        defaultButton.setOnMousePressed(new EventHandler<MouseEvent>() {
            @Override
            public void handle(MouseEvent event) {
                Alert confirmation = new Alert(Alert.AlertType.CONFIRMATION);
                confirmation.setTitle("Reset to Default?");
                confirmation.setHeaderText(null);
                confirmation.setContentText("Are you sure you wish reset the Settings to Default? ");
                confirmation.getDialogPane().getButtonTypes().setAll(ButtonType.YES, ButtonType.NO);
                Optional<ButtonType> response = confirmation.showAndWait();
                if (response.isPresent() && response.get() == ButtonType.YES) {
                    //reset to default
                    copy.reset();
                    settingsDialog.getDialogPane().setContent(getDialogLayout(copy));
                }
            }
        });

        //set result
        settingsDialog.setResultConverter(new Callback<ButtonType, Settings>() {
            @Override
            public Settings call(ButtonType param) {
                if (param == buttonTypeSave) {
                    copy.setWidth(widthSpinner.getValue());
                    copy.setHeight(heightSpinner.getValue());
                    copy.setThresholdSummary(thresholdSpinner.getValue());
                    copy.setMaxNoOfThreads(threadCountSpinner.getValue());
                    return copy;
                }
                return null;
            }
        });
        return settingsDialog;
    }

    /**
     * For the Settings Dialog, get its Content Root Layout: a GridPane with Text and Spinners for the different Settings
     * options. The given Settings are used to populate the fields in the layout.
     *
     * @param settings the given Settings.
     * @return the root Content layout of the Dialog.
     */
    private GridPane getDialogLayout(Settings settings) {
        GridPane gridPane = new GridPane();
        gridPane.setMaxWidth(Double.MAX_VALUE);
        gridPane.setPadding(new Insets(10));
        gridPane.setHgap(10);
        gridPane.setVgap(10);

        Text fileTitle = new Text("File Processing");
        gridPane.add(fileTitle, 0, 0);

        Text threadText = new Text("Maximum Number of Threads running in parallel: ");
        threadCountSpinner = new Spinner<>();
        SpinnerValueFactory.IntegerSpinnerValueFactory spinnerValueFactory = new SpinnerValueFactory.IntegerSpinnerValueFactory(1, 20, settings.getMaxNoOfThreads());
        spinnerValueFactory.setAmountToStepBy(1);
        spinnerValueFactory.setConverter(new IntegerStringConverter(threadCountSpinner));
        threadCountSpinner.setValueFactory(spinnerValueFactory);
        threadCountSpinner.setEditable(true);
        gridPane.add(threadText, 0, 1);
        gridPane.add(threadCountSpinner, 1, 1);

        Text thresholdText = new Text("Threshold of Text Summary: ");
        thresholdSpinner = new Spinner<>();
        spinnerValueFactory = new SpinnerValueFactory.IntegerSpinnerValueFactory(1, 20, settings.getThresholdSummary());
        spinnerValueFactory.setAmountToStepBy(1);
        spinnerValueFactory.setConverter(new IntegerStringConverter(thresholdSpinner));
        thresholdSpinner.setValueFactory(spinnerValueFactory);
        thresholdSpinner.setEditable(true);
        gridPane.add(thresholdText, 0, 2);
        gridPane.add(thresholdSpinner, 1, 2);

        Separator separator = new Separator();//by default its horizontal
        gridPane.add(separator, 0, 3, 2, 1);

        Text appearanceText = new Text("Appearance");
        gridPane.add(appearanceText, 0, 4);

        Text widthText = new Text("Width at Startup: ");
        widthSpinner = new Spinner<>();
        spinnerValueFactory = new SpinnerValueFactory.IntegerSpinnerValueFactory(720, 1920, settings.getWidth());
        spinnerValueFactory.setAmountToStepBy(5);
        spinnerValueFactory.setConverter(new IntegerStringConverter(widthSpinner));
        widthSpinner.setValueFactory(spinnerValueFactory);
        widthSpinner.setEditable(true);
        gridPane.add(widthText, 0, 5);
        gridPane.add(widthSpinner, 1, 5);

        Text heightText = new Text("Height at Startup: ");
        heightSpinner = new Spinner<>();
        spinnerValueFactory = new SpinnerValueFactory.IntegerSpinnerValueFactory(600, 1080, settings.getHeight());
        spinnerValueFactory.setAmountToStepBy(5);
        spinnerValueFactory.setConverter(new IntegerStringConverter(heightSpinner));
        heightSpinner.setValueFactory(spinnerValueFactory);
        heightSpinner.setEditable(true);
        gridPane.add(heightText, 0, 6);
        gridPane.add(heightSpinner, 1, 6);

        return gridPane;
    }

    /**
     * Private Class to deal with the conversion of Strings to Ints in the Spinner text fields.
     */
    private static class IntegerStringConverter extends StringConverter<Integer> {
        Spinner<Integer> spinner;

        /**
         * Constructor that takes in the Spinner where the changes will be applied on.
         *
         * @param spinner the given Spinner.
         */
        IntegerStringConverter(Spinner<Integer> spinner) {
            this.spinner = spinner;
        }

        /**
         * For the given Integer held in the Spinner, show its String version.
         *
         * @param object the given Integer.
         * @return the Integer as a String.
         */
        @Override
        public String toString(Integer object) {
            return String.valueOf(object);
        }

        /**
         * For the given String (in the textfield) attempt to turn it into an Integer, if that is not possible, use the
         * previous value held by the Spinner (which would be an Integer).
         *
         * @param string the text in the textfield.
         * @return an Integer made from the text in the textfield, or the previous available Integer.
         */
        @Override
        public Integer fromString(String string) {
            Integer integer;
            try {
                integer = Integer.parseInt(string);
            } catch (Exception e) {
                integer = spinner.getValue();
            }
            return integer;
        }
    }
}
\end{lstlisting}

\subsection{Helpers}
\subsubsection{TextFieldState.java}
\begin{lstlisting}
package frontend.helpers;

/**
 * Enum class that represents the state of a TextField in the FileConfirmationDialog in terms of its input being correct
 * or wrong
 */
public enum TextFieldState {
    CORRECT,
    WRONG
}
\end{lstlisting}
\subsection{Observers}
\subsubsection{DocumentReaderObserver.java}
\begin{lstlisting}
package frontend.observers;

/**
 * Interface to be implemented by the Observer of the DocumentReaderController.
 */
public interface DocumentReaderObserver {
    /**
     * Called to allow the Controller to inform the Observer when they need to close the window in which the layout
     * of the Controller resides.
     */
    void close();
}
\end{lstlisting}
\subsubsection{DocumentsLoadedObserver.java}
\begin{lstlisting}
package frontend.observers;

import backend.process.FileData;

/**
 * Interface for the Observers of the DocumentLoadedRowController.
 */
public interface DocumentsLoadedObserver {
    /**
     * Called when all the Results connected to the given FileData need to be removed. (i.e. removing the File given by
     * FileData from the Timeline).
     * @param fileData the FileData for the File where we want to remove its results from the Timeline.
     */
    void remove(FileData fileData);
}
\end{lstlisting}
\subsubsection{EditEventDialogObserver.java}
\begin{lstlisting}
package frontend.observers;

/**
 * Implemented by the Observer of the EditEventDialog Content. To inform, from the content of the Dialog, to the holder
 * of the Dialog to disable/enable the save button.
 */
public interface EditEventDialogObserver {
    /**
     * Called when the Dialog Save button should be disabled or enabled.
     */
    void disableSave(boolean disableSave);
}
\end{lstlisting}
\subsubsection{MenuBarObserver.java}
\begin{lstlisting}
package frontend.observers;

/**
 * Implemented by Observers of Controllers that contain a Menu bar, to handle the relevant events.
 */
public interface MenuBarObserver {
    /**
     * Called when the about menu item is pressed.
     */
    void showAbout();

    /**
     * Called when the close menu item is pressed.
     */
    void close();

    /**
     * Called when the timeline menu item is pressed.
     */
    void timeline();

    /**
     * Called when the preferences menu item is pressed.
     */
    void preferences();
}
\end{lstlisting}
\subsubsection{StartUpObserver.java}
\begin{lstlisting}
package frontend.observers;

/**
 * Interface for the Observer of the StartUpController. It extends the MenuBarObserver as the StartUp scene has a menu
 * bar.
 */
public interface StartUpObserver extends MenuBarObserver {//menu options included

    /**
     * Called when the load files button is pressed.
     */
    void loadFiles();
}
\end{lstlisting}
\subsubsection{TimelineObserver.java}
\begin{lstlisting}
package frontend.observers;

import backend.process.Result;

import java.util.List;

/**
 * Interface for the Observer of the Timeline scene.
 */
public interface TimelineObserver extends MenuBarObserver {
    /**
     * Called when the "Load Documents" button is pressed.
     */
    void loadDocuments();

    /**
     * Called when the "Save to ..." button is pressed.
     * @param resultList the List of Results to save.
     */
    void saveTo(List<Result> resultList);
}
\end{lstlisting}
\subsubsection{TimelineRowObserver.java}
\begin{lstlisting}
package frontend.observers;

import backend.process.Result;

/**
 * Implemented by the Observer of this TimelineRow (that is the holder of the ListView), so that the cells can communicate
 * with the overall listView.
 */
public interface TimelineRowObserver {
    /**
     * Called by a TimelineRow to inform the Observer (holder of the ListView), that the data it is holding has been
     * updated.
     *
     * @param updatedResult the updated data of this TimelineRow.
     * @param position      the position this cell had in the ListView.
     */
    void update(Result updatedResult, int position);

    /**
     * Called by a TimelineRow to inform the Observer (holder of the ListView), that the data it is holding has been
     * updated.
     *
     * @param previous      the Result object that was previously in this Row.
     * @param updatedResult the updated Result object of the TimelineRow.
     */
    void update(Result previous, Result updatedResult);

    /**
     * Called by a TimelineRow to inform the Observer (holder of the ListView), that this event (in the given position),
     * needs to be deleted from the list.
     *
     * @param position the position of the event to be deleted.
     */
    void delete(int position);

    /**
     * Called by a TimelineRow to inform the Observer (holder of the ListView), that this event (the given Result) needs
     * to be deleted from the list.
     *
     * @param result the given event (Result) to be deleted.
     */
    void delete(Result result);
}
\end{lstlisting}
\subsection{Main.java}
\begin{lstlisting}
package frontend;

import backend.helpers.ToJSON;
import backend.helpers.ToPDF;
import backend.process.FileData;
import backend.process.ProcessFiles;
import backend.process.Result;
import backend.system.BackEndSystem;
import backend.system.Settings;
import frontend.controllers.ListViewController;
import frontend.controllers.StartUpController;
import frontend.dialogs.AboutDialog;
import frontend.dialogs.FileConfirmationDialog;
import frontend.dialogs.SettingsDialog;
import frontend.observers.StartUpObserver;
import frontend.observers.TimelineObserver;
import javafx.application.Application;
import javafx.concurrent.Task;
import javafx.concurrent.WorkerStateEvent;
import javafx.event.EventHandler;
import javafx.fxml.FXMLLoader;
import javafx.scene.Scene;
import javafx.scene.control.Alert;
import javafx.scene.control.ButtonBar;
import javafx.scene.control.ButtonType;
import javafx.scene.control.Dialog;
import javafx.stage.FileChooser;
import javafx.stage.Stage;

import java.io.File;
import java.io.IOException;
import java.io.PrintWriter;
import java.util.ArrayList;
import java.util.List;
import java.util.Optional;
import java.util.function.Consumer;

/**
 * Main class that is used to run the program (i.e. show the UI that uses the backend). It handles loading the different
 * layouts, and being the observer of the layouts controllers.
 */
public class Main extends Application implements StartUpObserver, TimelineObserver {
    private final static String TAG = "MAIN: ";
    private Stage primaryStage;
    private StartUpController startUpController;
    private ListViewController listViewController;

    /**
     * Called to start showing the window of the program (i.e. the please load documents layout).
     *
     * @param primaryStage the root window of the application
     * @throws Exception
     */
    @Override
    public void start(Stage primaryStage) throws Exception {
        //need to start engine
        BackEndSystem.getInstance();//thread waits for this to be done
        Settings settings = BackEndSystem.getInstance().getSettings();
        System.out.println("Called getInstance");
        FXMLLoader fxmlLoader = new FXMLLoader(getClass().getResource("startup.fxml"));
        primaryStage.setScene(new Scene(fxmlLoader.load(), settings.getWidth(), settings.getHeight()));
        primaryStage.setTitle("Automated Timeline Extractor - Oliver Philip Hohn");
        startUpController = fxmlLoader.getController();
        startUpController.setObserver(this);
        primaryStage.show();
        this.primaryStage = primaryStage;
    }

    /**
     * First method that gets called when the program runs.
     *
     * @param args the arguments passed in when the instruction to run the program is given.
     */
    public static void main(String[] args) {
        launch(args);
    }

    /**
     * Produces a list of Files, picked by the user in the File Chooser.
     *
     * @param primaryStage the window where the program is running on.
     * @return a list of Files picked by the user in the File Chooser.
     */
    private List<File> loadFiles(Stage primaryStage) {
        FileChooser fileChooser = new FileChooser();
        fileChooser.setTitle("Choose Document Files");
        fileChooser.getExtensionFilters().addAll(new FileChooser.ExtensionFilter("Text Files", "*.txt", "*.pdf", "*.docx"));
        return fileChooser.showOpenMultipleDialog(primaryStage);
    }

    /**
     * For the given List of Files produce a List of FileData, where each FileData corresponds to its File (i.e. have the
     * same index in the list), and holds its Files name, and path.
     *
     * @param files the given List of Files.
     * @return the output List of FileData.
     */
    private List<FileData> getFileData(List<File> files) {
        ArrayList<FileData> toReturn = new ArrayList<>();
        for (File file : files) {
            toReturn.add(new FileData(file));
        }
        return toReturn;
    }

    /**
     * Called to produce a Task object, that will run a set of operations, when given to a Thread, in parallel. This
     * Task object will, for the given Lists of Files and FileData, produce the List of Results that emerge from
     * processing the text in the Files, and linking each Result object to its corresponding FileData.
     *
     * @param files     the given List of Files.
     * @param fileDatas the given List of FileData.
     * @return a Task object to run in a Thread in parallel, to process the given Files and produce Result objects.
     */
    private Task<List<Result>> prepareTask(List<File> files, List<FileData> fileDatas) {
        return new Task<List<Result>>() {
            @Override
            protected List<Result> call() throws Exception {
                ProcessFiles processFiles = new ProcessFiles();
                return processFiles.processFiles(files, fileDatas);
            }
        };
    }

    /**
     * For the global Stage, load the listView layout, set its Observer as Main.this, and hold its controller.
     *
     * @return the Controller of the listView layout.
     */
    private ListViewController showListView() {
        FXMLLoader fxmlLoader = new FXMLLoader(getClass().getResource("listView.fxml"));
        try {
            primaryStage.setScene(new Scene(fxmlLoader.load(), primaryStage.getWidth(), primaryStage.getHeight()));
            listViewController = fxmlLoader.getController();
            listViewController.setTimelineObserver(this);
            primaryStage.show();
        } catch (IOException e) {
            e.printStackTrace();
        }
        return listViewController;
    }


    /**
     * Called from the Controller of the initial layout (the "load-button-layout"), to indicate to Main, that it must
     * provide the user the option to allow the User to pick Files, pass those Files to the backend to process, and
     * then update the layout to show the timeline of the events depicted in the text of the Files.
     * The button to load Files is disabled, to stop the user from pressing it again, while the Files are being loaded.
     */
    @Override
    public void loadFiles() {
        List<File> files = loadFiles(primaryStage);
        if (files != null) {
            List<FileData> fileDatas = getFileData(files);
            Alert fileConfirmationDialog = new FileConfirmationDialog().getConfirmationFileDialog(fileDatas);
            Optional<ButtonType> response = fileConfirmationDialog.showAndWait();
            if (response.isPresent() && response.get() == ButtonType.OK) {
                //disable button
                startUpController.setDisableLoadDocumentsButton(true);
                System.out.println(TAG + "Process Files and set them in the Timeline");
                Task<List<Result>> task = prepareTask(files, fileDatas);
                task.setOnSucceeded(new EventHandler<WorkerStateEvent>() {
                    @Override
                    public void handle(WorkerStateEvent event) {
                        List<Result> results = task.getValue();
                        listViewController = showListView();
                        listViewController.setTimelineListView(results, fileDatas);
                        //stop showing the loading dialog as we have the other layout ready to show
                        startUpController.removeLoadingDialog();
                        startUpController = null;
                        primaryStage.show();
                    }
                });
                task.setOnRunning(new EventHandler<WorkerStateEvent>() {
                    @Override
                    public void handle(WorkerStateEvent event) {
                        //show loading dialog
                        startUpController.showLoadingDialog();
                    }
                });
                new Thread(task).start();
            } else {
                System.out.println(TAG + "Dont Process Files and set them in the Timeline");
            }
        }
    }

    /**
     * Called when the About MenuItem is pressed. Called from either the ListViewController or the StartUpController.
     */
    @Override
    public void showAbout() {
        System.out.println(TAG + "show about information");
        new AboutDialog().getAboutDialog().showAndWait();
    }

    /**
     * Called when the Close MenuItem is pressed. Should close the programs (root) window.
     * Called from either the ListViewController or the StartUpController.
     */
    @Override
    public void close() {
        System.out.println(TAG + "close program");
        primaryStage.close();
    }

    /**
     * Called when the Timeline MenuItem is pressed. Called from either the ListViewController or the StartUpController.
     * But it is only available from the ListView layout, as in the StartUp layout it is disabled.
     */
    @Override
    public void timeline() {
        System.out.println(TAG + "timeline options");
    }

    /**
     * Called when the Preferences MenuItem is pressed in the File Menu.
     */
    @Override
    public void preferences() {
        System.out.println(TAG + "Preferences pressed");
        try {
            Dialog<Settings> settingsDialog = new SettingsDialog().settingsDialog();//show a settings dialog
            Optional<Settings> response = settingsDialog.showAndWait();
            response.ifPresent(new Consumer<Settings>() {
                @Override
                public void accept(Settings settings) {
                    if (settings != null) {//then the user decided to save the settings
                        BackEndSystem.getInstance().setSettings(settings);
                    }//else dont apply the Settings to the System.
                }
            });
        } catch (CloneNotSupportedException e) {
            e.printStackTrace();
        }
    }

    /**
     * Called by the StartUpController, to indicate to Main, to allow the User to pick Files (from the File Chooser),
     * process them, and add them to the timeline (instead of setting them like in the previous method). The layout is
     * not changed, as we are already in the timeline layout.
     */
    @Override
    public void loadDocuments() {
        List<File> files = loadFiles(primaryStage);
        if (files != null) {
            List<FileData> fileDatas = getFileData(files);
            Alert fileConfirmationDialog = new FileConfirmationDialog().getConfirmationFileDialog(fileDatas);
            Optional<ButtonType> response = fileConfirmationDialog.showAndWait();
            if (response.isPresent() && response.get() == ButtonType.OK) {
                System.out.println(TAG + "Process Files and Add them to the Timeline");
                Task<List<Result>> task = prepareTask(files, fileDatas);
                task.setOnSucceeded(new EventHandler<WorkerStateEvent>() {
                    @Override
                    public void handle(WorkerStateEvent event) {
                        List<Result> results = task.getValue();
                        if (listViewController != null) {
                            listViewController.addToTimelineListView(results, fileDatas);
                            //stop showing loading dialog as we have added the results
                            listViewController.removeLoadingDialog();
                        }
                    }
                });
                task.setOnRunning(new EventHandler<WorkerStateEvent>() {
                    @Override
                    public void handle(WorkerStateEvent event) {
                        //show loading dialog
                        listViewController.showLoadingDialog();
                    }
                });
                new Thread(task).start();
            } else {
                System.out.println(TAG + "Don't process Files");
            }
        }
    }


    /**
     * Used to show the Alert the dialog to allow the User to pick in what format to save the List of Results (JSON or
     * PDF). Depending on the option selected, the FileChooser is shown, and then the User picks the location to save
     * the file. An Alert is shown if the User is trying to overwrite a File in use when saving.
     *
     * @param results the List of Results to Save.
     */
    @Override
    public void saveTo(List<Result> results) {
        //show alert to know what to save as
        Alert alert = new Alert(Alert.AlertType.CONFIRMATION);
        alert.setTitle("Save to...");
        alert.setHeaderText(null);
        alert.setContentText("Do you wish to save as PDF or JSON?");
        ButtonType buttonTypePDF = new ButtonType("Save to PDF");
        ButtonType buttonTypeJSON = new ButtonType("Save to JSON");
        ButtonType buttonTypeCancel = new ButtonType("Cancel", ButtonBar.ButtonData.CANCEL_CLOSE);
        alert.getDialogPane().getButtonTypes().setAll(buttonTypePDF, buttonTypeJSON, buttonTypeCancel);

        Optional<ButtonType> response = alert.showAndWait();
        if (response.isPresent() && response.get() == buttonTypePDF) {
            saveToPDF(results);
        } else if (response.isPresent() && response.get() == buttonTypeJSON) {
            saveToJSON(results);
        }
    }

    /**
     * Called when the Timeline being displayed (described by the List of Result objects) needs to be saved as a JSON.
     * The System should allow the user to pick a location to save the PDF, if it is overwriting a File in use by
     * another process the user is informed and given the option to select another location or close the process using
     * the File that they wish to overwrite.
     *
     * @param results the given List of Result objects.
     */
    private void saveToJSON(List<Result> results) {
        String json = ToJSON.toJSON(results);
        FileChooser fileChooser = new FileChooser();
        fileChooser.setTitle("Save Timeline As...");
        fileChooser.getExtensionFilters().addAll(new FileChooser.ExtensionFilter("JSON File", "*.json"));
        File file = fileChooser.showSaveDialog(primaryStage);
        if (file != null && listViewController != null) {
            try {
                PrintWriter printWriter = new PrintWriter(file);
                printWriter.write(json);
                printWriter.close();
            } catch (Exception e) {
                Alert alert = new Alert(Alert.AlertType.INFORMATION);//then inform the user,
                alert.setTitle("File In Use");
                alert.setHeaderText(null);
                alert.setContentText("The file: " + file.getName() + " is in use by another process.");
                alert.getDialogPane().getButtonTypes().setAll(ButtonType.OK);
                Optional<ButtonType> response = alert.showAndWait();
                if (response.isPresent() && response.get() == ButtonType.OK) {//and if they press OK, ie want to save
                    saveToJSON(results);//show them the file chooser to let them pick a different (or same location, if
                }
            }
        }
    }

    /**
     * Called when the Timeline being displayed (described by the List of Result objects) needs to be saved as a PDF.
     * The System should allow the user to pick a location to save the PDF, if it is overwriting a File in use by
     * another process the user is informed and given the option to select another location or close the process using
     * the File that they wish to overwrite.
     *
     * @param results the given List of Result objects.
     */
    private void saveToPDF(List<Result> results) {
        System.out.println(TAG + "Save To PDF pressed");
        FileChooser fileChooser = new FileChooser();
        fileChooser.setTitle("Save Timeline As...");
        fileChooser.getExtensionFilters().addAll(new FileChooser.ExtensionFilter("PDF File", "*.pdf"));
        File file = fileChooser.showSaveDialog(primaryStage);
        if (file != null && listViewController != null) {
            try {
                new ToPDF().saveToPDF(results, file);
            } catch (IOException e) {//if cant save the file, because it is most probably in use or it has been deleted
                Alert alert = new Alert(Alert.AlertType.INFORMATION);//then inform the user,
                alert.setTitle("File In Use");
                alert.setHeaderText(null);
                alert.setContentText("The file: " + file.getName() + " is in use by another process.");
                alert.getDialogPane().getButtonTypes().setAll(ButtonType.OK);
                Optional<ButtonType> response = alert.showAndWait();
                if (response.isPresent() && response.get() == ButtonType.OK) {//and if they press OK, ie want to save
                    saveToPDF(results);//show them the file chooser to let them pick a different (or same location, if
                }                               //they closed the process
            }
        }
    }
}
\end{lstlisting}

\section{src/main/resources/frontend}
\subsection{controllers}
\subsubsection{customErrorFields.css}
\begin{lstlisting}
.text-field:error {
  -fx-text-box-border: red ;
  -fx-focus-color: red ;
}

.text-area:error{
  -fx-text-box-border: red ;
  -fx-focus-color: red ;
}
\end{lstlisting}
\subsubsection{customResultRow.fxml}
\begin{lstlisting}
<?xml version="1.0" encoding="UTF-8"?>

<?import javafx.geometry.Insets?>
<?import javafx.scene.Group?>
<?import javafx.scene.control.Button?>
<?import javafx.scene.control.Label?>
<?import javafx.scene.control.Separator?>
<?import javafx.scene.layout.ColumnConstraints?>
<?import javafx.scene.layout.GridPane?>
<?import javafx.scene.layout.HBox?>
<?import javafx.scene.layout.RowConstraints?>
<?import javafx.scene.layout.VBox?>

<VBox fx:id="rootVBox" maxHeight="1.7976931348623157E308" maxWidth="1.7976931348623157E308" minHeight="-Infinity" minWidth="-Infinity" xmlns="http://javafx.com/javafx/8.0.111" xmlns:fx="http://javafx.com/fxml/1">
   <children>
      <Group>
         <children>
            <GridPane vgap="5.0">
              <columnConstraints>
                <ColumnConstraints hgrow="NEVER" minWidth="10.0" />
                <ColumnConstraints hgrow="NEVER" minWidth="10.0" />
              </columnConstraints>
              <rowConstraints>
                <RowConstraints minHeight="10.0" vgrow="SOMETIMES" />
                <RowConstraints minHeight="10.0" vgrow="ALWAYS" />
                  <RowConstraints prefHeight="30.0" vgrow="ALWAYS" />
              </rowConstraints>
               <children>
                  <Label text="Subjects: " GridPane.halignment="LEFT" GridPane.valignment="TOP" />
                  <Label text="Event: " GridPane.halignment="LEFT" GridPane.rowIndex="1" GridPane.valignment="TOP" />
                  <Label fx:id="subjectsLabel" maxHeight="-Infinity" maxWidth="-Infinity" minHeight="-Infinity" minWidth="-Infinity" prefWidth="230.0" text="[subjects]" wrapText="true" GridPane.columnIndex="1" GridPane.halignment="LEFT" GridPane.valignment="TOP">
                     <GridPane.margin>
                        <Insets right="5.0" />
                     </GridPane.margin>
                  </Label>
                  <Label fx:id="eventLabel" maxHeight="-Infinity" maxWidth="-Infinity" minHeight="-Infinity" minWidth="-Infinity" prefWidth="230.0" text="[event]" wrapText="true" GridPane.columnIndex="1" GridPane.halignment="LEFT" GridPane.rowIndex="1" GridPane.valignment="TOP" GridPane.vgrow="ALWAYS">
                     <GridPane.margin>
                        <Insets right="5.0" />
                     </GridPane.margin>
                  </Label>
                  <Label text="From: " GridPane.halignment="LEFT" GridPane.rowIndex="2" GridPane.valignment="TOP" />
                  <Label fx:id="fromLabel" maxHeight="-Infinity" maxWidth="-Infinity" minHeight="-Infinity" minWidth="-Infinity" prefWidth="230.0" text="[from]" wrapText="true" GridPane.columnIndex="1" GridPane.halignment="LEFT" GridPane.rowIndex="2" GridPane.valignment="TOP" GridPane.vgrow="NEVER" />
               </children>
            </GridPane>
         </children>
      </Group>
      <HBox alignment="TOP_RIGHT">
         <children>
            <Button fx:id="editButton" mnemonicParsing="false" text="Edit">
               <HBox.margin>
                  <Insets bottom="5.0" left="5.0" right="5.0" top="5.0" />
               </HBox.margin>
            </Button>
            <Button fx:id="viewButton" mnemonicParsing="false" text="View">
               <HBox.margin>
                  <Insets bottom="5.0" left="5.0" right="5.0" top="5.0" />
               </HBox.margin>
            </Button>
         </children>
      </HBox>
      <Separator prefWidth="200.0" />
   </children>
   <padding>
      <Insets bottom="5.0" left="5.0" right="5.0" top="5.0" />
   </padding>
</VBox>
\end{lstlisting}
\subsubsection{documentLoadedRow.fxml}
\begin{lstlisting}
<?xml version="1.0" encoding="UTF-8"?>

<?import javafx.geometry.Insets?>
<?import javafx.scene.control.Label?>
<?import javafx.scene.image.Image?>
<?import javafx.scene.image.ImageView?>
<?import javafx.scene.layout.ColumnConstraints?>
<?import javafx.scene.layout.GridPane?>
<?import javafx.scene.layout.HBox?>
<?import javafx.scene.layout.RowConstraints?>
<?import javafx.scene.text.Font?>

<GridPane fx:id="gridPane" maxHeight="-Infinity" maxWidth="-Infinity" minHeight="-Infinity" minWidth="-Infinity" xmlns="http://javafx.com/javafx/8.0.111" xmlns:fx="http://javafx.com/fxml/1">
  <columnConstraints>
    <ColumnConstraints hgrow="SOMETIMES" maxWidth="211.0" minWidth="10.0" prefWidth="211.0" />
    <ColumnConstraints hgrow="SOMETIMES" maxWidth="94.0" minWidth="10.0" prefWidth="39.0" />
  </columnConstraints>
  <rowConstraints>
    <RowConstraints minHeight="10.0" prefHeight="30.0" vgrow="SOMETIMES" />
  </rowConstraints>
   <children>
      <ImageView fx:id="removingImageView" pickOnBounds="true" preserveRatio="true" GridPane.columnIndex="1" GridPane.halignment="RIGHT" GridPane.valignment="CENTER">
         <image>
            <Image url="@iccloseblack.png" />
         </image>
      </ImageView>
      <HBox alignment="CENTER_LEFT" prefHeight="100.0" prefWidth="200.0">
         <children>
            <Label text="·">
               <font>
                  <Font size="36.0" />
               </font>
            </Label>
            <Label fx:id="documentLabel" text="Document1.pdf">
               <padding>
                  <Insets left="5.0" />
               </padding>
            </Label>
         </children>
      </HBox>
   </children>
   <padding>
      <Insets bottom="10.0" left="10.0" right="10.0" top="10.0" />
   </padding>
</GridPane>
\end{lstlisting}
\subsubsection{documentReader.fxml}
\begin{lstlisting}
<?xml version="1.0" encoding="UTF-8"?>

<?import javafx.geometry.Insets?>
<?import javafx.scene.control.*?>
<?import javafx.scene.layout.BorderPane?>
<?import javafx.scene.text.Font?>
<?import org.fxmisc.richtext.InlineCssTextArea?>
<BorderPane fx:id="rootBorderPane" maxHeight="-Infinity" maxWidth="-Infinity" minHeight="-Infinity" minWidth="-Infinity" prefHeight="400.0" prefWidth="600.0" xmlns="http://javafx.com/javafx/8.0.111" xmlns:fx="http://javafx.com/fxml/1">
   <center>
         <InlineCssTextArea fx:id="documentInlineCssTextArea" editable="false" wrapText="true" BorderPane.alignment="TOP_LEFT">
            <BorderPane.margin>
               <Insets bottom="5.0" left="5.0" right="5.0" top="5.0" />
            </BorderPane.margin>
            <padding>
               <Insets bottom="2.0" left="2.0" right="2.0" top="2.0" />
            </padding>
         <font>
            <Font size="18.0" />
         </font>
         </InlineCssTextArea>
   </center>
   <top>
      <MenuBar BorderPane.alignment="CENTER">
        <menus>
          <Menu mnemonicParsing="false" text="File">
            <items>
              <MenuItem fx:id="closeMenuItem" mnemonicParsing="false" text="Close" />
            </items>
          </Menu>
          <Menu mnemonicParsing="false" text="Edit">
            <items>
              <MenuItem fx:id="copyMenuItem" mnemonicParsing="false" text="Copy" />
            </items>
          </Menu>
        </menus>
      </MenuBar>
   </top>
</BorderPane>
\end{lstlisting}
\subsubsection{editEventDialog.fxml}
\begin{lstlisting}
<?xml version="1.0" encoding="UTF-8"?>

<?import javafx.geometry.Insets?>
<?import javafx.scene.control.Label?>
<?import javafx.scene.control.TextArea?>
<?import javafx.scene.control.TextField?>
<?import javafx.scene.image.Image?>
<?import javafx.scene.image.ImageView?>
<?import javafx.scene.layout.BorderPane?>
<?import javafx.scene.layout.ColumnConstraints?>
<?import javafx.scene.layout.GridPane?>
<?import javafx.scene.layout.HBox?>
<?import javafx.scene.layout.RowConstraints?>
<?import javafx.scene.layout.VBox?>

<GridPane fx:id="rootGridPane" alignment="CENTER" hgap="10.0" maxHeight="-Infinity" maxWidth="-Infinity" minHeight="-Infinity" minWidth="-Infinity" prefHeight="400.0" prefWidth="600.0" vgap="10.0" xmlns="http://javafx.com/javafx/8.0.111" xmlns:fx="http://javafx.com/fxml/1">
  <columnConstraints>
    <ColumnConstraints hgrow="SOMETIMES" minWidth="10.0" prefWidth="100.0" />
    <ColumnConstraints hgrow="SOMETIMES" minWidth="10.0" prefWidth="100.0" />
  </columnConstraints>
  <rowConstraints>
    <RowConstraints minHeight="10.0" prefHeight="30.0" vgrow="SOMETIMES" />
    <RowConstraints maxHeight="46.0" minHeight="10.0" prefHeight="46.0" vgrow="SOMETIMES" />
    <RowConstraints maxHeight="305.0" minHeight="10.0" prefHeight="294.0" vgrow="SOMETIMES" />
  </rowConstraints>
   <children>
      <Label text="Date: " GridPane.hgrow="NEVER" GridPane.vgrow="NEVER" />
      <Label text="Subjects: " GridPane.hgrow="NEVER" GridPane.rowIndex="1" GridPane.valignment="CENTER" GridPane.vgrow="NEVER" />
      <Label text="Event: " GridPane.hgrow="NEVER" GridPane.rowIndex="2" GridPane.valignment="TOP" GridPane.vgrow="NEVER">
         <padding>
            <Insets bottom="5.0" right="5.0" top="5.0" />
         </padding>
      </Label>
      <BorderPane GridPane.columnIndex="1" GridPane.halignment="CENTER" GridPane.hgrow="ALWAYS" GridPane.valignment="CENTER" GridPane.vgrow="NEVER">
         <left>
            <TextField fx:id="firstDateTextField" prefHeight="31.0" prefWidth="214.0" promptText="dd-MM-yyyy" text="dd-MM-yyyy" BorderPane.alignment="CENTER" />
         </left>
         <center>
            <Label alignment="CENTER" text="-&gt;" BorderPane.alignment="CENTER" />
         </center>
         <right>
            <TextField fx:id="secondDateTextField" prefWidth="214.0" promptText="dd-MM-yyyy" BorderPane.alignment="CENTER" />
         </right>
         <padding>
            <Insets bottom="10.0" left="10.0" right="10.0" top="10.0" />
         </padding>
      </BorderPane>
      <VBox spacing="5.0" GridPane.columnIndex="1" GridPane.rowIndex="1">
         <children>
            <BorderPane prefHeight="200.0" prefWidth="200.0" VBox.vgrow="ALWAYS">
               <center>
                  <TextField fx:id="subjectTextField" BorderPane.alignment="CENTER">
                     <BorderPane.margin>
                        <Insets right="5.0" />
                     </BorderPane.margin>
                  </TextField>
               </center>
               <right>
                  <ImageView fx:id="addImageView" pickOnBounds="true" preserveRatio="true" BorderPane.alignment="CENTER">
                     <image>
                        <Image url="@icaddcircleoutline.png" />
                     </image>
                  </ImageView>
               </right>
               <padding>
                  <Insets top="10.0" />
               </padding>
            </BorderPane>
            <HBox fx:id="subjectsHBox" prefHeight="100.0" prefWidth="200.0" spacing="5.0" />
         </children>
         <padding>
            <Insets bottom="10.0" left="10.0" right="10.0" top="10.0" />
         </padding>
      </VBox>
      <VBox alignment="TOP_RIGHT" prefHeight="200.0" prefWidth="100.0" spacing="2.0" GridPane.columnIndex="1" GridPane.rowIndex="2">
         <children>
            <TextArea fx:id="eventTextArea" prefHeight="200.0" prefWidth="200.0" wrapText="true" VBox.vgrow="ALWAYS" />
            <Label fx:id="maxCharacterLabel" text="Maximum 100 characters" textFill="#797979" />
         </children>
         <padding>
            <Insets bottom="10.0" left="10.0" right="10.0" top="10.0" />
         </padding>
      </VBox>
   </children>
   <padding>
      <Insets bottom="5.0" left="5.0" right="5.0" top="5.0" />
   </padding>
</GridPane>
\end{lstlisting}
\subsubsection{icaddcircleoutline.png}
\begin{figure}
\caption{PNG Image used to indicate the adding of subjects to an event.}
\includegraphics{icaddcircleoutline.png}
\centering
\end{figure}
\subsubsection{iccloseblack.png}
\begin{figure}
\caption{PNG Image used for the action of deleting documents from the timeline.}
\includegraphics{iccloseblack.png}
\centering
\end{figure}
\subsubsection{listViewTheme.css}
\begin{lstlisting}
.list-cell:even {
    -fx-background-color: white;
    -fx-text-fill: black;
    -fx-padding:0px 0px 0px 10px;
}

.list-cell:odd{
    -fx-background-color: white;
    -fx-text-fill: black;
    -fx-padding:0px 10px 0px 0px;
}

.list-cell:filled:selected:focused, .list-cell:filled:selected {
    -fx-focus-color: transparent;
}

.label{
    -fx-text-fill:black;
}
\end{lstlisting}
\subsubsection{listViewThemeTimeline.css}
\begin{lstlisting}
.list-cell {
    -fx-background-color: white;
    -fx-text-fill: black;
}

.list-cell:filled:selected:focused, .list-cell:filled:selected {
    -fx-focus-color: transparent;
}

.label{
    -fx-text-fill:black;
}
\end{lstlisting}
\subsubsection{rangeDataLayout.fxml}
\begin{lstlisting}
<?xml version="1.0" encoding="UTF-8"?>

<?import javafx.geometry.Insets?>
<?import javafx.scene.control.Label?>
<?import javafx.scene.control.ListView?>
<?import javafx.scene.layout.VBox?>
<?import javafx.scene.text.Font?>

<VBox fx:id="rootVBox" maxHeight="-Infinity" maxWidth="-Infinity" xmlns="http://javafx.com/javafx/8.0.111" xmlns:fx="http://javafx.com/fxml/1">
   <children>
      <Label fx:id="dateLabel" text="DATE">
         <font>
            <Font name="System Bold" size="19.0" />
         </font>
      </Label>
      <ListView fx:id="resultsListView" maxHeight="-Infinity" maxWidth="-Infinity" prefHeight="160.0" prefWidth="350.0" stylesheets="@listViewThemeTimeline.css" VBox.vgrow="ALWAYS">
         <VBox.margin>
            <Insets bottom="5.0" left="15.0" right="5.0" top="5.0" />
         </VBox.margin></ListView>
   </children>
   <padding>
      <Insets bottom="5.0" left="5.0" right="5.0" top="5.0" />
   </padding>
</VBox>
\end{lstlisting}
\subsubsection{timelineRowEven.fxml}
\begin{lstlisting}
<?xml version="1.0" encoding="UTF-8"?>

<?import javafx.geometry.Insets?>
<?import javafx.scene.Group?>
<?import javafx.scene.control.Button?>
<?import javafx.scene.control.Label?>
<?import javafx.scene.layout.BorderPane?>
<?import javafx.scene.layout.HBox?>
<?import javafx.scene.layout.Region?>
<?import javafx.scene.layout.VBox?>
<?import javafx.scene.shape.Line?>

<Group fx:id="group" xmlns="http://javafx.com/javafx/8.0.111" xmlns:fx="http://javafx.com/fxml/1">
   <children>
      <HBox alignment="CENTER_LEFT" layoutX="10.0" layoutY="10.0" maxHeight="-Infinity" maxWidth="-Infinity" minHeight="-Infinity" minWidth="-Infinity">
         <children>
            <BorderPane fx:id="borderPane" maxHeight="-Infinity" maxWidth="-Infinity" minHeight="-Infinity" minWidth="-Infinity" prefWidth="370.0">
               <padding>
                  <Insets bottom="15.0" left="15.0" right="15.0" top="15.0" />
               </padding>
               <bottom>
                  <BorderPane>
                     <left>
                        <Button fx:id="editButton" mnemonicParsing="false" text="Edit Event" BorderPane.alignment="CENTER" />
                     </left>
                     <right>
                        <Button fx:id="viewButton" mnemonicParsing="false" text="Go to Document" BorderPane.alignment="CENTER" />
                     </right>
                     <padding>
                        <Insets bottom="5.0" top="5.0" />
                     </padding>
                  </BorderPane>
               </bottom>
               <left>
                  <Group BorderPane.alignment="CENTER">
                     <children>
                        <VBox prefWidth="370.0">
                           <children>
                              <Label fx:id="eventNumberLabel" text="Event ##" />
                              <Label fx:id="dateLabel" maxHeight="-Infinity" maxWidth="-Infinity" minHeight="-Infinity" minWidth="-Infinity" prefWidth="340.0" text="Date: dd-MM-yyy -&gt; dd-MM-yyyy" wrapText="true" />
                              <Label fx:id="subjectsLabel" maxHeight="-Infinity" maxWidth="-Infinity" minHeight="-Infinity" minWidth="-Infinity" prefWidth="340.0" text="Subjects: Football, Gifts, London" wrapText="true" />
                              <Label fx:id="eventLabel" maxHeight="-Infinity" maxWidth="-Infinity" minHeight="-Infinity" minWidth="-Infinity" prefWidth="340.0" text="Event: played football with friends" wrapText="true" />
                              <Label fx:id="fromLabel" maxHeight="-Infinity" maxWidth="-Infinity" minHeight="-Infinity" minWidth="-Infinity" prefWidth="340.0" text="From: filename.ext" wrapText="true" />
                           </children>
                           <padding>
                              <Insets bottom="5.0" top="5.0" />
                           </padding>
                        </VBox>
                     </children>
                  </Group>
               </left>
            </BorderPane>
            <Line startX="-100.0" strokeWidth="3.0" />
            <Line endY="250.0" nodeOrientation="LEFT_TO_RIGHT" strokeWidth="3.0" />
            <Region maxHeight="-Infinity" maxWidth="-Infinity" minHeight="-Infinity" minWidth="-Infinity" prefHeight="200.0" prefWidth="483.0" />
         </children>
         <padding>
            <Insets right="10.0" />
         </padding>
      </HBox>
   </children>
</Group>
\end{lstlisting}
\subsubsection{timelineRowOdd.fxml}
\begin{lstlisting}
<?xml version="1.0" encoding="UTF-8"?>

<?import javafx.geometry.Insets?>
<?import javafx.scene.Group?>
<?import javafx.scene.control.Button?>
<?import javafx.scene.control.Label?>
<?import javafx.scene.layout.BorderPane?>
<?import javafx.scene.layout.HBox?>
<?import javafx.scene.layout.Region?>
<?import javafx.scene.layout.VBox?>
<?import javafx.scene.shape.Line?>

<Group fx:id="group" xmlns="http://javafx.com/javafx/8.0.111" xmlns:fx="http://javafx.com/fxml/1">
    <children>
        <HBox alignment="CENTER_LEFT" layoutX="10.0" layoutY="10.0" maxHeight="-Infinity" maxWidth="-Infinity" minHeight="-Infinity" minWidth="-Infinity">
            <children>
            <Region maxHeight="-Infinity" maxWidth="-Infinity" minHeight="-Infinity" minWidth="-Infinity" prefHeight="200.0" prefWidth="483.0" />
                <Line endY="250.0" nodeOrientation="LEFT_TO_RIGHT" strokeWidth="3.0" />
                <Line startX="-100.0" strokeWidth="3.0" />
            <BorderPane fx:id="borderPane" maxHeight="-Infinity" maxWidth="-Infinity" minHeight="-Infinity" minWidth="-Infinity" prefWidth="370.0">
               <padding>
                  <Insets bottom="15.0" left="15.0" right="15.0" top="15.0" />
               </padding>
               <bottom>
                  <BorderPane>
                     <left>
                        <Button fx:id="editButton" mnemonicParsing="false" text="Edit Event" BorderPane.alignment="CENTER" />
                     </left>
                     <right>
                        <Button fx:id="viewButton" mnemonicParsing="false" text="Go to Document" BorderPane.alignment="CENTER" />
                     </right>
                     <padding>
                        <Insets bottom="5.0" top="5.0" />
                     </padding>
                  </BorderPane>
               </bottom>
               <left>
                  <Group BorderPane.alignment="CENTER">
                     <children>
                        <VBox prefWidth="370.0">
                           <children>
                              <Label fx:id="eventNumberLabel" text="Event ##" />
                              <Label fx:id="dateLabel" maxHeight="-Infinity" maxWidth="-Infinity" minHeight="-Infinity" minWidth="-Infinity" prefWidth="340.0" text="Date: dd-MM-yyy -&gt; dd-MM-yyyy" wrapText="true" />
                              <Label fx:id="subjectsLabel" maxHeight="-Infinity" maxWidth="-Infinity" minHeight="-Infinity" minWidth="-Infinity" prefWidth="340.0" text="Subjects: Football, Gifts, London" wrapText="true" />
                              <Label fx:id="eventLabel" maxHeight="-Infinity" maxWidth="-Infinity" minHeight="-Infinity" minWidth="-Infinity" prefWidth="340.0" text="Event: played football with friends" wrapText="true" />
                               <Label fx:id="fromLabel" maxHeight="-Infinity" maxWidth="-Infinity" minHeight="-Infinity" minWidth="-Infinity" prefWidth="340.0" text="From: filename.ext" wrapText="true" />
                           </children>
                           <padding>
                              <Insets bottom="5.0" top="5.0" />
                           </padding>
                        </VBox>
                     </children>
                  </Group>
               </left>
            </BorderPane>
            </children>
            <padding>
                <Insets right="10.0" />
            </padding>
        </HBox>
    </children>
</Group>
\end{lstlisting}

\subsection{dialogs}
\subsubsection{customErrorFields.css}
\begin{lstlisting}
.text-field:error {
  -fx-text-box-border: red ;
  -fx-focus-color: red ;
}

.text-area:error{
  -fx-text-box-border: red ;
  -fx-focus-color: red ;
}
\end{lstlisting}
\subsubsection{loadingDialog.css}
\begin{lstlisting}
.please-wait-text {
    -fx-font-weight: bold;
    -fx-font-size: 18px;
}

.processing-text {
    -fx-font-size: 16px;
}

.loading-dialog {
    -fx-background-color: #ffffff;
    -fx-border-color: -fx-text-box-border;
    -fx-border-width: 2;
}
\end{lstlisting}
\subsection{documentListViewTheme.css}
\begin{lstlisting}
.list-cell {
    -fx-background-color: white;
    -fx-text-fill: black;
}


.list-cell:filled:selected:focused, .list-cell:filled:selected {
    -fx-focus-color: transparent;

}

.label{
    -fx-text-fill:black;
}
\end{lstlisting}
\subsection{listView.fxml}
\begin{lstlisting}
<?xml version="1.0" encoding="UTF-8"?>

<?import javafx.geometry.Insets?>
<?import javafx.scene.control.Button?>
<?import javafx.scene.control.Label?>
<?import javafx.scene.control.ListView?>
<?import javafx.scene.control.Menu?>
<?import javafx.scene.control.MenuBar?>
<?import javafx.scene.control.MenuItem?>
<?import javafx.scene.control.RadioMenuItem?>
<?import javafx.scene.control.ScrollPane?>
<?import javafx.scene.control.SeparatorMenuItem?>
<?import javafx.scene.layout.BorderPane?>
<?import javafx.scene.layout.ColumnConstraints?>
<?import javafx.scene.layout.GridPane?>
<?import javafx.scene.layout.HBox?>
<?import javafx.scene.layout.RowConstraints?>
<?import javafx.scene.layout.StackPane?>
<?import javafx.scene.layout.VBox?>
<?import javafx.scene.text.Font?>

<StackPane fx:id="stackPane" maxHeight="-Infinity" maxWidth="-Infinity" minHeight="-Infinity" minWidth="-Infinity" xmlns="http://javafx.com/javafx/8.0.111" xmlns:fx="http://javafx.com/fxml/1" fx:controller="frontend.controllers.ListViewController">
   <children>
      <VBox fx:id="vBox" maxHeight="1.7976931348623157E308" maxWidth="1.7976931348623157E308" minHeight="-Infinity" minWidth="-Infinity" prefHeight="800.0" prefWidth="1024.0" StackPane.alignment="CENTER">
         <children>
            <MenuBar>
               <menus>
                  <Menu mnemonicParsing="false" text="File">
                     <items>
                        <MenuItem mnemonicParsing="false" onAction="#close" text="Close" />
                        <MenuItem mnemonicParsing="false" onAction="#about" text="About" />
                        <SeparatorMenuItem mnemonicParsing="false" />
                        <MenuItem mnemonicParsing="false" onAction="#preferences" text="Preferences..." />
                     </items>
                  </Menu>
                  <Menu mnemonicParsing="false" text="Timeline">
                     <items>
                        <RadioMenuItem fx:id="dateView" mnemonicParsing="false" text="Date View" />
                        <RadioMenuItem fx:id="rangeView" mnemonicParsing="false" text="Range View" />
                     </items>
                  </Menu>
               </menus>
            </MenuBar>
            <GridPane VBox.vgrow="ALWAYS">
               <columnConstraints>
                  <ColumnConstraints hgrow="SOMETIMES" maxWidth="370.0" minWidth="10.0" />
                  <ColumnConstraints hgrow="SOMETIMES" minWidth="10.0" />
                  <ColumnConstraints hgrow="SOMETIMES" minWidth="10.0" prefWidth="100.0" />
                  <ColumnConstraints hgrow="SOMETIMES" minWidth="10.0" prefWidth="100.0" />
                  <ColumnConstraints hgrow="SOMETIMES" minWidth="10.0" prefWidth="100.0" />
               </columnConstraints>
               <rowConstraints>
                  <RowConstraints minHeight="10.0" prefHeight="30.0" vgrow="SOMETIMES" />
               </rowConstraints>
               <children>
                  <BorderPane GridPane.rowSpan="2147483647">
                     <top>
                        <VBox prefWidth="100.0" BorderPane.alignment="CENTER">
                           <children>
                              <Label text="Documents Loaded:">
                                 <padding>
                                    <Insets bottom="5.0" right="5.0" top="5.0" />
                                 </padding>
                              </Label>
                              <ListView fx:id="documentListView" maxHeight="-Infinity" maxWidth="-Infinity" prefHeight="280.0" prefWidth="310.0" stylesheets="@documentListViewTheme.css" />
                           </children>
                           <padding>
                              <Insets bottom="10.0" left="20.0" right="20.0" top="10.0" />
                           </padding>
                        </VBox>
                     </top>
                     <bottom>
                        <HBox alignment="CENTER" prefHeight="100.0" prefWidth="200.0" BorderPane.alignment="CENTER">
                           <children>
                              <BorderPane prefHeight="200.0" prefWidth="200.0" HBox.hgrow="ALWAYS">
                                 <left>
                                    <Button fx:id="loadDocumentsButton" maxWidth="-Infinity" minWidth="-Infinity" mnemonicParsing="false" prefWidth="120.0" text="Load Documents" textAlignment="CENTER" wrapText="true" BorderPane.alignment="CENTER">
                                       <font>
                                          <Font size="18.0" />
                                       </font>
                                    </Button>
                                 </left>
                                 <right>
                                    <Button fx:id="saveToButton" maxWidth="-Infinity" minWidth="-Infinity" mnemonicParsing="false" prefWidth="120.0" text="Save to... " textAlignment="CENTER" wrapText="true" BorderPane.alignment="CENTER">
                                       <font>
                                          <Font size="18.0" />
                                       </font>
                                    </Button>
                                 </right>
                                 <padding>
                                    <Insets bottom="10.0" left="20.0" right="20.0" top="10.0" />
                                 </padding>
                              </BorderPane>
                           </children>
                        </HBox>
                     </bottom>
                  </BorderPane>
                  <ListView fx:id="timelineListView" GridPane.columnIndex="1" GridPane.columnSpan="2147483647" GridPane.halignment="RIGHT" GridPane.hgrow="ALWAYS" GridPane.rowSpan="2147483647">
                     <padding>
                        <Insets bottom="10.0" left="10.0" right="10.0" top="10.0" />
                     </padding>
                  </ListView>
                  <ScrollPane fx:id="scrollPane" prefHeight="200.0" prefWidth="200.0" GridPane.columnIndex="1" GridPane.columnSpan="2147483647" GridPane.rowSpan="2147483647" GridPane.valignment="TOP" />
               </children>
            </GridPane>
         </children>
      </VBox>
   </children>
</StackPane>
\end{lstlisting}
\subsection{startup.fxml}
\begin{lstlisting}
<?xml version="1.0" encoding="UTF-8"?>

<?import javafx.geometry.Insets?>
<?import javafx.scene.control.Button?>
<?import javafx.scene.control.Menu?>
<?import javafx.scene.control.MenuBar?>
<?import javafx.scene.control.MenuItem?>
<?import javafx.scene.control.SeparatorMenuItem?>
<?import javafx.scene.layout.BorderPane?>
<?import javafx.scene.layout.StackPane?>
<?import javafx.scene.layout.VBox?>

<StackPane fx:id="stackPane" maxHeight="-Infinity" maxWidth="-Infinity" minHeight="-Infinity" minWidth="-Infinity" prefHeight="400.0" prefWidth="600.0" xmlns="http://javafx.com/javafx/8.0.111" xmlns:fx="http://javafx.com/fxml/1" fx:controller="frontend.controllers.StartUpController">
   <children>
      <VBox fx:id="vBox" maxHeight="1.7976931348623157E308" maxWidth="1.7976931348623157E308" StackPane.alignment="CENTER">
         <children>
            <MenuBar>
               <menus>
                  <Menu mnemonicParsing="false" text="File">
                     <items>
                        <MenuItem mnemonicParsing="false" onAction="#close" text="Close" />
                        <MenuItem mnemonicParsing="false" onAction="#about" text="About" />
                        <SeparatorMenuItem mnemonicParsing="false" />
                        <MenuItem mnemonicParsing="false" onAction="#preferences" text="Preferences..." />
                     </items>
                  </Menu>
                  <Menu disable="true" mnemonicParsing="false" text="Timeline">
                     <items>
                        <MenuItem mnemonicParsing="false" onAction="#timeline" text="Delete" />
                     </items>
                  </Menu>
               </menus>
            </MenuBar>
            <BorderPane VBox.vgrow="ALWAYS">
               <center>
                  <Button fx:id="loadDocumentsButton" mnemonicParsing="false" onAction="#loadDocuments" prefHeight="111.0" prefWidth="283.0" text="Get started by loading in documents!" BorderPane.alignment="CENTER" />
               </center>
               <padding>
                  <Insets bottom="10.0" left="10.0" right="10.0" top="10.0" />
               </padding>
               <VBox.margin>
                  <Insets />
               </VBox.margin>
            </BorderPane>
         </children>
      </VBox>
   </children>
</StackPane>
\end{lstlisting}


\section{src/test/backend}
\subsubsection{EngineTest.java}
\begin{lstlisting}
package backend;

import backend.process.Engine;
import backend.process.Result;
import backend.process.TimelineDate;
import org.junit.Assert;
import org.junit.Test;

import java.text.ParseException;
import java.text.SimpleDateFormat;
import java.util.ArrayList;

/**
 * Test class for the backend.process.Engine. Gives in a sample text, and compares that the generated and predicted backend.process.Result objects are equal
 */
public class EngineTest {
    private final static SimpleDateFormat simpleDateFormat = new SimpleDateFormat("yyyy-MM-dd");

    /**
     * Gives backend.process.Engine sample text and compares the expected and generated output.
     * Only compares the output by the date (checking that it is picking out the right sentence, not looking at sentence trimming).
     *
     * @throws ParseException thrown when we create the Dates for the expected TimelineDates and Results.
     */
    @Test
    public void testDates() throws ParseException {
        String sampleText = "On the 12th of December I ran tests on my final year project. The tests did not go so well. " +
                "Yesterday I played games. " + "It was fun playing games! Tomorrow I am going to study. Last week I went to watch a football match.";
        ArrayList<Result> expectedResults = new ArrayList<>();

        Result result1 = new Result();
        TimelineDate timelineDate1 = new TimelineDate();
        timelineDate1.setDate1(simpleDateFormat.parse("2016-12-12"));
        result1.setTimelineDate(timelineDate1);
        expectedResults.add(result1);

        Result result2 = new Result();
        TimelineDate timelineDate2 = new TimelineDate();
        timelineDate2.setDate1(simpleDateFormat.parse("2016-12-22"));
        result2.setTimelineDate(timelineDate2);
        expectedResults.add(result2);

        Result result3 = new Result();
        TimelineDate timelineDate3 = new TimelineDate();
        timelineDate3.setDate1(simpleDateFormat.parse("2016-12-24"));
        result3.setTimelineDate(timelineDate3);
        expectedResults.add(result3);

        Result result4 = new Result();
        TimelineDate timelineDate4 = new TimelineDate();
        timelineDate4.setDate1(simpleDateFormat.parse("2016-12-12"));
        timelineDate4.setDate2(simpleDateFormat.parse("2016-12-18"));
        result4.setTimelineDate(timelineDate4);
        expectedResults.add(result4);

        Engine engine = new Engine();
        ArrayList<Result> results = engine.getResults(sampleText, "2016-12-23");
        compareExpectedToActualDate(results, expectedResults);
    }

    /**
     * Tests the Engines ability to pick out Subjects from a given Sample text, and checks it to the expected Subjects
     * for that text.
     */
    @Test
    public void testSubjects() {
        //picks up last 16 as a date (issue with casual language used)
        String sampleText = "Manchester City face Monaco in the Champions League, with the first leg on 21 February. In their last season, Liverpool narrowly missed out on the 2013/14 Premier League title to Manchester City, while Leicester won last year's title only playing domestically. ";
        ArrayList<Result> expectedResults = new ArrayList<>();//just looking at subjects, not dates
        Result result1 = new Result();//two sentences, both with dates, so two results made
        Result result2 = new Result();//should pick out the team names: Monaco, Liverpool, Manchester City, Leicester, and Competitions: Champions League, Premier League
        //for the first sentence, the teams mentioned
        result1.addSubject("Champions League");
        result1.addSubject("Monaco");
        result1.addSubject("Manchester City");

        result2.addSubject("Liverpool");
        result2.addSubject("Manchester City");
        result2.addSubject("Premier League");
        result2.addSubject("Leicester");

        expectedResults.add(result1);
        expectedResults.add(result2);

        Engine engine = new Engine();
        ArrayList<Result> actualResults = engine.getResults(sampleText, "26-12-2016");

        compareExpectedToActualSubject(actualResults, expectedResults);
    }


    /**
     * Compares an actual list of Results with a expected one, by just looking at the dates picked out. Checking that the backend.process.Engine
     * is picking out the right sentence and producing the right dates, not caring about the sentence trimming or subject picking.
     *
     * @param actualResults   the list of produced backend.process.Result objects
     * @param expectedResults the list of expected backend.process.Result objects
     */
    private void compareExpectedToActualDate(ArrayList<Result> actualResults, ArrayList<Result> expectedResults) {
        System.out.println("Actual Results: " + actualResults);
        Assert.assertEquals(expectedResults.size(), actualResults.size());//if this does not hold then the test fails
        for (int i = 0; i < actualResults.size(); i++) {
            Assert.assertEquals(expectedResults.get(i).getTimelineDate(), actualResults.get(i).getTimelineDate());
        }
    }

    /**
     * Compares the Subjects (only) in the list of Expected backend.process.Result objects to the list of Actual backend.process.Result objects.
     * Compares each actual backend.process.Result object individually to its expected counterpart.
     *
     * @param actualResults   the list of produced backend.process.Result objects
     * @param expectedResults the list of expected backend.process.Result objects
     */
    private void compareExpectedToActualSubject(ArrayList<Result> actualResults, ArrayList<Result> expectedResults) {
        System.out.println("Actual Results: " + actualResults);
        System.out.println("Expected Results: " + expectedResults);
        Assert.assertEquals(expectedResults.size(), actualResults.size());
        for (int i = 0; i < actualResults.size(); i++) {
            Assert.assertEquals(expectedResults.get(i).getSubjects(), actualResults.get(i).getSubjects());
        }
    }
}
\end{lstlisting}
\subsubsection{ProcessFileTest.java}
\begin{lstlisting}
package backend;

import backend.process.*;
import javafx.concurrent.Task;
import org.junit.Assert;
import org.junit.Test;

import java.io.File;
import java.net.URISyntaxException;
import java.text.ParseException;
import java.text.SimpleDateFormat;
import java.util.ArrayList;
import java.util.List;

/**
 * Tests the the extracting of text in files and passing it to the backend.process.Engine, and the fact that Threads are being created
 * to process the Files separately.
 */

public class ProcessFileTest {
    private SimpleDateFormat simpleDateFormat = new SimpleDateFormat("yyyy-MM-dd");
    private List<Result> actualResults;

    /**
     * Chaining tests (all in same class running one after the other) for ProcessFile will overlap, as when the next
     * test runs, the other has not finished.
     * <p>
     * Test for processing a test file with sample text, and comparing its output to expected backend.process.Result objects.
     *
     * @throws ParseException       for the Dates that we create for the expected Results.
     * @throws InterruptedException as we are putting the Thread that runs this test to sleep to let ProcessFile finish running.
     */
    @Test
    public void testSampleFileProcessTXT() throws ParseException, InterruptedException, URISyntaxException {
        List<Result> expectedResults = new ArrayList<>();
        actualResults = null;

        Result result1 = new Result();
        TimelineDate timelineDate1 = new TimelineDate();
        timelineDate1.setDate1(simpleDateFormat.parse("2016-12-12"));
        result1.setTimelineDate(timelineDate1);
        expectedResults.add(result1);

        Result result2 = new Result();
        TimelineDate timelineDate2 = new TimelineDate();
        timelineDate2.setDate1(simpleDateFormat.parse("2017-01-27"));
        result2.setTimelineDate(timelineDate2);
        expectedResults.add(result2);

        Result result3 = new Result();
        TimelineDate timelineDate3 = new TimelineDate();
        timelineDate3.setDate1(simpleDateFormat.parse("2017-01-29"));
        result3.setTimelineDate(timelineDate3);
        expectedResults.add(result3);

        Result result4 = new Result();
        TimelineDate timelineDate4 = new TimelineDate();
        timelineDate4.setDate1(simpleDateFormat.parse("2017-01-16"));
        timelineDate4.setDate2(simpleDateFormat.parse("2017-01-22"));
        result4.setTimelineDate(timelineDate4);
        expectedResults.add(result4);

        File testFile = new File(getClass().getResource("testfile1.txt").toURI());
        ArrayList<File> files = new ArrayList<>();
        files.add(testFile);

        ArrayList<FileData> fileDatas = new ArrayList<>();
        FileData fileData = new FileData("testfile1.txt", getClass().getResource("testfile1.txt").toString());
        fileData.setCreationDate("28-01-2017");
        fileDatas.add(fileData);

        ProcessFiles processFiles = new ProcessFiles();
        actualResults = processFiles.processFiles(files, fileDatas);
        compareExpectedToActual(actualResults, expectedResults);
    }

    /**
     * Tests that more than one Thread is running after calling ProcessFile (as a minimum it will create one extra Thread
     * to process the passed in files).
     *
     * @throws InterruptedException as we put the Thread that runs this test to sleep, so that ProcessFile can finish running.
     */
    @Test
    public void testMultiThread() throws InterruptedException, URISyntaxException {//check that processfile makes two more threads run
        // at least 2 Threads should be running (main and the one created by process file).
        actualResults = null;
        File testFile1 = new File(getClass().getResource("testfile1.txt").toURI());
        File testFile2 = new File(getClass().getResource("testfile2.txt").toURI());
        File testFile3 = new File(getClass().getResource("testfile3.txt").toURI());

        ArrayList<File> files = new ArrayList<>();
        files.add(testFile1);
        files.add(testFile2);
        files.add(testFile3);

        ArrayList<FileData> fileDatas = new ArrayList<>();
        FileData fileData = new FileData("testfile1.txt", getClass().getResource("testfile1.txt").toString());
        fileData.setCreationDate("28-01-2017");
        fileDatas.add(fileData);
        fileData = new FileData("testfile2.txt", getClass().getResource("testfile2.txt").toString());
        fileData.setCreationDate("28-01-2017");
        fileDatas.add(fileData);
        fileData = new FileData("testfile3.txt", getClass().getResource("testfile3.txt").toString());
        fileData.setCreationDate("28-01-2017");
        fileDatas.add(fileData);

        ProcessFiles processFiles = new ProcessFiles();
        Task<Boolean> task = new Task<Boolean>() {
            @Override
            protected Boolean call() throws Exception {
                processFiles.processFiles(files, fileDatas);
                return null;
            }
        };
        new Thread(task).start();
        System.out.println("Actual Threads running" + Thread.activeCount());
        Assert.assertEquals(true, Thread.activeCount() >= 2);
        Thread.sleep(10000);//give the backend.process.Engine time to finish running
        Assert.assertEquals(true, true);//seems to be an error that if you finish with a thread sleep it will stop running the other operations on other threads (i.e. releasing semaphores and setting state to FINISHED).
    }

    /**
     * Tests the processing of a sample docx file, by checking the processed Result against expected Results. Only looks
     * at dates picked out (checking only it is processing the File correctly, i.e. picking out the correct sentences, and
     * their dates).
     *
     * @throws InterruptedException as we are putting the Thread that runs this test to sleep to let ProcessFile finish running.
     * @throws ParseException       for the Dates that we create for the expected Results.
     */
    @Test
    public void testSampleFileProcessDocx() throws InterruptedException, ParseException, URISyntaxException {
        actualResults = null;
        ArrayList<Result> expectedResults = new ArrayList<>();
        //base date is 29/12/2016 (creation date of file)
        Result result1 = new Result();
        TimelineDate timelineDate1 = new TimelineDate();
        timelineDate1.setDate1(simpleDateFormat.parse("2017-01-27"));
        result1.setTimelineDate(timelineDate1);
        expectedResults.add(result1);

        Result result2 = new Result();
        TimelineDate timelineDate2 = new TimelineDate();
        timelineDate2.setDate1(simpleDateFormat.parse("2008-01-01"));//second sentence talks about donations between 2008 and spring of this year (2016)
        timelineDate2.setDate2(simpleDateFormat.parse("2017-05-31"));//last day of spring
        result2.setTimelineDate(timelineDate2);
        expectedResults.add(result2);

        Result result3 = new Result();
        TimelineDate timelineDate3 = new TimelineDate();
        timelineDate3.setDate1(simpleDateFormat.parse("1980-01-01"));//donated since the early 1980s (range so 1980 ->1989)
        timelineDate3.setDate2(simpleDateFormat.parse("1989-12-31"));//last day of the 1980s
        result3.setTimelineDate(timelineDate3);
        expectedResults.add(result3);

        Result result4 = new Result();
        TimelineDate timelineDate4 = new TimelineDate();
        timelineDate4.setDate1(simpleDateFormat.parse("1996-01-01"));//in 1996 an event happened (doesnt specify day or month, just year)
        result4.setTimelineDate(timelineDate4);
        expectedResults.add(result4);

        File testFile = new File(getClass().getResource("testfile4.docx").toURI());
        ArrayList<File> files = new ArrayList<>();
        files.add(testFile);

        ArrayList<FileData> fileDatas = new ArrayList<>();
        FileData fileData = new FileData("testfile4.docx", getClass().getResource("testfile4.docx").toString());
        fileData.setCreationDate("28-01-2017");
        fileDatas.add(fileData);

        ProcessFiles processFiles = new ProcessFiles();
        actualResults = processFiles.processFiles(files, fileDatas);
        compareExpectedToActual(actualResults, expectedResults);
    }

    /**
     * Test for processing PDF Files, by checking the processed Results to the expected Results. Checks only their dates
     * as we are interested in it picking out the right sentences and their dates in the documents parsed in.
     *
     * @throws InterruptedException as we are putting the Thread that runs this test to sleep to allow time for ProcessFile to finish (runs on separate Thread).
     * @throws ParseException       for the Dates that we create for the expected Results.
     */
    @Test
    public void testSampleFilePDF() throws InterruptedException, ParseException, URISyntaxException {
        actualResults = null;
        ArrayList<Result> expectedResults = new ArrayList<>();
        //base date 2016-12-29

        //on Friday = 2016-12-30
        //1967 = 1967-01-01
        Result result1 = new Result();
        TimelineDate timelineDate1 = new TimelineDate();
        timelineDate1.setDate1(simpleDateFormat.parse("2017-01-27"));
        result1.setTimelineDate(timelineDate1);
        expectedResults.add(result1);

        Result result2 = new Result();
        TimelineDate timelineDate2 = new TimelineDate();
        timelineDate2.setDate1(simpleDateFormat.parse("1967-01-01"));
        result2.setTimelineDate(timelineDate2);
        expectedResults.add(result2);

        File testFile = new File(getClass().getResource("testfile5.pdf").toURI());//bbc article (http://www.bbc.com/news/world-us-canada-38451258)
        ArrayList<File> files = new ArrayList<>();
        files.add(testFile);

        ArrayList<FileData> fileDatas = new ArrayList<>();
        FileData fileData = new FileData("testfile5.pdf", getClass().getResource("testfile5.pdf").toString());
        fileData.setCreationDate("28-01-2017");
        fileDatas.add(fileData);

        ProcessFiles processFiles = new ProcessFiles();
        actualResults = processFiles.processFiles(files, fileDatas);
        System.out.println(actualResults);
        compareExpectedToActual(actualResults, expectedResults);
    }

    /**
     * Compares an actual list of Results with a expected one, by just looking at the dates picked out. Checking that the backend.process.Engine
     * is picking out the right sentence and producing the right dates, not caring about the sentence trimming or subject picking.
     *
     * @param actualResults   the list of produced backend.process.Result objects
     * @param expectedResults the list of expected backend.process.Result objects
     */
    private void compareExpectedToActual(List<Result> actualResults, List<Result> expectedResults) {
        System.out.println("Actual Results: " + actualResults);
        Assert.assertEquals(actualResults.size(), expectedResults.size());//if this does not hold then the test fails
        for (int i = 0; i < actualResults.size(); i++) {
            Assert.assertEquals(actualResults.get(i).getTimelineDate(), expectedResults.get(i).getTimelineDate());
        }
    }
}
\end{lstlisting}
\subsubsection{ProduceRangesTest.java}
\begin{lstlisting}
package backend;

import backend.process.Result;
import backend.process.TimelineDate;
import backend.ranges.ProduceRanges;
import backend.ranges.Range;
import org.junit.Assert;
import org.junit.Test;

import java.text.ParseException;
import java.text.SimpleDateFormat;
import java.util.ArrayList;
import java.util.List;

/**
 * Test class for the processing of Result Lists into Range Trees.
 */
public class ProduceRangesTest {

    /**
     * Tests a simple Range Tree formed by processing the fake List of Results, by matching it to an expected Range Tree.
     *
     * @throws ParseException when setting the Dates for the fake Results.
     */
    @Test
    public void testSingleTree() throws ParseException {
        List<Result> resultList = new ArrayList<>();

        SimpleDateFormat simpleDateFormat = new SimpleDateFormat("dd-MM-yyyy G");

        List<Range> expectedTrees = new ArrayList<>();
        Range range = new Range(simpleDateFormat.parse("01-01-0001 AD"), simpleDateFormat.parse("31-12-9999 AD"));
        Range range1980 = new Range(simpleDateFormat.parse("01-01-1980 AD"), simpleDateFormat.parse("31-12-1989 AD"));
        Range range2008 = new Range(simpleDateFormat.parse("01-01-2008 AD"), simpleDateFormat.parse("31-05-2016 AD"));
        Range range2015 = new Range(simpleDateFormat.parse("01-01-2015 AD"), null);
        range2008.addChild(range2015);
        range.addChild(range1980);
        range.addChild(range2008);
        expectedTrees.add(range);

        Result result1;
        TimelineDate timelineDate1;
        result1 = new Result();
        timelineDate1 = new TimelineDate();
        timelineDate1.setDate1(simpleDateFormat.parse("01-01-0001 AD"));
        timelineDate1.setDate2(simpleDateFormat.parse("31-12-9999 AD"));
        timelineDate1.getRange();
        result1.setTimelineDate(timelineDate1);
        resultList.add(result1);
        range.add(result1);
        result1 = new Result();
        timelineDate1 = new TimelineDate();
        timelineDate1.setDate1(simpleDateFormat.parse("01-01-1980 AD"));
        timelineDate1.setDate2(simpleDateFormat.parse("31-12-1989 AD"));
        timelineDate1.getRange();
        result1.setTimelineDate(timelineDate1);
        resultList.add(result1);
        range1980.add(result1);
        result1 = new Result();
        timelineDate1 = new TimelineDate();
        timelineDate1.setDate1(simpleDateFormat.parse("01-01-2008 AD"));
        timelineDate1.setDate2(simpleDateFormat.parse("31-05-2016 AD"));
        timelineDate1.getRange();
        result1.setTimelineDate(timelineDate1);
        resultList.add(result1);
        range2008.add(result1);
        result1 = new Result();
        timelineDate1 = new TimelineDate();
        timelineDate1.setDate1(simpleDateFormat.parse("01-01-2015 AD"));
        timelineDate1.getRange();
        result1.setTimelineDate(timelineDate1);
        resultList.add(result1);
        range2015.add(result1);

        ProduceRanges produceRanges = new ProduceRanges();
        produceRanges.produceRanges(resultList);
        List<Range> actualTrees = produceRanges.getTrees();

        checkTrees(expectedTrees, actualTrees);

    }

    /**
     * Tests that multiple Trees are formed when the List of fake Results contains at least 2 disjoint Results (based on
     * their Dates), and then checks that the Trees formed match the expected List of Trees.
     *
     * @throws ParseException when setting the Dates for the fake Results.
     */
    @Test
    public void testSeparateTrees() throws ParseException {
        List<Result> resultList = new ArrayList<>();

        SimpleDateFormat simpleDateFormat = new SimpleDateFormat("dd-MM-yyyy G");

        List<Range> expectedTrees = new ArrayList<>();
        Range range = new Range(simpleDateFormat.parse("01-01-0001 AD"), simpleDateFormat.parse("31-12-9999 AD"));
        Range range1980 = new Range(simpleDateFormat.parse("01-01-1980 AD"), simpleDateFormat.parse("31-12-1989 AD"));
        Range range2008 = new Range(simpleDateFormat.parse("01-01-2008 AD"), simpleDateFormat.parse("31-05-2016 AD"));
        Range range2015 = new Range(simpleDateFormat.parse("01-01-2015 AD"), null);
        range2008.addChild(range2015);
        range.addChild(range1980);
        range.addChild(range2008);
        Range range400BC = new Range(simpleDateFormat.parse("01-01-499 BC"), simpleDateFormat.parse("31-12-400 BC"));
        expectedTrees.add(range400BC);
        expectedTrees.add(range);

        Result result1;
        TimelineDate timelineDate1;
        result1 = new Result();
        timelineDate1 = new TimelineDate();
        timelineDate1.setDate1(simpleDateFormat.parse("01-01-0001 AD"));
        timelineDate1.setDate2(simpleDateFormat.parse("31-12-9999 AD"));
        timelineDate1.getRange();
        result1.setTimelineDate(timelineDate1);
        resultList.add(result1);
        range.add(result1);
        result1 = new Result();
        timelineDate1 = new TimelineDate();
        timelineDate1.setDate1(simpleDateFormat.parse("01-01-1980 AD"));
        timelineDate1.setDate2(simpleDateFormat.parse("31-12-1989 AD"));
        timelineDate1.getRange();
        result1.setTimelineDate(timelineDate1);
        resultList.add(result1);
        range1980.add(result1);
        result1 = new Result();
        timelineDate1 = new TimelineDate();
        timelineDate1.setDate1(simpleDateFormat.parse("01-01-2008 AD"));
        timelineDate1.setDate2(simpleDateFormat.parse("31-05-2016 AD"));
        timelineDate1.getRange();
        result1.setTimelineDate(timelineDate1);
        resultList.add(result1);
        range2008.add(result1);
        result1 = new Result();
        timelineDate1 = new TimelineDate();
        timelineDate1.setDate1(simpleDateFormat.parse("01-01-2015 AD"));
        timelineDate1.getRange();
        result1.setTimelineDate(timelineDate1);
        resultList.add(result1);
        range2015.add(result1);
        result1 = new Result();
        timelineDate1 = new TimelineDate();
        timelineDate1.setDate1(simpleDateFormat.parse("01-01-0499 BC"));
        timelineDate1.setDate2(simpleDateFormat.parse("31-12-0400 BC"));
        timelineDate1.getRange();
        result1.setTimelineDate(timelineDate1);
        resultList.add(result1);
        range400BC.add(result1);


        ProduceRanges produceRanges = new ProduceRanges();
        produceRanges.produceRanges(resultList);
        List<Range> actualTrees = produceRanges.getTrees();

        checkTrees(expectedTrees, actualTrees);
    }

    /**
     * Tests the two Trees formed based on a list of fake Results. One of the Trees has a depth of 2, such that it contains
     * a Range with a sub-Range with a sub-sub-Range (more complex than previous tests). Checks that the Results produce
     * the expected two Range Trees (one in BC and one in AD), with the AD Range Tree containing 10 results.
     *
     * @throws ParseException when producing the Dates for the fake Results.
     */
    @Test
    public void testComplexSeparateTrees() throws ParseException {
        SimpleDateFormat simpleDateFormat = new SimpleDateFormat("dd-MM-yyyy G");
        List<Range> expectedTrees = new ArrayList<>();

        Range range01To99 = new Range(simpleDateFormat.parse("01-01-0001 AD"), simpleDateFormat.parse("31-12-9999 AD"));
        Range range1980 = new Range(simpleDateFormat.parse("01-01-1980 AD"), simpleDateFormat.parse("31-12-1989 AD"));
        Range range2008 = new Range(simpleDateFormat.parse("01-01-2008 AD"), simpleDateFormat.parse("31-05-2016 AD"));
        Range range2015 = new Range(simpleDateFormat.parse("01-01-2015 AD"), null);
        Range range201612To20 = new Range(simpleDateFormat.parse("12-12-2016 AD"), simpleDateFormat.parse("20-12-2016 AD"));
        Range range201612To18 = new Range(simpleDateFormat.parse("12-12-2016 AD"), simpleDateFormat.parse("18-12-2016 AD"));
        Range range201615To20 = new Range(simpleDateFormat.parse("15-12-2016 AD"), simpleDateFormat.parse("20-12-2016 AD"));
        Range range201601 = new Range(simpleDateFormat.parse("01-11-2016 AD"), null);
        Range range1996 = new Range(simpleDateFormat.parse("01-01-1996 AD"), null);
        Range range201623 = new Range(simpleDateFormat.parse("23-12-2016 AD"), null);
        Range range0499 = new Range(simpleDateFormat.parse("01-01-0499 BC"), simpleDateFormat.parse("31-12-0400 BC"));


        range01To99.getChildren().add(range1980);
        range01To99.getChildren().add(range1996);
        range01To99.getChildren().add(range2008);
        range01To99.getChildren().add(range201601);
        range01To99.getChildren().add(range201612To20);
        range01To99.getChildren().add(range201623);
        range2008.getChildren().add(range2015);
        range201612To20.getChildren().add(range201612To18);
        range201612To20.getChildren().add(range201615To20);

        expectedTrees.add(range0499);
        expectedTrees.add(range01To99);

        List<Result> results = new ArrayList<>();
        Result result;
        TimelineDate timelineDate;
        result = new Result();
        timelineDate = new TimelineDate();
        timelineDate.setDate1(simpleDateFormat.parse("01-01-0001 AD"));
        timelineDate.setDate2(simpleDateFormat.parse("31-12-9999 AD"));
        timelineDate.getRange();
        result.setTimelineDate(timelineDate);
        range01To99.add(result);
        range01To99.add(result);
        results.add(result);
        results.add(result);
        result = new Result();
        timelineDate = new TimelineDate();
        timelineDate.setDate1(simpleDateFormat.parse("01-01-1980 AD"));
        timelineDate.setDate2(simpleDateFormat.parse("31-12-1989 AD"));
        timelineDate.getRange();
        result.setTimelineDate(timelineDate);
        range1980.add(result);
        results.add(result);
        result = new Result();
        timelineDate = new TimelineDate();
        timelineDate.setDate1(simpleDateFormat.parse("01-01-2008 AD"));
        timelineDate.setDate2(simpleDateFormat.parse("31-05-2016 AD"));
        timelineDate.getRange();
        result.setTimelineDate(timelineDate);
        range2008.add(result);
        results.add(result);
        result = new Result();
        timelineDate = new TimelineDate();
        timelineDate.setDate1(simpleDateFormat.parse("01-01-2015 AD"));
        timelineDate.getRange();
        result.setTimelineDate(timelineDate);
        range2015.add(result);
        results.add(result);
        result = new Result();
        timelineDate = new TimelineDate();
        timelineDate.setDate1(simpleDateFormat.parse("12-12-2016 AD"));
        timelineDate.setDate2(simpleDateFormat.parse("18-12-2016 AD"));
        timelineDate.getRange();
        result.setTimelineDate(timelineDate);
        range201612To18.add(result);
        results.add(result);
        result = new Result();
        timelineDate = new TimelineDate();
        timelineDate.setDate1(simpleDateFormat.parse("15-12-2016 AD"));
        timelineDate.setDate2(simpleDateFormat.parse("20-12-2016 AD"));
        timelineDate.getRange();
        result.setTimelineDate(timelineDate);
        range201615To20.add(result);
        results.add(result);
        result = new Result();
        timelineDate = new TimelineDate();
        timelineDate.setDate1(simpleDateFormat.parse("01-11-2016 AD"));
        timelineDate.getRange();
        result.setTimelineDate(timelineDate);
        range201601.add(result);
        results.add(result);
        result = new Result();
        timelineDate = new TimelineDate();
        timelineDate.setDate1(simpleDateFormat.parse("01-01-1996 AD"));
        timelineDate.getRange();
        result.setTimelineDate(timelineDate);
        range1996.add(result);
        results.add(result);
        result = new Result();
        timelineDate = new TimelineDate();
        timelineDate.setDate1(simpleDateFormat.parse("23-12-2016 AD"));
        timelineDate.getRange();
        result.setTimelineDate(timelineDate);
        range201623.add(result);
        results.add(result);
        result = new Result();
        timelineDate = new TimelineDate();
        timelineDate.setDate1(simpleDateFormat.parse("01-01-0499 BC"));
        timelineDate.setDate2(simpleDateFormat.parse("31-12-0400 BC"));
        timelineDate.getRange();
        result.setTimelineDate(timelineDate);
        range0499.add(result);
        results.add(result);

        ProduceRanges produceRanges = new ProduceRanges();
        produceRanges.produceRanges(results);

        List<Range> actual = produceRanges.getTrees();
        checkTrees(expectedTrees, actual);
    }

    /**
     * Asserts that both lists of Ranges, the expected and actual, are equal, ie they are the same size, and that each
     * contain the same Range objects at the same indices.
     *
     * @param expected the expected list of Ranges.
     * @param actual   the actual list of Ranges.
     */
    private void checkTrees(List<Range> expected, List<Range> actual) {
        Assert.assertEquals(expected.size(), actual.size());
        for (int i = 0; i < expected.size(); i++) {
            Assert.assertEquals(expected.get(i), actual.get(i));
        }
    }
}
\end{lstlisting}
\subsubsection{SystemStateTest.java}
\begin{lstlisting}
package backend;

import backend.process.FileData;
import backend.process.ProcessFiles;
import backend.system.BackEndSystem;
import backend.system.SystemState;
import org.junit.Assert;
import org.junit.Test;

import java.io.File;
import java.net.URISyntaxException;
import java.util.ArrayList;

/**
 * Test class to test the change of System States as files are being Processed. Checks for STARTED, PROCESSING, and PROCESSED
 * System States, as these are the three main different states.
 */
public class SystemStateTest {

    /**
     * Checks the initial state, then when processing, and then when it should be finished.
     *
     * @throws InterruptedException as we make the Test thread wait for ProcessFile to finish processing the Files.
     */
    @Test
    public void testSystemStatePROCESSING() throws InterruptedException, URISyntaxException {
        BackEndSystem.getInstance().setSystemState(SystemState.STARTED);
        Assert.assertEquals(SystemState.STARTED, BackEndSystem.getInstance().getSystemState());

        ProcessFiles processFiles = new ProcessFiles();
        File testFile = new File(getClass().getResource("testfile1.txt").toURI());
        ArrayList<File> files = new ArrayList<>();
        files.add(testFile);

        ArrayList<FileData> fileDatas = new ArrayList<>();
        for(File file: files){
            fileDatas.add(new FileData(file));
        }

        Runnable newRunnable = new Runnable() {
            @Override
            public void run() {
                processFiles.processFiles(files, fileDatas);
            }
        };

        Thread thread = new Thread(newRunnable);
        thread.start();
        Thread.sleep(100);
        System.out.println("System state: " + BackEndSystem.getInstance().getSystemState());
        Assert.assertEquals(SystemState.PROCESSING, BackEndSystem.getInstance().getSystemState());

        Thread.sleep(5000);

    }


    /**
     * Checks that the state is set to FINISHED when files have been processed, and the result has been returned.
     */
    @Test
    public void testSystemStateFINISHED() throws URISyntaxException {
        BackEndSystem.getInstance().setSystemState(SystemState.STARTED);
        Assert.assertEquals(SystemState.STARTED, BackEndSystem.getInstance().getSystemState());

        ProcessFiles processFiles = new ProcessFiles();
        File testFile = new File(getClass().getResource("testfile1.txt").toURI());
        ArrayList<File> files = new ArrayList<>();
        files.add(testFile);

        ArrayList<FileData> fileDatas = new ArrayList<>();
        for(File file: files){
            fileDatas.add(new FileData(file));
        }

        processFiles.processFiles(files, fileDatas);

        Assert.assertEquals(SystemState.FINISHED, BackEndSystem.getInstance().getSystemState());

    }
}
\end{lstlisting}
\subsubsection{TimelineDateTest.java}
\begin{lstlisting}
package backend;

import backend.process.TimelineDate;
import edu.stanford.nlp.time.SUTime;
import org.junit.Assert;
import org.junit.Test;

import java.text.ParseException;
import java.text.SimpleDateFormat;
import java.util.Date;

/**
 * Test for the parsing of Normalized Entity Tag in backend.process.TimelineDate, and how they are processed.
 */
public class TimelineDateTest {
    private static final String PRESENT_REF = "PRESENT_REF";
    private static final String PAST_REF = "PAST_REF";
    private static final String FUTURE_REF = "FUTURE_REF";
    private SimpleDateFormat simpleDateFormat = new SimpleDateFormat("yyyy-MM-dd");
    private String baseDate = "2016-12-30";//doesn't affect the output of most tests, but has to be still passed in.

    /**
     * Passes in a simple date of the format yyyy-MM-dd into backend.process.TimelineDate, and checks the output date.
     *
     * @throws ParseException when comparing the expected Date to check.
     */
    @Test
    public void testTimelineDateProcessSimple() throws ParseException {
        String dateStr = "2016-12-24";
        Date expectedDate = simpleDateFormat.parse(dateStr);

        TimelineDate timelineDate = new TimelineDate();
        timelineDate.parse(dateStr, baseDate);//should make a Date object of the same type, not much processing.

        Assert.assertEquals(expectedDate, timelineDate.getDate1());
        Assert.assertEquals(null, timelineDate.getDate2());//should have only picked up one date
    }

    /**
     * Passes in a normalized entity tag for a year range, 198X, and checks the output to the expected start and end Date
     * of the year range.
     *
     * @throws ParseException when creating the expected Dates to check.
     */
    @Test
    public void testTimelineDateProcessYearRange() throws ParseException {
        String yearRange = "198X";
        String expectedStartDate = "1980-01-01";
        String expectedEndDate = "1989-12-31";

        Date expectedStart = simpleDateFormat.parse(expectedStartDate);
        Date expectedEnd = simpleDateFormat.parse(expectedEndDate);

        TimelineDate timelineDate = new TimelineDate();
        timelineDate.parse(yearRange, baseDate);

        Assert.assertEquals(expectedStart, timelineDate.getDate1());
        Assert.assertEquals(expectedEnd, timelineDate.getDate2());

    }

    /**
     * Passes in a normalized entity tag for an actual year range, 1980-01-01/2016-10-25, and checks the output to the
     * expected start and end Date.
     *
     * @throws ParseException when creating the expected Dates.
     */
    @Test
    public void testTimelineDateProcessActualDateRange() throws ParseException {
        String input = "1980-01-01/2016-10-25";
        String expectedStartDate = "1980-01-01";
        String expectedEndDate = "2016-10-25";

        Date expectedStart = simpleDateFormat.parse(expectedStartDate);
        Date expectedEnd = simpleDateFormat.parse(expectedEndDate);

        TimelineDate timelineDate = new TimelineDate();
        timelineDate.parse(input, baseDate);

        System.out.println(timelineDate);
        Assert.assertEquals(expectedStart, timelineDate.getDate1());
        Assert.assertEquals(expectedEnd, timelineDate.getDate2());
    }

    /**
     * Passes in a year and a season to backend.process.TimelineDate, and compares the output to an expected start and end Dates.
     *
     * @throws ParseException when creating the expected Dates.
     */
    @Test
    public void testTimelineDateProcessSeason() throws ParseException {
        String input = "1980-SP";
        String expectedStartDate = "1980-03-01";
        String expectedEndDate = "1980-05-31";

        Date expectedStart = simpleDateFormat.parse(expectedStartDate);
        Date expectedEnd = simpleDateFormat.parse(expectedEndDate);

        TimelineDate timelineDate = new TimelineDate();
        timelineDate.parse(input, baseDate);

        System.out.println(timelineDate);

        Assert.assertEquals(expectedStart, timelineDate.getDate1());
        Assert.assertEquals(expectedEnd, timelineDate.getDate2());
    }

    /**
     * Passes in a year with a week number to backend.process.TimelineDate, and compares the output to an expected start and end Dates.
     *
     * @throws ParseException when creating the expected Dates.
     */
    @Test
    public void testTimelineDateProcessWeekNumber() throws ParseException {
        String input = "2016-W47";
        String expectedStartDate = "2016-11-21";
        String expectedEndDate = "2016-11-27";

        Date expectedStart = simpleDateFormat.parse(expectedStartDate);
        Date expectedEnd = simpleDateFormat.parse(expectedEndDate);

        TimelineDate timelineDate = new TimelineDate();
        timelineDate.parse(input, baseDate);

        System.out.println(timelineDate);

        Assert.assertEquals(expectedStart, timelineDate.getDate1());
        Assert.assertEquals(expectedEnd, timelineDate.getDate2());
    }

    /**
     * Passes in a year and a month, and checks the output to an expected start Date.
     *
     * @throws ParseException when creating the expected Date.
     */
    @Test
    public void testTimelineDateProcessMonth() throws ParseException {
        String input = "2016-10";
        String expectedStartDate = "2016-10-01";

        Date expectedStart = simpleDateFormat.parse(expectedStartDate);

        TimelineDate timelineDate = new TimelineDate();
        timelineDate.parse(input, baseDate);

        System.out.println(timelineDate);

        Assert.assertEquals(expectedStart, timelineDate.getDate1());
    }

    /**
     * Tests the resulting date produced when a normalized entity tag for now (PRESENT_REF) is passed into TimelineDate.
     *
     * @throws ParseException when creating the expected Date.
     */
    @Test
    public void testTimelineDateProcessNow() throws ParseException {
        //the base date passed in should be the date used when it processes REFERENCE_NOW
        //base date being
        Date expectedDate = simpleDateFormat.parse(baseDate);

        TimelineDate timelineDate = new TimelineDate();
        timelineDate.parse(PRESENT_REF, baseDate);

        Assert.assertEquals(expectedDate, timelineDate.getDate1());
        Assert.assertEquals(null, timelineDate.getDate2());//shouldnt have a second date
    }

    /**
     * Tests the processing of the INTERSECT data (duration) that can come with a normalized entity tag for Dates.
     * Tests for a simple INTERSECT (4 Years).
     *
     * @throws ParseException when creating the expected Date.
     */
    @Test
    public void testTimelineDateProcessINTERSECT() throws ParseException {
        String expectedResultDuration = "Period: 4 Year(s)";
        Date expectedResultDate = simpleDateFormat.parse(baseDate);
        String input = baseDate + " INTERSECT P4Y";

        TimelineDate timelineDate = new TimelineDate();
        timelineDate.parse(input, baseDate);

        Assert.assertEquals(expectedResultDate, timelineDate.getDate1());
        Assert.assertEquals(null, timelineDate.getDate2());
        Assert.assertEquals(expectedResultDuration, timelineDate.getDurationData());
    }

    /**
     * Tests the processing of a more complex INTERSECT data (duration), including time.
     *
     * @throws ParseException when creating the expected Date.
     */
    @Test
    public void testTimelineDateProcessINTERSECTInclTime() throws ParseException {
        String expectedResultDuration = "Period: 3 Year(s) 6 Month(s) 4 Day(s) Time: 12 Hour(s) 30 Minute(s) 5 Second(s)";
        Date expectedResultDate = simpleDateFormat.parse(baseDate);
        String input = baseDate + " INTERSECT P3Y6M4DT12H30M5S";

        TimelineDate timelineDate = new TimelineDate();
        timelineDate.parse(input, baseDate);

        Assert.assertEquals(expectedResultDate, timelineDate.getDate1());
        Assert.assertEquals(null, timelineDate.getDate2());
        Assert.assertEquals(expectedResultDuration, timelineDate.getDurationData());
    }

    /**
     * Tests the auto correct for Dates, where the day value has been over-estimated (that day value for that month
     * does not exist).
     *
     * @throws ParseException when creating the expected Date.
     */
    @Test
    public void testTimelineDateProcessAutoCorrectDates() throws ParseException {
        String inputWrongDate = "2016-12-32";
        String expectedDateStr = "2016-12-31";
        Date expectedDate = simpleDateFormat.parse(expectedDateStr);

        TimelineDate timelineDate = new TimelineDate();
        timelineDate.parse(inputWrongDate, baseDate);

        Assert.assertEquals(expectedDate, timelineDate.getDate1());
        Assert.assertEquals(null, timelineDate.getDate2());
    }

    /**
     * Tests the auto correct for Dates, where the day value has been over-estimated, such that the month value has to
     * also be updated.
     *
     * @throws ParseException when creating the expected Date.
     */
    @Test
    public void testTimelineDateProcessAutoCorrectDatesAndMonths() throws ParseException {
        String inputWrongDate = "2016-12-00";
        String expectedDateStr = "2016-11-30";
        Date expectedDate = simpleDateFormat.parse(expectedDateStr);

        TimelineDate timelineDate = new TimelineDate();
        timelineDate.parse(inputWrongDate, baseDate);

        Assert.assertEquals(expectedDate, timelineDate.getDate1());
        Assert.assertEquals(null, timelineDate.getDate2());
    }

    /**
     * Tests the processing of a full BC normalized entity Date, of the format -yyyy-MM-dd.
     *
     * @throws ParseException when creating the expected Date.
     */
    @Test
    public void testTimelineDateBCFullDate() throws ParseException {
        SimpleDateFormat simpleDateFormatBC = new SimpleDateFormat("yyyy-MM-dd G");
        String input = "-0004-12-31";
        String expectedDateStr = "0004-12-31 BC";
        Date expectedDate = simpleDateFormatBC.parse(expectedDateStr);

        TimelineDate timelineDate = new TimelineDate();
        timelineDate.parse(input, baseDate);

        Assert.assertEquals(expectedDate, timelineDate.getDate1());
        Assert.assertNull(timelineDate.getDate2());
    }

    /**
     * Tests the processing of only a BC year normalized entity Date, of the format -yyyy, to see how TimelineDate gives
     * it a Date value (i.e. assumes it means the start of the year, so month 01 and day 01)
     *
     * @throws ParseException when creating the expected Date.
     */
    @Test
    public void testTimelineDateBCYearDate() throws ParseException {
        SimpleDateFormat simpleDateFormatBC = new SimpleDateFormat("yyyy-MM-dd G");
        String input = "-0004";
        String expectedDateStr = "0004-01-01 BC";
        Date expectedDate = simpleDateFormatBC.parse(expectedDateStr);

        TimelineDate timelineDate = new TimelineDate();
        timelineDate.parse(input, baseDate);

        Assert.assertEquals(expectedDate, timelineDate.getDate1());
        Assert.assertNull(timelineDate.getDate2());
    }

    /**
     * Tests the processing of a BC Date, that has a day value that is wrong for that month. Looks at how this Date is
     * reduced until it correct (assuming we over-estimated the day value).
     *
     * @throws ParseException when creating the expected Date.
     */
    @Test
    public void testTimelineDateBCWrongDate() throws ParseException {
        SimpleDateFormat simpleDateFormatBC = new SimpleDateFormat("yyyy-MM-dd G");
        String input = "-0004-12-32";
        String expectedDateStr = "0004-12-31 BC";
        Date expectedDate = simpleDateFormatBC.parse(expectedDateStr);

        TimelineDate timelineDate = new TimelineDate();
        timelineDate.parse(input, baseDate);

        Assert.assertEquals(expectedDate, timelineDate.getDate1());
        Assert.assertNull(timelineDate.getDate2());
    }

    /**
     * Tests the building of a range of Dates, for a BC date. Testing how the system infers the start and end Dates for
     * BC dates. (Should be the same as AD dates, just that the year values to towards 0 as it is B.C.)
     *
     * @throws ParseException when creating the expected Date.
     */
    @Test
    public void testTimelineDateBCDateRange() throws ParseException {
        SimpleDateFormat simpleDateFormatBC = new SimpleDateFormat("yyyy-MM-dd G");
        String input = "-04XX-12-32";
        String expectedDateStr1 = "0499-12-31 BC";
        String expectedDateStr2 = "0400-12-31 BC";
        Date expectedDate1 = simpleDateFormatBC.parse(expectedDateStr1);
        Date expectedDate2 = simpleDateFormatBC.parse(expectedDateStr2);

        TimelineDate timelineDate = new TimelineDate();
        timelineDate.parse(input, baseDate);

        Assert.assertEquals(expectedDate1, timelineDate.getDate1());
        Assert.assertEquals(expectedDate2, timelineDate.getDate2());
    }

    /**
     * Tests the case when the input is a week number with the day of the week.
     *
     * @throws ParseException when the expected date is created.
     */
    @Test
    public void testWeekDate() throws ParseException {
        String input = "2016-W47-7";//the seventh day in week 47 in iso 8601 is the sunday
        String expectedDateStr = "2016-11-27";//this the actual date for sunday
        Date expectedDate = simpleDateFormat.parse(expectedDateStr);
        TimelineDate timelineDate = new TimelineDate();
        timelineDate.parse(input, baseDate);
        Assert.assertEquals(expectedDate, timelineDate.getDate1());
        Assert.assertEquals(null, timelineDate.getDate2());
    }

    /**
     * Tests that when a Past reference is passed, a range of dates from the start of AD years until the base date is
     * created.
     *
     * @throws ParseException when the expected dates are created.
     */
    @Test
    public void testPastRef() throws ParseException {
        String input = PAST_REF;
        String expectedDateStr1 = "0001-01-01";
        String expectedDateStr2 = baseDate;//past ref points from start of year in AD to base date
        Date expectedDate1 = simpleDateFormat.parse(expectedDateStr1);
        Date expectedDate2 = simpleDateFormat.parse(expectedDateStr2);
        TimelineDate timelineDate = new TimelineDate();
        timelineDate.parse(input, baseDate);
        Assert.assertEquals(expectedDate1, timelineDate.getDate1());
        Assert.assertEquals(expectedDate2, timelineDate.getDate2());
    }

    /**
     * Tests that when a Future reference is passed, a range of dates from the base date until the end of times is
     * created.
     *
     * @throws ParseException when creating the expected dates.
     */
    @Test
    public void testFutureRef() throws ParseException {
        String input = FUTURE_REF;
        String expectedDateStr1 = baseDate;
        String expectedDateStr2 = "9999-12-31";
        Date expectedDate1 = simpleDateFormat.parse(expectedDateStr1);
        Date expectedDate2 = simpleDateFormat.parse(expectedDateStr2);
        TimelineDate timelineDate = new TimelineDate();
        timelineDate.parse(input, baseDate);
        Assert.assertEquals(expectedDate1, timelineDate.getDate1());
        Assert.assertEquals(expectedDate2, timelineDate.getDate2());
    }

    /**
     * Tests that TimelineDate enforces the Rules of expanding the Range of dates.
     * I.e. when pass a date which is earlier than the current date1, then date1 should be updated, or if a date is
     * passed in that is further away in the future than date2, then date2 should be updated. However, if a date is
     * passed in that is within the range, then the date1, and date2 values should not be changed.
     *
     * @throws ParseException when creating the expected Dates.
     */
    @Test
    public void testEnforceRule() throws ParseException {
        String input = baseDate;
        Date expectedDate1 = simpleDateFormat.parse(input);
        Date expectedDate2 = null;
        TimelineDate timelineDate = new TimelineDate();
        timelineDate.parse(input, baseDate);
        Assert.assertEquals(expectedDate1, timelineDate.getDate1());
        Assert.assertNull(timelineDate.getDate2());
        input = "2015-12-15";
        timelineDate.parse(input, baseDate);
        expectedDate1 = simpleDateFormat.parse(input);
        Assert.assertEquals(expectedDate1, timelineDate.getDate1());
        Assert.assertNull(timelineDate.getDate2());
        input = "2077-01-05";
        timelineDate.parse(input, baseDate);
        expectedDate2 = simpleDateFormat.parse(input);
        Assert.assertEquals(expectedDate1, timelineDate.getDate1());
        Assert.assertEquals(expectedDate2, timelineDate.getDate2());
        input = "2055-12-24";
        timelineDate.parse(input, baseDate);//dates shouldnt have changed as range hasnt expanded
        Assert.assertEquals(expectedDate1, timelineDate.getDate1());
        Assert.assertEquals(expectedDate2, timelineDate.getDate2());
    }

}
\end{lstlisting}
\subsubsection{ToJSONTest.java}
\begin{lstlisting}
package backend;

import backend.helpers.ToJSON;
import backend.process.FileData;
import backend.process.Result;
import backend.process.TimelineDate;
import org.junit.Assert;
import org.junit.Test;

import java.text.ParseException;
import java.text.SimpleDateFormat;
import java.util.ArrayList;

/**
 * Test Class to test the toJSON class (ie taking a list of Results and producing JSON Strings).
 */
public class ToJSONTest {
    private static final SimpleDateFormat simpleDateFormat = new SimpleDateFormat("yyyy-MM-dd G");

    /**
     * Tests the JSON String produced for a List of one Result object which has all the data initialised (i.e. all the
     * fields in the JSON should be populated).
     *
     * @throws ParseException when producing the dates (date1, date2 from TimelineDate) for the test Result object.
     */
    @Test
    public void testCompleteJSON() throws ParseException {
        Result result = new Result();
        TimelineDate timelineDate = new TimelineDate();
        timelineDate.setDate1(simpleDateFormat.parse(returnDate("2016", "02", "14", false)));
        timelineDate.setDate2(simpleDateFormat.parse(returnDate("2016", "02", "15", false)));
        result.setTimelineDate(timelineDate);
        result.addSubject("Valentines Day");
        result.addSubject("Party");
        result.setEvent("On Valentines Day we had a huge party!");
        FileData fileData = new FileData("party.txt", "FAKEPATH");
        fileData.setCreationDate("02-02-2017");
        result.setFileData(fileData);

        ArrayList<Result> results = new ArrayList<>();
        results.add(result);
        //process the full result

        String actualResultJson = ToJSON.toJSON(results);
        String expectedJson = "[{\"date1\":\"14-02-2016 AD\",\"date2\":\"15-02-2016 AD\",\"subjects\":[\"Valentines Day\",\"Party\"],\"event\":\"On Valentines Day we had a huge party!\",\"from\":{\"filename\":\"party.txt\",\"baseDate\":\"02-02-2017\"}}]";
        Assert.assertEquals(expectedJson, actualResultJson);
    }

    /**
     * Tests the JSON String produced for a List of one Result object, where only some of its data fields have been
     * populated (date1, event, and from fields should have data).
     *
     * @throws ParseException when producing the date1 in TimelineDate for the test Result object.
     */
    @Test
    public void testPartialJSON() throws ParseException {
        Result result = new Result();
        TimelineDate timelineDate = new TimelineDate();
        timelineDate.setDate1(simpleDateFormat.parse(returnDate("2017", "02", "14", false)));
        result.setTimelineDate(timelineDate);
        result.setEvent("On Valentines Day we had a huge party!");
        FileData fileData = new FileData("party.txt", "FAKEPATH");
        fileData.setCreationDate("02-02-2017");
        result.setFileData(fileData);

        ArrayList<Result> results = new ArrayList<>();
        results.add(result);

        String actualResultJson = ToJSON.toJSON(results);
        String expectedJson = "[{\"date1\":\"14-02-2017 AD\",\"subjects\":[],\"event\":\"On Valentines Day we had a huge party!\",\"from\":{\"filename\":\"party.txt\",\"baseDate\":\"02-02-2017\"}}]";
        Assert.assertEquals(expectedJson, actualResultJson);
    }

    /**
     * Tests the JSON String produced for a List of two Result objects, where one is fully populated (i.e. all fields
     * have non null values) and the other is partially populated (i.e. only date1, event and from fields are populated).
     *
     * @throws ParseException when producing the date1 and/or date2 in TimelineDate for the test Result objects (since
     *                        its a list).
     */
    @Test
    public void testMultipleResultJSON() throws ParseException {
        ArrayList<Result> results = new ArrayList<>();

        Result result = new Result();
        TimelineDate timelineDate = new TimelineDate();
        timelineDate.setDate1(simpleDateFormat.parse(returnDate("2016", "02", "14", false)));
        timelineDate.setDate2(simpleDateFormat.parse(returnDate("2016", "02", "15", false)));
        result.setTimelineDate(timelineDate);
        result.addSubject("Valentines Day");
        result.addSubject("Party");
        result.setEvent("On Valentines Day we had a huge party!");
        FileData fileData = new FileData("party.txt", "FAKEPATH");
        fileData.setCreationDate("02-02-2017");
        result.setFileData(fileData);
        results.add(result);

        result = new Result();
        timelineDate = new TimelineDate();
        timelineDate.setDate1(simpleDateFormat.parse(returnDate("2017", "02", "14", false)));
        result.setTimelineDate(timelineDate);
        result.setEvent("On Valentines Day we had a huge party!");
        fileData = new FileData("party.txt", "FAKEPATH");
        fileData.setCreationDate("02-02-2017");
        result.setFileData(fileData);
        results.add(result);

        String actualResultsJson = ToJSON.toJSON(results);
        String expectedJson = "[{\"date1\":\"14-02-2016 AD\",\"date2\":\"15-02-2016 AD\",\"subjects\":[\"Valentines Day\",\"Party\"],\"event\":\"On Valentines Day we had a huge party!\",\"from\":{\"filename\":\"party.txt\",\"baseDate\":\"02-02-2017\"}},
        	{\"date1\":\"14-02-2017 AD\",\"subjects\":[],\"event\":\"On Valentines Day we had a huge party!\",\"from\":{\"filename\":\"party.txt\",\"baseDate\":\"02-02-2017\"}}]";
        Assert.assertEquals(expectedJson, actualResultsJson);
    }

    /**
     * Tests the JSON String produced for a List of one Result object, that has just been initialised, and had no data
     * set to it.
     */
    @Test
    public void testBareResult() {
        Result result = new Result();
        ArrayList<Result> results = new ArrayList<>();
        results.add(result);

        String actualResultJson = ToJSON.toJSON(results);
        String expectedJson = "[{\"subjects\":[],\"event\":\"\",\"from\":{}}]";
        Assert.assertEquals(expectedJson, actualResultJson);
    }

    /**
     * Tests the JSON String produced for an empty List.
     */
    @Test
    public void testNoResults() {
        ArrayList<Result> results = new ArrayList<>();
        String actualResultJson = ToJSON.toJSON(results);
        String expectedJson = "[]";
        Assert.assertEquals(expectedJson, actualResultJson);
    }

    /**
     * Produces a date String of the format yyyy-MM-dd, for the given input.
     *
     * @param year  the year of the date
     * @param month the month of the date
     * @param day   the day of the date
     * @return a String of the format yyyy-MM-dd
     */
    private String returnDate(String year, String month, String day, boolean isBC) {
        if (isBC) {
            return String.format("\%s-\%s-\%s BC", year, month, day);
        }
        return String.format("\%s-\%s-\%s AD", year, month, day);
    }

}
\end{lstlisting}
\section{src/test/resources/backend}
\subsection{testfile1.txt, textfile2.txt, and textfile3.txt}
\begin{lstlisting}
On the 12th of December I ran tests on my final year project. The tests did not go so well. Yesterday I played games. It was fun playing games! Tomorrow I am going to study. Last week I went to watch a football match.
\end{lstlisting}
\subsection{testfile4.docx}
\begin{lstlisting}
On Friday, the Washington Post came out with the latest from its long-running investigation into Trump's charitable donations. 
In its latest story, the paper called 420-plus charities with some connection to Trump but found only one personal gift from him between 2008 and the spring of this year. 
But the Post did find nearly $8m that Trump has donated from his own pocket since the early 1980s.
One of the bizarre episodes the paper recounts is that in 1996, Trump showed up without an invitation to a charity for the Association to Benefit Children where he took a seat on the stage that had been reserved for a major donor, despite not being a donor himself. 
In response to the piece, the Trump campaign told the Post that he "has personally donated tens of millions of dollars... to charitable causes".
\end{lstlisting}
\subsection{testfile5.pdf}
\begin{lstlisting}
On Friday, the US chose not to veto a UN Security Council resolution calling for an end to Israeli settlement construction, leading to an angry response from Israel. The issue of Jewish settlements is one of the most contentious between Israel and the Palestinians, who see them as an obstacle to peace and the creation of a viable Palestinian state. More than 500,000 Jews live in about 140 settlements built since Israel's 1967 occupation of the West Bank and East Jerusalem. The settlements are considered illegal under international law, though Israel disputes this.
\end{lstlisting}

\section{build.gradle}
\begin{lstlisting}
apply plugin: 'idea'
apply plugin: 'java'
apply plugin: 'application'

mainClassName = 'frontend.Main'

repositories {
	mavenCentral()
}

dependencies {
    //added slf4j  because one of the libraries points to it, hence throwing the error Failed to load class “org.slf4j.impl.StaticLoggerBinder”
    //when running the build
   	 compile group: 'org.slf4j', name: 'slf4j-api', version: '1.7.7'
    	compile group: 'org.slf4j', name: 'slf4j-simple', version: '1.7.7'
	compile group: 'edu.stanford.nlp', name: 'stanford-corenlp', version: '3.7.0'
	compile group: 'edu.stanford.nlp', name: 'stanford-corenlp', version: '3.7.0', classifier: 'models-english'
	compile group: 'org.apache.pdfbox', name: 'pdfbox', version: '2.0.4'
	compile group: 'org.apache.poi', name: 'poi', version: '3.15'
	compile group: 'org.apache.xmlbeans', name: 'xmlbeans', version: '2.6.0'
	compile group: 'org.fxmisc.richtext', name: 'richtextfx', version: '0.6.10'
	compile group: 'org.apache.commons', name: 'commons-lang3', version: '3.5'
	compile group: 'org.apache.poi', name: 'poi-ooxml', version: '3.15'
	compile group: 'com.google.code.gson', name: 'gson', version: '2.8.0'
	compile group: 'joda-time', name: 'joda-time', version: '2.9.7'

	testCompile 'junit:junit:4.12'
    	testCompile group: 'org.slf4j', name: 'slf4j-api', version: '1.7.7'
    	testCompile group: 'org.slf4j', name: 'slf4j-simple', version: '1.7.7'
	testCompile group: 'edu.stanford.nlp', name: 'stanford-corenlp', version: '3.7.0'
	testCompile group: 'edu.stanford.nlp', name: 'stanford-corenlp', version: '3.7.0', classifier: 'models-english'
	testCompile group: 'org.apache.pdfbox', name: 'pdfbox', version: '2.0.4'
	testCompile group: 'org.apache.poi', name: 'poi', version: '3.15'
	testCompile group: 'org.apache.xmlbeans', name: 'xmlbeans', version: '2.6.0'
	testCompile group: 'org.fxmisc.richtext', name: 'richtextfx', version: '0.6.10'
	testCompile group: 'org.apache.commons', name: 'commons-lang3', version: '3.5'
	testCompile group: 'org.apache.poi', name: 'poi-ooxml', version: '3.15'
	testCompile group: 'com.google.code.gson', name: 'gson', version: '2.8.0'
	testCompile group: 'joda-time', name: 'joda-time', version: '2.9.7'

}

jar {
	baseName = 'timeline-extraction'
	version = '0.1.0'
}

configurations.all {
    resolutionStrategy {
        force 'xml-apis:xml-apis:1.4.01'
    }
}

task copyTestResources(type: Copy) {
    from "${projectDir}/src/test/resources"
    into "${buildDir}/classes/test"
}
processTestResources.dependsOn copyTestResources

task showMeCache {
    doLast {
        configurations.compile.each { println it }
    }
}
\end{lstlisting}